\documentclass[11pt,a4paper,oneside,openany]{book}

\usepackage[T1]{fontenc}
\usepackage[utf8]{inputenc}
\usepackage{libertinus}
\usepackage{microtype}
\usepackage[a4paper,margin=27mm,headheight=15pt]{geometry}
\usepackage{amsmath,amssymb,amsthm,mathtools}
\usepackage{bm}
\usepackage{mathrsfs}
\usepackage{booktabs,longtable,array,tabularx}
\usepackage{enumitem}
\usepackage{xcolor}
\usepackage{tikz}
\usetikzlibrary{arrows.meta,positioning,calc}
\usepackage[most]{tcolorbox}
\usepackage{fancyhdr}
\usepackage{xurl}
\usepackage{hyperref}
\usepackage[nameinlink,noabbrev,capitalise]{cleveref}
\usepackage{makeidx}
\makeindex

\definecolor{deepolive}{RGB}{67,79,45}
\definecolor{softolive}{RGB}{236,239,226}
\definecolor{darkteal}{RGB}{31,83,84}
\definecolor{warmgray}{RGB}{92,87,78}
\definecolor{softcoral}{RGB}{246,225,219}
\definecolor{linkblue}{RGB}{31,76,122}

\hypersetup{
  colorlinks=true,
  linkcolor=deepolive,
  citecolor=darkteal,
  urlcolor=linkblue,
  pdftitle={A Unified Calculus of Combinatorial Formulae},
  pdfauthor={A proof-complete synthesis of five drafts},
  pdfsubject={Stirling, Bell, Eulerian, Bernoulli, Faa di Bruno, Lagrange-Good inversion, and asymptotic formulae}
}

\setlist{itemsep=2pt,topsep=4pt,parsep=1pt}
\allowdisplaybreaks[3]
\raggedbottom
\setcounter{tocdepth}{2}
\setcounter{secnumdepth}{3}

\pagestyle{fancy}
\fancyhf{}
\fancyhead[L]{\small\nouppercase{\leftmark}}
\fancyhead[R]{\small\thepage}
\renewcommand{\headrulewidth}{0.4pt}
\fancypagestyle{plain}{%
  \fancyhf{}%
  \fancyfoot[C]{\thepage}%
  \renewcommand{\headrulewidth}{0pt}%
}

\newtheorem{theorem}{Theorem}[chapter]
\newtheorem{lemma}[theorem]{Lemma}
\newtheorem{proposition}[theorem]{Proposition}
\newtheorem{corollary}[theorem]{Corollary}
\newtheorem{identity}[theorem]{Identity}
\theoremstyle{definition}
\newtheorem{definition}[theorem]{Definition}
\newtheorem{example}[theorem]{Example}
\newtheorem{algorithm}[theorem]{Algorithm}
\theoremstyle{remark}
\newtheorem{remark}[theorem]{Remark}
\newtheorem{historical}[theorem]{Historical note}

\newtcolorbox{masterbox}[1][]{
  enhanced,breakable,colback=softolive,colframe=deepolive,
  boxrule=0.7pt,arc=1mm,left=2mm,right=2mm,top=1.5mm,bottom=1.5mm,
  title={Master principle},fonttitle=\bfseries,#1
}
\newtcolorbox{auditbox}[1][]{
  enhanced,breakable,colback=softcoral,colframe=warmgray,
  boxrule=0.6pt,arc=1mm,left=2mm,right=2mm,top=1.5mm,bottom=1.5mm,
  title={Editorial normalization},fonttitle=\bfseries,#1
}

\newcommand{\N}{\mathbb N}
\newcommand{\Z}{\mathbb Z}
\newcommand{\Q}{\mathbb Q}
\newcommand{\R}{\mathbb R}
\newcommand{\C}{\mathbb C}
\newcommand{\F}{\mathbb F}
\newcommand{\Bell}[1]{\mathsf B_{#1}}
\newcommand{\BellP}[2]{\mathcal B_{#1,#2}}
\newcommand{\BellC}[1]{\mathcal B_{#1}}
\newcommand{\BerP}[2]{\mathrm B_{#1}(#2)}
\newcommand{\EulP}[2]{\mathrm E_{#1}(#2)}
\newcommand{\GenP}[2]{\mathrm G_{#1}(#2)}
\newcommand{\OBellP}[2]{\widehat{\mathcal B}_{#1,#2}}
\newcommand{\stirone}[2]{\genfrac{[}{]}{0pt}{}{#1}{#2}}
\newcommand{\stirtwo}[2]{\genfrac{\{} {\}}{0pt}{}{#1}{#2}}
\newcommand{\eulerian}[2]{\genfrac{\langle}{\rangle}{0pt}{}{#1}{#2}}
\newcommand{\eulertwo}[2]{\left\langle\!\left\langle\genfrac{}{}{0pt}{}{#1}{#2}\right\rangle\!\right\rangle}
\newcommand{\falling}[2]{#1^{\underline{#2}}}
\newcommand{\rising}[2]{#1^{\overline{#2}}}
\newcommand{\coeff}[2]{[#1^{#2}]}
\newcommand{\res}{\operatorname*{Res}}
\newcommand{\Li}{\operatorname{Li}}
\newcommand{\Var}{\operatorname{Var}}
\newcommand{\E}{\mathbb E}
\newcommand{\1}{\mathbf 1}
\newcommand{\wt}{\operatorname{wt}}
\newcommand{\cyc}{\operatorname{cyc}}
\newcommand{\setpart}{\operatorname{Part}}
\newcommand{\dd}{\,\mathrm d}
\newcommand{\bitand}{\mathbin{\&}}
\DeclareMathOperator{\tr}{tr}
\DeclareMathOperator{\lcm}{lcm}
\DeclareMathOperator{\sgn}{sgn}
\DeclareMathOperator{\id}{id}
\DeclareMathOperator{\diag}{diag}

\title{\Huge\bfseries A Unified Calculus of Combinatorial Formulae\\[4mm]
\Large Stirling, Bell, Eulerian, Bernoulli, and Catalan Structures;\\
Fa\`a di Bruno, Lagrange--Good Inversion, and Asymptotic Methods}
\author{A proof-complete synthesis of five drafts and two reference dossiers}
\date{September 2026}

\begin{document}
\frontmatter
\hypersetup{pageanchor=false}
\maketitle
\hypersetup{pageanchor=true}

\chapter*{Abstract}
\addcontentsline{toc}{chapter}{Abstract}
This monograph develops a unified calculus for the coefficients that occur when one
changes polynomial bases, partitions labelled sets, differentiates or composes
functions, and reverses formal or analytic power series.  Its principal objects are
the signed and unsigned Stirling numbers of the first kind, the Stirling numbers of
the second kind, Lah numbers, Bell numbers, exponential and ordinary Bell
polynomials, Eulerian numbers of types $A$ and $B$, second-order Eulerian numbers,
Fa\`a di Bruno's formula, and the Lagrange inversion and reversion theorems.

The treatment is deliberately proof-oriented.  Definitions are connected by
bijections and change-of-basis arguments; generating functions are derived rather
than merely recorded; recurrences, convolutions, congruences, determinant formulae,
probabilistic interpretations, special-function identities, and asymptotic formulae
are proved from a small collection of reusable principles.  Applications include
finite differences, differential-operator normal ordering, moments and cumulants,
symmetric polynomials, cycle indices, Hermite polynomials, polynomial sequences of
binomial type, negative polylogarithms, Bernoulli-number identities, Poisson moments,
Catalan and Fuss--Catalan series, the Lambert $W$ function, and Laplace-type
asymptotic integrals.

The source archive uses the letter $B$ for Bell numbers, Bell polynomials, Bernoulli
numbers, type-$B$ Eulerian numbers, and binary-tree numbers. Here these are separated
notationally. Five overlapping drafts have been merged at section level rather than
concatenated. New chapters integrate complementary material from the supplied
Wikipedia and MathWorld folders and from primary online references: coordinate-free
and multivariate Fa\`a di Bruno formulas, Good's inversion theorem, partition-lattice
M\"obius inversion, Sheffer and Riordan theory, Bernoulli--Euler--Genocchi and Cauchy
polynomials, complementary Bell numbers, Euler--Maclaurin summation, Darboux formulas,
Pochhammer parameter derivatives, polygamma and Barnes $G$ functions, Kepler's
equation, Narayana numbers, and associahedral face enumeration. A final concordance
records the role of each draft, the deduplication rules, and independent consistency
checks.

\chapter*{How to read this work}
\addcontentsline{toc}{chapter}{How to read this work}
The chapters are arranged by mathematical dependence, not by the order of the source
files.  The two master ideas are:
\begin{enumerate}
  \item a coefficient is often easiest to compute after changing to the basis in
        which the relevant operator is diagonal; and
  \item an exponential generating function remembers the decomposition of a
        labelled object into components.
\end{enumerate}
Stirling numbers implement the first idea.  Bell polynomials implement the second.
Eulerian numbers mediate between powers and shifted binomial coefficients.  Fa\`a di
Bruno's formula is the differential form of labelled composition, while Lagrange
inversion is the coefficient form of implicit substitution.

A result stated for formal power series requires no convergence.  When an analytic
version is asserted, its domain hypotheses are written explicitly.  Statements in the five drafts that report historical computations---for example, lists of Bell
primes or verified exact modular periods---are distinguished from theorems and are
accompanied by a certificate procedure rather than misrepresented as symbolic
identities.

\begin{auditbox}
Three systematic repairs are used throughout.
\begin{enumerate}
\item Bell numbers are denoted $\Bell n$, complete exponential Bell polynomials by
      $\BellC n$, partial Bell polynomials by $\BellP nk$, and Bernoulli numbers by
      $\beta_n$.
\item Type-$B$ Eulerian numbers are indexed by descents in positions
      $0,1,\ldots,n-1$ with $\pi_0=0$.  This removes the one-step shift in one formula
      in the archive.
\item The archived Moser--Wyman display contains unnamed quantities
      $B,P_i,Q_i$.  We give the exact saddle-point scheme that defines every
      correction coefficient recursively and state the leading expansion in closed
      form; no undefined placeholder is retained.
\end{enumerate}
\end{auditbox}

\tableofcontents

\chapter*{Notation}
\addcontentsline{toc}{chapter}{Notation}
For $n\in\N$, write
\[
 \falling{x}{n}=x(x-1)\cdots(x-n+1),\qquad
 \rising{x}{n}=x(x+1)\cdots(x+n-1),
\]
with both expressions equal to $1$ when $n=0$.  Coefficient extraction is denoted by
$\coeff z n F(z)$, and $\res_z F(z)\dd z$ means the coefficient of $z^{-1}$ in a
formal Laurent series.  Our principal arrays are
\[
 s(n,k),\qquad \stirone nk=(-1)^{n-k}s(n,k),\qquad
 \stirtwo nk,\qquad \eulerian nk,
\]
where $s(n,k)$ is the signed first-kind Stirling number, $\stirone nk$ its unsigned
version, $\stirtwo nk$ the second-kind Stirling number, and $\eulerian nk$ the
ordinary Eulerian number.  Bell numbers are $\Bell n$; exponential partial and
complete Bell polynomials are $\BellP nk$ and $\BellC n$.

Unless stated otherwise, sums over an impossible index range are zero, and
$0^0=1$ in coefficient formulae.  Bernoulli numbers use
\[
 \frac{t}{e^t-1}=\sum_{n\ge0}\beta_n\frac{t^n}{n!},
 \qquad \beta_1=-\frac12.
\]

\mainmatter

\part{Coefficient Calculus and Labelled Structures}

\chapter{Formal power series and coefficient extraction}
\label{ch:formal}

\section{Formal series as an algebraic environment}
Let $R$ be a commutative $\Q$-algebra.  The ring $R[[z]]$ consists of expressions
$F(z)=\sum_{n\ge0}f_nz^n$ with termwise addition and Cauchy multiplication.  A series
is invertible precisely when $f_0$ is invertible in $R$.  If $G(0)=0$, then the
composition $F(G(z))$ is well-defined because each coefficient receives only
finitely many contributions.

\begin{lemma}[Coefficient rules]\label{lem:coeff-rules}
For $F(z)=\sum f_nz^n$ and $G(z)=\sum g_nz^n$,
\begin{align}
 \coeff z n F(z)G(z)&=\sum_{j=0}^n f_jg_{n-j},\label{eq:cauchy}\\
 \coeff z n F'(z)&=(n+1)f_{n+1},\label{eq:der-coeff}\\
 \coeff z n \frac{F(z)}{1-az}&=\sum_{j=0}^n a^{n-j}f_j.\label{eq:geom-conv}
\end{align}
If $F$ is analytic on $|z|\le r$, then also
\begin{equation}\label{eq:cauchy-coeff}
 f_n=\frac{1}{2\pi i}\int_{|z|=r}\frac{F(z)}{z^{n+1}}\dd z,
 \qquad |f_n|\le r^{-n}\max_{|z|=r}|F(z)|.
\end{equation}
\end{lemma}
\begin{proof}
The first identity is the definition of Cauchy multiplication.  The second follows
by differentiating term by term, and the third by multiplying with
$\sum_{m\ge0}a^mz^m$.  The analytic formula is Cauchy's coefficient formula, and the
bound follows by estimating the contour integral by its length $2\pi r$.
\end{proof}

\section{Formal residues}
For a Laurent series $A(z)$, set
$\res_z A(z)\dd z=\coeff z{-1}A(z)$.  Then
\begin{equation}\label{eq:res-coeff}
 \coeff z n F(z)=\res_z\frac{F(z)}{z^{n+1}}\dd z,
 \qquad \res_z A'(z)\dd z=0.
\end{equation}
The second identity is simply the fact that the derivative of $z^m$ has no
$z^{-1}$ term.

\begin{theorem}[Formal change of variables for residues]\label{thm:res-subst}
Let $u=u(z)\in zR[[z]]$ have invertible linear coefficient.  For every Laurent series
$A$ for which the composition is defined,
\begin{equation}\label{eq:formal-res-subst}
 \res_z A(u(z))u'(z)\dd z=\res_u A(u)\dd u.
\end{equation}
\end{theorem}
\begin{proof}
By linearity it suffices to take $A(u)=u^m$.  If $m\ne-1$, then
$u(z)^mu'(z)=\frac{1}{m+1}(u(z)^{m+1})'$ has residue zero.  For $m=-1$,
$u'(z)/u(z)=z^{-1}+H(z)$ with $H\in R[[z]]$, because
$u(z)=az(1+O(z))$.  Thus both sides have residue $1$.
\end{proof}

\begin{masterbox}
The combination of \cref{eq:res-coeff,eq:formal-res-subst} turns implicit
substitution into coefficient extraction.  This is the entire algebraic engine of
Lagrange inversion; no analytic contour is needed until one asks about convergence.
\end{masterbox}

\section{Ordinary and exponential generating functions}
For a sequence $(a_n)$, its ordinary and exponential generating functions are
\[
 A(z)=\sum_{n\ge0}a_nz^n,
 \qquad \widehat A(z)=\sum_{n\ge0}a_n\frac{z^n}{n!}.
\]
Ordinary multiplication gives unweighted convolution.  Exponential multiplication
gives binomial convolution:
\begin{equation}\label{eq:egf-product}
 \widehat A(z)\widehat B(z)
 =\sum_{n\ge0}\left(\sum_{j=0}^n\binom nj a_jb_{n-j}\right)\frac{z^n}{n!}.
\end{equation}
This binomial factor is the number of ways to choose which labels belong to the
first component.

\section{Finite differences and factorial bases}
Let $E$ be the shift operator $(Ef)(x)=f(x+1)$ and
$\Delta=E-1$.  Then
\begin{equation}\label{eq:finite-diff-falling}
 \Delta\falling{x}{k}=k\falling{x}{k-1},
 \qquad
 \sum_{i=0}^{n-1}\falling{i}{k}=\frac{\falling{n}{k+1}}{k+1}.
\end{equation}
\begin{proof}
The first follows from
$\falling{x+1}{k}-\falling{x}{k}=k\falling{x}{k-1}$ after factoring the common
product.  Summing this telescoping identity gives the second formula.
\end{proof}

\begin{theorem}[Newton expansion]\label{thm:newton-expansion}
Every polynomial $p$ of degree at most $d$ has the unique expansion
\begin{equation}\label{eq:newton-expansion}
 p(x)=\sum_{k=0}^d\frac{\Delta^kp(0)}{k!}\falling{x}{k}
     =\sum_{k=0}^d\Delta^kp(0)\binom{x}{k}.
\end{equation}
\end{theorem}
\begin{proof}
The factorial polynomials form a basis because their degrees are distinct and their
leading coefficients are one.  If $p=\sum a_k\falling{x}{k}$, applying $\Delta^j$
and setting $x=0$ kills terms with $k>j$ by a remaining factor $x$, kills terms
with $k<j$ by excessive differencing, and gives $j!a_j$ for $k=j$.
\end{proof}

\section{Labelled products and the exponential formula}
A labelled combinatorial class is one whose atoms carry distinct labels.  Suppose a
connected component of size $j$ has total weight $x_j$.  A set of $k$ such components
with total size $n$ has weight enumerator $\BellP nk(x_1,x_2,\ldots)$; a set of any
number of components has enumerator $\BellC n(x_1,\ldots,x_n)$.

\begin{theorem}[The exponential formula]\label{thm:exponential-formula}
Let
\[
 C(z)=\sum_{j\ge1}x_j\frac{z^j}{j!}.
\]
Then
\begin{align}
 \frac{C(z)^k}{k!}
  &=\sum_{n\ge k}\BellP nk(x_1,\ldots,x_{n-k+1})\frac{z^n}{n!},
  \label{eq:partial-bell-egf}\\
 e^{uC(z)}
  &=\sum_{n\ge k\ge0}\BellP nk(x_1,\ldots)u^k\frac{z^n}{n!},
  \label{eq:bivariate-bell-egf}\\
 e^{C(z)}
  &=\sum_{n\ge0}\BellC n(x_1,\ldots,x_n)\frac{z^n}{n!}.
  \label{eq:complete-bell-egf}
\end{align}
\end{theorem}
\begin{proof}
Expand $C(z)^k/k!$ multinomially.  If $j_i$ components have size $i$, then
$\sum_i j_i=k$ and $\sum_iij_i=n$, and the coefficient of $z^n$ is
\[
 \prod_i\frac1{j_i!}\left(\frac{x_i}{i!}\right)^{j_i}.
\]
Multiplication by $n!$ labels the atoms.  This proves the first formula.  Summing it
with weight $u^k$ proves the second, and setting $u=1$ proves the third.
\end{proof}

\begin{corollary}[Partition-type count]\label{cor:partition-type}
The number of partitions of an $n$-element labelled set with exactly $j_i$ blocks of
size $i$ is
\begin{equation}\label{eq:partition-type-count}
 \frac{n!}{\prod_{i\ge1}(i!)^{j_i}j_i!},
 \qquad \sum_iij_i=n.
\end{equation}
\end{corollary}
\begin{proof}
Order the $n$ labels, cut the ordering into blocks of prescribed sizes, then divide
by $i!$ for the internal order of each size-$i$ block and by $j_i!$ for the order of
blocks of equal size.
\end{proof}

\part{Stirling Numbers, Lah Numbers, and Basis Changes}

\chapter{The two Stirling triangles}
\label{ch:stirling-def}

\section{First-kind numbers: cycles and factorial coefficients}
\begin{definition}
The signed Stirling numbers of the first kind $s(n,k)$ and the unsigned numbers
$\stirone nk$ are defined by
\begin{align}
 \falling{x}{n}&=\sum_{k=0}^n s(n,k)x^k,\label{eq:signed-first-def}\\
 \rising{x}{n}&=\sum_{k=0}^n\stirone nkx^k,\label{eq:unsigned-first-def}\\
 s(n,k)&=(-1)^{n-k}\stirone nk.\label{eq:first-sign}
\end{align}
\end{definition}
The sign relation follows from
$\rising{x}{n}=(-1)^n\falling{-x}{n}$.

\begin{theorem}[Permutation interpretation]\label{thm:first-cycle}
$\stirone nk$ is the number of permutations of $[n]$ having exactly $k$ cycles.
Consequently,
\begin{equation}\label{eq:first-recurrence}
 \stirone{n+1}{k}=n\stirone nk+\stirone n{k-1},
\end{equation}
with $\stirone00=1$ and zero boundary values outside $0\le k\le n$.
\end{theorem}
\begin{proof}
Given a permutation of $[n+1]$ with $k$ cycles, either $n+1$ is a singleton cycle,
leaving a permutation of $[n]$ with $k-1$ cycles, or it is inserted immediately
after one of the $n$ existing elements in the cyclic notation of a permutation of
$[n]$ with $k$ cycles.  These alternatives are disjoint and reversible.
The same recurrence follows algebraically by multiplying
$\rising{x}{n}$ by $x+n$ and comparing coefficients.
\end{proof}

For $n=3$, the cycle counts are $\stirone33=1$, $\stirone32=3$, and
$\stirone31=2$, corresponding to the identity, the three transpositions, and the two
$3$-cycles.  Summing over the possible number of cycles gives
\begin{equation}\label{eq:first-row-sum}
 \sum_{k=0}^n\stirone nk=n!,
\end{equation}
which also follows from \cref{eq:unsigned-first-def} at $x=1$.

\section{Second-kind numbers: set partitions and surjections}
\begin{definition}
The Stirling number of the second kind $\stirtwo nk$ is the number of partitions of
$[n]$ into exactly $k$ nonempty blocks.
\end{definition}

\begin{theorem}[Second-kind recurrence]\label{thm:second-recurrence}
For $n,k\ge0$,
\begin{equation}\label{eq:second-recurrence}
 \stirtwo{n+1}{k}=k\stirtwo nk+\stirtwo n{k-1},
\end{equation}
with $\stirtwo00=1$ and zero boundary values outside $0\le k\le n$.
\end{theorem}
\begin{proof}
In a partition of $[n+1]$ into $k$ blocks, the new label either joins one of the $k$
blocks of a partition of $[n]$, or forms a singleton block beside a partition of
$[n]$ into $k-1$ blocks.
\end{proof}

\begin{theorem}[Surjection formula]\label{thm:second-explicit}
For all $n,k\ge0$,
\begin{equation}\label{eq:second-explicit}
 \stirtwo nk=\frac1{k!}\sum_{j=0}^k(-1)^{k-j}\binom kjj^n
 =\sum_{j=0}^k\frac{(-1)^{k-j}j^n}{(k-j)!j!}.
\end{equation}
\end{theorem}
\begin{proof}
There are $k!\stirtwo nk$ surjections from $[n]$ onto a fixed $k$-element codomain:
a set partition into fibres followed by an ordering of the fibres.  Inclusion--
exclusion over the set of omitted codomain values gives
$\sum_{j=0}^k(-1)^{k-j}\binom kjj^n$.  Division by $k!$ proves the formula.
\end{proof}

The simple diagonals follow at once:
\begin{equation}\label{eq:second-simple}
 \stirtwo nn=1,\qquad \stirtwo n1=1\ (n\ge1),\qquad
 \stirtwo n{n-1}=\binom n2,\qquad
 \stirtwo n2=2^{n-1}-1.
\end{equation}
The last identity chooses a nonempty proper subset as one block and divides by two
for exchanging the blocks.


\section{Alternative and reverse recurrences}
The ordinary recurrence builds each row from the preceding one.  Several less
familiar identities in the source archive instead recover an entry from entries to
its right or from earlier entries in the same column.  They were labelled there as
conjectures; in fact they are exact identities.

\begin{theorem}[Alternative second-kind recurrences]
\label{thm:second-reverse-recurrences}
For $1\le k<n$,
\begin{align}
 \stirtwo nk
 &=\frac{k^n}{k!}-\sum_{r=1}^{k-1}\frac{\stirtwo nr}{(k-r)!},
 \label{eq:second-triangular-explicit}\\
 (n-k)\stirtwo nk
 &=\sum_{j=2}^{n-k+1}(j-2)!\binom{-k}{j}
   \stirtwo n{k+j-1},
 \label{eq:second-reverse-row}\\
 (n-k)\stirtwo nk
 &=\sum_{j=2}^{n-k+1}(-1)^j\binom nj
   \stirtwo{n-j+1}{k}.
 \label{eq:second-reverse-column}
\end{align}
The terminal value in either reverse recursion is $\stirtwo nn=1$.
\end{theorem}
\begin{proof}
Substitute $x=k$ in the basis identity
$k^n=\sum_r\stirtwo nr\falling{k}{r}$.  Since
$\falling{k}{r}=k!/(k-r)!$ for $r\le k$, division by $k!$ and isolation of the
$r=k$ term gives \eqref{eq:second-triangular-explicit}.

For the row reversal, use the bivariate exponential generating function
\[
 F(x,y)=\exp\!\bigl(y(e^x-1)\bigr)
 =\sum_{n,k\ge0}\stirtwo nk\frac{x^n}{n!}y^k,
 \qquad u=e^x-1.
\]
The generating function of the left side of
\eqref{eq:second-reverse-row} is
\[
 xF_x-yF_y=yF\bigl(xe^x-u\bigr).
\]
On the other hand, because
$\binom{-k}{j}=(-1)^j\binom{k+j-1}{j}$, the generating function of its right side is
\begin{align*}
 y\sum_{j\ge2}\frac{(-1)^j}{j(j-1)}\,\partial_y^jF
 &=yF\sum_{j\ge2}\frac{(-1)^ju^j}{j(j-1)}\\
 &=yF\bigl((1+u)\log(1+u)-u\bigr)\\
 &=yF\bigl(xe^x-u\bigr).
\end{align*}
Here the middle series is obtained by differentiating it twice, summing the
geometric series, and restoring the two zero initial coefficients.  Equality of
coefficients proves \eqref{eq:second-reverse-row}.

For the column reversal fix $k$ and put
\[
 F_k(x)=\sum_{n\ge k}\stirtwo nk\frac{x^n}{n!}
       =\frac{(e^x-1)^k}{k!}.
\]
The two sides of \eqref{eq:second-reverse-column}, after multiplication by
$x^n/n!$ and summation over $n$, become respectively
\[
 xF_k'(x)-kF_k(x)
 \quad\text{and}\quad
 (e^{-x}-1+x)F_k'(x).
\]
They are equal because direct differentiation gives
$(1-e^{-x})F_k'(x)=kF_k(x)$.
\end{proof}

The factor interchange mentioned in the source produces a parallel pair for the
unsigned first-kind triangle.

\begin{theorem}[Reverse recurrences of the first kind]
\label{thm:first-reverse-recurrences}
For $1\le k<n$,
\begin{align}
 (n-k)\stirone nk
 &=\sum_{j=2}^{n-k+1}(-1)^j\binom{-k}{j}
   \stirone n{k+j-1},
 \label{eq:first-reverse-row}\\
 (n-k)\stirone nk
 &=\sum_{j=2}^{n-k+1}(j-2)!\binom nj
   \stirone{n-j+1}{k}.
 \label{eq:first-reverse-column}
\end{align}
\end{theorem}
\begin{proof}
Let $C_n(y)=\rising{y}{n}=\sum_r\stirone nr y^r$.  Since
\[
 C_n(1+y)=\frac{\rising{y}{n+1}}{y},
\]
the coefficient of $y^{k-1}$ on the two sides is, respectively,
\[
 \sum_{r\ge k-1}\binom r{k-1}\stirone nr
 \quad\text{and}\quad
 \stirone{n+1}{k}=n\stirone nk+\stirone n{k-1}.
\]
Remove the terms $r=k-1$ and $r=k$ and put $r=k+j-1$.  Because
$(-1)^j\binom{-k}{j}=\binom{k+j-1}{j}$, the remainder is precisely
\eqref{eq:first-reverse-row}.

For fixed $k$, set
\[
 C_k(x)=\sum_{n\ge k}\stirone nk\frac{x^n}{n!}
       =\frac{[-\log(1-x)]^k}{k!}.
\]
The exponential generating function of the right side of
\eqref{eq:first-reverse-column} is
\[
 C_k'(x)\sum_{j\ge2}\frac{x^j}{j(j-1)}
 =\bigl(x+(1-x)\log(1-x)\bigr)C_k'(x).
\]
The left side has generating function $xC_k'-kC_k$.  Their equality reduces to
\[
 -(1-x)\log(1-x)\,C_k'(x)=kC_k(x),
\]
which follows immediately from the displayed formula for $C_k$.
\end{proof}
\section{Tables and common boundary conventions}
The first rows are
\[
\begin{array}{c|rrrrrrr}
 n\backslash k&0&1&2&3&4&5&6\\\hline
0&1\\
1&0&1\\
2&0&1&1\\
3&0&2&3&1\\
4&0&6&11&6&1\\
5&0&24&50&35&10&1\\
6&0&120&274&225&85&15&1
\end{array}
\qquad
\begin{array}{c|rrrrrrr}
 n\backslash k&0&1&2&3&4&5&6\\\hline
0&1\\
1&0&1\\
2&0&1&1\\
3&0&1&3&1\\
4&0&1&7&6&1\\
5&0&1&15&25&10&1\\
6&0&1&31&90&65&15&1
\end{array}
\]
for $\stirone nk$ and $\stirtwo nk$, respectively.

\chapter{Factorial bases and inverse matrices}
\label{ch:bases}

\section{The basic changes of basis}
\begin{theorem}[Stirling basis transformations]\label{thm:stirling-basis}
For every $n\ge0$,
\begin{align}
 \falling{x}{n}&=\sum_{k=0}^ns(n,k)x^k,\label{eq:fall-to-power}\\
 \rising{x}{n}&=\sum_{k=0}^n\stirone nkx^k,\label{eq:rise-to-power}\\
 x^n&=\sum_{k=0}^n\stirtwo nk\falling{x}{k},\label{eq:power-to-fall}\\
 x^n&=\sum_{k=0}^n(-1)^{n-k}\stirtwo nk\rising{x}{k}.\label{eq:power-to-rise}
\end{align}
Equivalently,
\begin{equation}\label{eq:binomial-basis}
 x^n=\sum_{k=0}^n\stirtwo nk k!\binom{x}{k},\qquad
 \binom{x}{n}=\frac1{n!}\sum_{k=0}^ns(n,k)x^k.
\end{equation}
\end{theorem}
\begin{proof}
The first two identities are definitions.  In the Newton expansion
\cref{eq:newton-expansion} for $p(x)=x^n$, one has
\[
 \Delta^kx^n\big|_{x=0}=\sum_{j=0}^k(-1)^{k-j}\binom kjj^n
 =k!\stirtwo nk
\]
by \cref{eq:second-explicit}.  This gives \cref{eq:power-to-fall}.  Replacing $x$
by $-x$ and multiplying by $(-1)^n$ gives \cref{eq:power-to-rise}.  Division of
$\falling{x}{n}$ by $n!$ gives the second binomial-basis formula.
\end{proof}


\begin{corollary}[Shifted factorial evaluations]
\label{cor:shifted-stirling-evaluations}
For every $n\ge0$,
\begin{align}
 \sum_{k=1}^{n+1}\stirtwo{n+1}{k}\falling{x-1}{k-1}&=x^n,
 \label{eq:shifted-power-to-fall}\\
 \sum_{k=0}^{n}\stirtwo nk\falling{n}{k}&=n^n.
 \label{eq:stirling-n-to-n}
\end{align}
\end{corollary}
\begin{proof}
Apply \eqref{eq:power-to-fall} with exponent $n+1$ and use
$\falling{x}{k}=x\falling{x-1}{k-1}$.  Both sides are divisible by $x$ as
polynomials, and cancellation gives \eqref{eq:shifted-power-to-fall}.  Formula
\eqref{eq:stirling-n-to-n} is \eqref{eq:power-to-fall} evaluated at $x=n$.
\end{proof}
\begin{corollary}[Inverse matrices]\label{cor:stirling-inverse}
The lower triangular matrices $s=(s(n,k))$ and $S=(\stirtwo nk)$ are mutual inverses:
\begin{align}
 \sum_{j=k}^ns(n,j)\stirtwo jk&=\delta_{nk},\label{eq:stirling-inv1}\\
 \sum_{j=k}^n\stirtwo nj s(j,k)&=\delta_{nk}.\label{eq:stirling-inv2}
\end{align}
Equivalently,
\begin{align}
 \sum_{j=k}^n(-1)^{n-j}\stirone nj\stirtwo jk&=\delta_{nk},\\
 \sum_{j=k}^n(-1)^{j-k}\stirtwo nj\stirone jk&=\delta_{nk}.
\end{align}
\end{corollary}
\begin{proof}
Expand $\falling{x}{n}$ into powers and each power back into falling factorials.
Uniqueness of coordinates in the factorial basis gives the first matrix product.
The reverse substitution gives the second.
\end{proof}

\section{Power sums by basis adaptation}
The factorial basis makes discrete antidifferentiation immediate.  Combining
\cref{eq:power-to-fall,eq:finite-diff-falling} yields
\begin{equation}\label{eq:power-sum-stirling}
 \sum_{i=0}^{N}i^m
 =\sum_{k=0}^m\stirtwo mk\frac{\falling{N+1}{k+1}}{k+1}.
\end{equation}
For $m=4$ this is
\[
 \frac12\falling{N+1}{2}+\frac73\falling{N+1}{3}
 +\frac64\falling{N+1}{4}+\frac15\falling{N+1}{5},
\]
which expands to the familiar quartic power-sum polynomial.


\section{Factorial moments: Poisson variables and permutation fixed points}
The change from ordinary powers to falling factorials is especially effective for
integer-valued random variables.

\begin{theorem}[Poisson moments]
\label{thm:poisson-stirling-moments}
If $X$ has the Poisson distribution of mean $\lambda$, then
\begin{equation}
\label{eq:poisson-stirling-moments}
 \E X^n=\sum_{k=0}^{n}\stirtwo nk\lambda^k.
\end{equation}
In particular, $\E X^n=\Bell n$ when $\lambda=1$.
\end{theorem}
\begin{proof}
The $k$th falling-factorial moment is
\[
 \E\falling{X}{k}
 =e^{-\lambda}\sum_{r\ge k}\frac{\falling r k\lambda^r}{r!}
 =\lambda^ke^{-\lambda}\sum_{q\ge0}\frac{\lambda^q}{q!}
 =\lambda^k.
\]
Take expectations in the basis identity
$X^n=\sum_k\stirtwo nk\falling Xk$.
\end{proof}

\begin{theorem}[Moments of the number of fixed points]
\label{thm:fixed-point-stirling-moments}
Let $Y_m$ be the number of fixed points of a uniformly random permutation of
$[m]$.  Then
\begin{equation}
\label{eq:fixed-point-stirling-moments}
 \E Y_m^n=\sum_{k=0}^{m}\stirtwo nk
 =\sum_{k=0}^{\min(m,n)}\stirtwo nk.
\end{equation}
Thus, when $m\ge n$, the moment equals the Bell number $\Bell n$.
\end{theorem}
\begin{proof}
The random variable $\falling{Y_m}{k}$ counts ordered $k$-tuples of distinct fixed
points.  There are $\falling m k$ candidate tuples, and a specified tuple is fixed
with probability $(m-k)!/m!$.  Hence
\[
 \E\falling{Y_m}{k}=1\quad(0\le k\le m),
 \qquad
 \E\falling{Y_m}{k}=0\quad(k>m).
\]
Taking expectations in \eqref{eq:power-to-fall} proves the formula.  The upper
limit $m$ is essential: it records the finite size of the permutation, even when
$n>m$.
\end{proof}
\section{Lah numbers}
\begin{definition}
For $1\le k\le n$, the unsigned Lah number is
\begin{equation}\label{eq:lah-def}
 L(n,k)=\binom{n-1}{k-1}\frac{n!}{k!},
\end{equation}
with $L(0,0)=1$ and zero otherwise outside the triangle.
\end{definition}
Combinatorially, $L(n,k)$ counts partitions of $[n]$ into $k$ nonempty linearly
ordered lists, with the collection of lists unordered.

\begin{theorem}[Conversion between rising and falling factorials]\label{thm:lah-conv}
\begin{align}
 \rising{x}{n}&=\sum_{k=0}^nL(n,k)\falling{x}{k},\label{eq:lah-rise-fall}\\
 \falling{x}{n}&=\sum_{k=0}^n(-1)^{n-k}L(n,k)\rising{x}{k},\label{eq:lah-fall-rise}\\
 \sum_{j=k}^n(-1)^{j-k}L(n,j)L(j,k)&=\delta_{nk}.\label{eq:lah-inv}
\end{align}
Moreover,
\begin{equation}\label{eq:lah-stirling-product}
 L(n,k)=\sum_{j=k}^n\stirone nj\stirtwo jk.
\end{equation}
\end{theorem}
\begin{proof}
A direct generating-function calculation gives
\[
 \sum_{n\ge k}L(n,k)\frac{z^n}{n!}
 =\frac1{k!}\left(\frac{z}{1-z}\right)^k.
\]
Therefore
\[
 \sum_{n\ge0}\left(\sum_kL(n,k)\falling{x}{k}\right)\frac{z^n}{n!}
 =\sum_{k\ge0}\binom{x}{k}\left(\frac{z}{1-z}\right)^k
 =(1-z)^{-x}
 =\sum_{n\ge0}\rising{x}{n}\frac{z^n}{n!},
\]
proving \cref{eq:lah-rise-fall}.  Replace $x$ by $-x$ to obtain the inverse formula,
and compose the two transformations to obtain \cref{eq:lah-inv}.  Finally, expand a
rising factorial into powers by first-kind numbers and powers into falling factorials
by second-kind numbers; uniqueness gives \cref{eq:lah-stirling-product}.
\end{proof}

The inequalities
\begin{equation}\label{eq:stirling-lah-ineq}
 \stirtwo nk\le\stirone nk\le L(n,k)
\end{equation}
follow from explicit injections: a set partition may be canonically turned into a
permutation with one cycle per block by writing each block with its smallest member
first and arranging cycles by minima; forgetting the cyclic identifications embeds
permutations with $k$ cycles into unordered collections of $k$ lists.


\chapter{Generating functions, transforms, and operators}
\label{ch:stirling-gf}

\section{Exponential generating functions}
\begin{theorem}[Column generating functions]\label{thm:stirling-egfs}
For every fixed $k\ge0$,
\begin{align}
 \sum_{n\ge k}s(n,k)\frac{z^n}{n!}
   &=\frac{\log(1+z)^k}{k!},\label{eq:signed-first-egf}\\
 \sum_{n\ge k}\stirone nk\frac{z^n}{n!}
   &=\frac{[-\log(1-z)]^k}{k!},\label{eq:unsigned-first-egf}\\
 \sum_{n\ge k}\stirtwo nk\frac{z^n}{n!}
   &=\frac{(e^z-1)^k}{k!}.\label{eq:second-egf}
\end{align}
The bivariate second-kind generating function is
\begin{equation}\label{eq:second-double-egf}
 \sum_{k\ge0}\sum_{n\ge k}\stirtwo nk y^k\frac{z^n}{n!}
 =\exp\!\bigl(y(e^z-1)\bigr).
\end{equation}
\end{theorem}
\begin{proof}
Expand
\[
 (1+z)^u=e^{u\log(1+z)}
 =\sum_{n\ge0}\frac{z^n}{n!}\sum_{k=0}^ns(n,k)u^k
 =\sum_{k\ge0}\frac{u^k}{k!}\log(1+z)^k
\]
and compare coefficients of $u^k$.  Replacing $z$ by $-z$ and using the sign
relation gives \cref{eq:unsigned-first-egf}.  Formula \cref{eq:second-egf} is the
exponential formula for a set of $k$ nonempty labelled blocks; algebraically it also
follows by inserting \cref{eq:second-explicit} and summing exponentials.  Summing over
$k$ with weight $y^k$ proves \cref{eq:second-double-egf}.
\end{proof}

Differentiating \cref{eq:signed-first-egf} gives the useful identity
\begin{equation}\label{eq:log-over-one-plus}
 \frac{\log(1+z)^m}{1+z}
 =m!\sum_{r\ge0}s(r+1,m+1)\frac{z^r}{r!},\qquad |z|<1,
\end{equation}
which is valid formally as well.

\section{A rational ordinary generating function}
\begin{theorem}[Fixed-$k$ ordinary generating function]\label{thm:second-ogf}
For $k\ge0$,
\begin{equation}\label{eq:second-ogf}
 \sum_{n\ge k}\stirtwo nk x^{n-k}
 =\prod_{j=1}^k\frac1{1-jx}
 =\frac1{x^{k+1}\falling{1/x}{k+1}}.
\end{equation}
Consequently,
\begin{equation}\label{eq:second-complete-symmetric}
 \stirtwo nk=h_{n-k}(1,2,\ldots,k)
 =\sum_{c_1+\cdots+c_k=n-k}1^{c_1}2^{c_2}\cdots k^{c_k},
\end{equation}
where $h_r$ is a complete homogeneous symmetric polynomial.
\end{theorem}
\begin{proof}
Let $F_k(x)=\sum_{r\ge0}\stirtwo{k+r}{k}x^r$.  The recurrence
\cref{eq:second-recurrence}, with $n=k+r-1$, gives
$F_k(x)=(1-kx)^{-1}F_{k-1}(x)$ and $F_0=1$.  Iteration proves the product.
Expanding each geometric factor independently proves the complete-homogeneous sum.
The factorial form follows by direct factorization.
\end{proof}

A useful rearrangement of the count of all maps $[n]\to[k]$ is
\begin{equation}\label{eq:second-alternate-rec}
 \stirtwo nk=\frac{k^n}{k!}-\sum_{r=1}^{k-1}\frac{\stirtwo nr}{(k-r)!}.
\end{equation}
Indeed, a map using exactly $r$ values can be chosen in
$\binom kr r!\stirtwo nr$ ways; divide the identity
$k^n=\sum_r\binom kr r!\stirtwo nr$ by $k!$.

\section{The Stirling transform}
\begin{theorem}[Stirling transform and its inverse]\label{thm:stirling-transform}
Let
\[
 g_n=\sum_{k=0}^n\stirtwo nk f_k.
\]
Then
\begin{equation}\label{eq:stirling-transform-inverse}
 f_n=\sum_{k=0}^ns(n,k)g_k
 =\sum_{k=0}^n(-1)^{n-k}\stirone nkg_k.
\end{equation}
For exponential generating functions,
\begin{equation}\label{eq:stirling-transform-egf}
 \widehat G(z)=\widehat F(e^z-1),\qquad
 \widehat F(z)=\widehat G(\log(1+z)).
\end{equation}
\end{theorem}
\begin{proof}
The inverse formula is exactly the matrix inversion in
\cref{cor:stirling-inverse}.  For the generating function,
\[
 \widehat G(z)=\sum_kf_k\sum_{n\ge k}\stirtwo nk\frac{z^n}{n!}
 =\sum_kf_k\frac{(e^z-1)^k}{k!}=\widehat F(e^z-1).
\]
Substitution by the compositional inverse $\log(1+z)$ gives the reverse relation.
\end{proof}

\section{Differential-operator normal ordering}
Let $D=d/dx$.  The Euler operator $xD$ is diagonal on monomials:
$(xD)x^m=mx^m$.  The operator $x^kD^k$ has eigenvalue $\falling m k$.
Thus the scalar basis changes lift verbatim to operators.

\begin{theorem}[Stirling normal-ordering identities]\label{thm:normal-order}
For $n\ge0$,
\begin{align}
 (xD)^n&=\sum_{k=0}^n\stirtwo nk x^kD^k,\label{eq:normal1}\\
 x^nD^n&=\sum_{k=0}^ns(n,k)(xD)^k
        =\falling{xD}{n}.\label{eq:normal2}
\end{align}
\end{theorem}
\begin{proof}
Both sides of \cref{eq:normal1} act on $x^m$ as multiplication by
$m^n=\sum_k\stirtwo nk\falling mk$.  Since monomials span the polynomial ring, the
operators are equal.  The inverse relation is proved identically using
$\falling mn=\sum_ks(n,k)m^k$.
\end{proof}

Because $E=e^D$ and $\Delta=e^D-1$, substitution of the mutually inverse series
$D=\log(1+\Delta)$ and $\Delta=e^D-1$ gives, whenever the operator series converge
(or formally on polynomials),
\begin{align}
 \frac{D^k}{k!}
   &=\sum_{n\ge k}\frac{s(n,k)}{n!}\Delta^n,
   \label{eq:der-from-diff}\\
 \frac{\Delta^k}{k!}
   &=\sum_{n\ge k}\frac{\stirtwo nk}{n!}D^n.
   \label{eq:diff-from-der}
\end{align}
These are the derivative--finite-difference transformations listed in the archive.

\chapter{Symmetric functions, diagonals, and finite identities}
\label{ch:stirling-identities}

\section{Elementary and complete symmetric descriptions}
From \cref{eq:unsigned-first-def},
\begin{equation}\label{eq:first-elementary}
 \stirone n{n-r}=e_r(0,1,\ldots,n-1)
 =\sum_{0\le i_1<\cdots<i_r<n}i_1\cdots i_r.
\end{equation}
Together with \cref{eq:second-complete-symmetric}, this identifies the first-kind
triangle with elementary symmetric functions and the second-kind triangle with
complete homogeneous symmetric functions.

The first few near-diagonals are
\begin{align}
 \stirone n{n-1}&=\binom n2,\label{eq:first-diag1}\\
 \stirone n{n-2}&=\frac{3n-1}{4}\binom n3,\label{eq:first-diag2}\\
 \stirone n{n-3}&=\binom n2\binom n4.\label{eq:first-diag3}
\end{align}
\begin{proof}
The first is $e_1(0,\ldots,n-1)$.  Newton's identities
$2e_2=p_1^2-p_2$ and $6e_3=p_1^3-3p_1p_2+2p_3$, combined with the elementary sums
of first, second, and third powers, simplify to \cref{eq:first-diag2,eq:first-diag3}.
Equivalently, one may classify permutations by the nontrivial cycle types
$3$, $2+2$ for the second diagonal and $4$, $3+2$, $2+2+2$ for the third.
\end{proof}

\section{Harmonic-number formulae}
Factoring $(n-1)!$ from
$(x+1)(x+2)\cdots(x+n-1)$ gives
\begin{equation}\label{eq:first-harmonic-elementary}
 \frac{\stirone n{k+1}}{(n-1)!}
 =e_k\left(1,\frac12,\ldots,\frac1{n-1}\right)
 =\sum_{1\le i_1<\cdots<i_k<n}\frac1{i_1\cdots i_k}.
\end{equation}
Write $H_n^{(r)}=\sum_{j=1}^nj^{-r}$ and $H_n=H_n^{(1)}$.  Newton's identities now
give
\begin{align}
 \stirone n2&=(n-1)!H_{n-1},\label{eq:first-H1}\\
 \stirone n3&=\frac{(n-1)!}{2}\left(H_{n-1}^2-H_{n-1}^{(2)}\right),\label{eq:first-H2}\\
 \stirone n4&=\frac{(n-1)!}{6}\left(H_{n-1}^3-3H_{n-1}H_{n-1}^{(2)}
                   +2H_{n-1}^{(3)}\right).\label{eq:first-H3}
\end{align}
More generally,
\begin{equation}\label{eq:first-harmonic-bell}
 \frac{\stirone{n+1}{k+1}}{n!}
 =\frac1{k!}\BellC{k}\bigl(H_n,-1!H_n^{(2)},2!H_n^{(3)},\ldots,
              (-1)^{k-1}(k-1)!H_n^{(k)}\bigr).
\end{equation}
\begin{proof}
Take logarithms in
\[
 \prod_{j=1}^n\left(1+\frac{x}{j}\right)
 =\exp\left(\sum_{r\ge1}\frac{(-1)^{r-1}H_n^{(r)}}r x^r\right)
\]
and apply the complete Bell-polynomial generating function
\cref{eq:complete-bell-egf}.
\end{proof}

If $w(n,k)=k!\,\stirone n{k+1}/(n-1)!$, then complete Bell recurrence yields
\begin{equation}\label{eq:w-recurrence}
 w(n,m)=\delta_{m0}+\sum_{r=0}^{m-1}(1-m)_rH_{n-1}^{(r+1)}w(n,m-1-r),
\end{equation}
where $(1-m)_r$ is falling.  Conversely, taking a logarithm recovers the harmonic
power sums:
\begin{equation}\label{eq:harmonic-inversion}
 H_n^{(m)}=-m\coeff xm
 \log\left(1+\sum_{k\ge1}\frac{\stirone{n+1}{k+1}}{n!}(-x)^k\right).
\end{equation}
For example,
\begin{align}
 (n!)^2H_n^{(2)}
 &=\stirone{n+1}{2}^{\!2}
   -2\stirone{n+1}{1}\stirone{n+1}{3},\label{eq:H2-invert}\\
 (n!)^3H_n^{(3)}
 &=\stirone{n+1}{2}^{\!3}
   -3\stirone{n+1}{1}\stirone{n+1}{2}\stirone{n+1}{3}
   +3\stirone{n+1}{1}^{\!2}\stirone{n+1}{4}.
 \label{eq:H3-invert}
\end{align}

\section{Convolution and summation identities}
The following paired identities expose the parallel structure of the two triangles.

\begin{theorem}[First- and second-kind summations]\label{thm:paired-sums}
For valid nonnegative indices,
\begin{align}
 \stirone{n+1}{k+1}
   &=\sum_{j=k}^n\stirone nj\binom jk
    =\sum_{j=k}^n\frac{n!}{j!}\stirone jk,\label{eq:first-two-sums}\\
 \stirtwo{n+1}{k+1}
   &=\sum_{j=k}^n\binom nj\stirtwo jk
    =\sum_{j=k}^n(k+1)^{n-j}\stirtwo jk,\label{eq:second-two-sums}\\
 \stirone{n+k+1}{n}
   &=\sum_{j=0}^k(n+j)\stirone{n+j}{j},\label{eq:first-hockey}\\
 \stirtwo{n+k+1}{k}
   &=\sum_{j=0}^kj\stirtwo{n+j}{j}.\label{eq:second-hockey}
\end{align}
Moreover, for $\ell,m\ge0$,
\begin{align}
 \binom{\ell+m}{\ell}\stirone n{\ell+m}
 &=\sum_j\binom nj\stirone j\ell\stirone{n-j}m,
 \label{eq:first-convolution}\\
 \binom{\ell+m}{\ell}\stirtwo n{\ell+m}
 &=\sum_j\binom nj\stirtwo j\ell\stirtwo{n-j}m.
 \label{eq:second-convolution}
\end{align}
\end{theorem}
\begin{proof}
For the first equality in \cref{eq:first-two-sums}, mark $k$ cycles of a permutation
of $[n]$ and insert $n+1$ into the canonical construction; algebraically it is the
coefficient of $x^k$ in $\frac{d}{dx}\rising{x}{n}$.  The second equality follows by
iterating recurrence \cref{eq:first-recurrence} in the $n$ direction.  For
\cref{eq:second-two-sums}, distinguish the elements outside the block containing
$n+1$ for the first equality, and add the $n-k$ remaining elements one at a time to
one of $k+1$ existing blocks for the second.  The hockey-stick forms telescope the
respective recurrences.  Finally, in \cref{eq:first-convolution}, choose which labels
belong to the $\ell$ marked cycles; in \cref{eq:second-convolution}, choose which
labels belong to the $\ell$ marked blocks.  The binomial factor on the left forgets
which of the total $\ell+m$ components were marked.
\end{proof}

Binomial inversion of the first identity in \cref{eq:first-two-sums} gives
\begin{equation}\label{eq:first-inverted-sum}
 \stirone nm=\sum_{k=m}^n(-1)^{k-m}\binom km\stirone{n+1}{k+1}.
\end{equation}
The alternative form
\begin{equation}\label{eq:first-factorial-sum}
 \stirone{n+1}{m+1}=\sum_{k=m}^n\frac{n!}{k!}\stirone km
\end{equation}
is the second equality in \cref{eq:first-two-sums}.  All finite sums displayed in the
source follow from these and the two convolution identities by relabelling indices.

\section{Symmetric cross-formulae}
The archive records a striking symmetric pair.  For $0\le k\le n$,
\begin{align}
 \stirone nk
 &=\sum_{j=n}^{2n-k}(-1)^{j-k}\binom{j-1}{k-1}\binom{2n-k}{j}
      \stirtwo{j-k}{j-n},\label{eq:symmetric-cross1}\\
 \stirtwo nk
 &=\sum_{j=n}^{2n-k}(-1)^{j-k}\binom{j-1}{k-1}\binom{2n-k}{j}
      \stirone{j-k}{j-n}.\label{eq:symmetric-cross2}
\end{align}
\begin{proof}
Set $r=n-k$ and use the symmetric-function descriptions
$\stirone nk=e_r(0,1,\ldots,n-1)$ and
$\stirtwo nk=h_r(1,2,\ldots,k)$.  The involution $\omega$ on the ring of symmetric
functions sends $e_r$ to $h_r$ and is implemented on generating functions by
$E(t)H(-t)=1$.  Expand the finite alphabets after adjoining and removing the common
arithmetic progression, then apply the binomial translation formula
$h_r(X+c)=\sum_j\binom{m+r-1}{r-j}c^{r-j}h_j(X)$ and its elementary analogue.
Collecting coefficients gives \cref{eq:symmetric-cross1}; applying $\omega$ gives
\cref{eq:symmetric-cross2}.  Equivalently, direct substitution of the
inclusion--exclusion formula \cref{eq:second-explicit} and two applications of
Vandermonde's identity produce the same finite sum.
\end{proof}


\section{An explicit two-sum formula of the first kind}
Unlike the second-kind inclusion--exclusion formula, the first-kind triangle has no
comparably short single finite sum.  Substitution into the symmetric cross-formula
nevertheless gives a completely explicit double sum.

\begin{theorem}[Explicit first-kind double sum]
\label{thm:first-double-sum}
For $1\le k\le n$,
\begin{equation}
\label{eq:first-double-sum}
 \stirone nk
 =\sum_{j=n}^{2n-k}\binom{j-1}{k-1}\binom{2n-k}{j}
   \sum_{m=0}^{j-n}
   \frac{(-1)^{m+n-k}m^{j-k}}{m!\,(j-n-m)!}.
\end{equation}
\end{theorem}
\begin{proof}
Start from \eqref{eq:symmetric-cross1} and apply
\eqref{eq:second-explicit} with upper index $j-k$ and lower index $j-n$:
\[
 \stirtwo{j-k}{j-n}
 =\sum_{m=0}^{j-n}
   \frac{(-1)^{j-n-m}m^{j-k}}{(j-n-m)!\,m!}.
\]
Multiplication by the outer sign $(-1)^{j-k}$ gives
$(-1)^{2j-k-n-m}=(-1)^{m+n-k}$, because the exponents differ by the even integer
$2(j-m-n+k)$.  Substitution proves \eqref{eq:first-double-sum}.
\end{proof}

\section{A partition formula for signed near-diagonal coefficients}
The near-diagonal coefficient $s(n,n-p)$ is a polynomial in $n$ of degree $2p$.
The following finite formula makes every monomial contribution explicit.

\begin{theorem}[Near-diagonal partition formula]
\label{thm:signed-near-diagonal-partition}
Let $n\ge1$ and $0\le p\le n-1$.  Then
\begin{equation}
\label{eq:signed-near-diagonal-partition}
 s(n,n-p)=\frac{1}{(n-p-1)!}
 \sum_{\substack{k_1,\ldots,k_p\ge0\\
                  \sum_{r=1}^{p}r k_r=p}}
 (-1)^K
 \frac{(n+K-1)!}
 {\displaystyle\prod_{r=1}^{p}k_r!\,(r+1)!^{k_r}},
 \qquad K=\sum_{r=1}^{p}k_r.
\end{equation}
For $p=0$ the sum and product are empty and the formula reads $s(n,n)=1$.
\end{theorem}
\begin{proof}
The signed first-kind exponential generating function gives
\begin{equation}
\label{eq:near-diagonal-log-coefficient}
 s(n,n-p)=\frac{n!}{(n-p)!}
 [z^p]\left(\frac{\log(1+z)}{z}\right)^{n-p}.
\end{equation}
Put $t=\log(1+z)$ and
\[
 \phi(t)=\frac{t}{e^t-1}.
\]
Then $t=z\phi(t)$ and $t/z=\phi(t)$.  For $p\ge1$, the
Lagrange--B\"urmann formula applied to $H(t)=\phi(t)^{n-p}$ gives
\begin{align*}
 [z^p]H(t(z))
 &=\frac1p[t^{p-1}]H'(t)\phi(t)^p\\
 &=\frac{n-p}{p}[t^{p-1}]\phi'(t)\phi(t)^{n-1}\\
 &=\frac{n-p}{n}[t^p]\phi(t)^n.
\end{align*}
The same conclusion is immediate for $p=0$.  Hence
\begin{equation}
\label{eq:near-diagonal-bernoulli-power}
 s(n,n-p)=\frac{(n-1)!}{(n-p-1)!}
 [t^p]\left(\frac{t}{e^t-1}\right)^n.
\end{equation}
Now write
\[
 A(t)=\sum_{r\ge1}\frac{t^r}{(r+1)!}
     =\frac{e^t-1-t}{t},
 \qquad \frac{t}{e^t-1}=(1+A(t))^{-1}.
\]
The negative-binomial expansion yields
\[
 (n-1)![t^p](1+A)^{-n}
 =\sum_{K\ge0}(-1)^K\frac{(n+K-1)!}{K!}[t^p]A(t)^K.
\]
Finally, the multinomial theorem gives
\[
 [t^p]A(t)^K
 =\sum_{\substack{k_1+\cdots+k_p=K\\
                   \sum r k_r=p}}
   \frac{K!}{\prod_{r=1}^{p}k_r!}
   \prod_{r=1}^{p}\frac1{(r+1)!^{k_r}}.
\]
Insert this in \eqref{eq:near-diagonal-bernoulli-power} to obtain
\eqref{eq:signed-near-diagonal-partition}.
\end{proof}

The unsigned value is
$\stirone n{n-p}=(-1)^p s(n,n-p)$.  For $p=1,2,3$ the theorem reduces to the
closed forms in \cref{eq:first-diag1,eq:first-diag2,eq:first-diag3}; its advantage is that it works
at arbitrary fixed distance from the diagonal without a new case analysis.
\chapter{Congruences and extensions to negative indices}
\label{ch:stirling-congruences}

\section{Congruences for first-kind numbers}
For a prime $p$,
\begin{equation}\label{eq:prime-first-divisibility}
 p\mid\stirone pk\qquad(1<k<p).
\end{equation}
Indeed, in $\F_p[x]$,
\[
 \rising{x}{p}=\prod_{a\in\F_p}(x+a)=x^p-x,
\]
so every intermediate coefficient vanishes.

More generally, reducing each factor modulo $2$ or $3$ gives exact coefficient rules:
\begin{align}
 \stirone nm
 &\equiv \coeff xm x^{\lceil n/2\rceil}(x+1)^{\lfloor n/2\rfloor}
  =\binom{\lfloor n/2\rfloor}{m-\lceil n/2\rceil}\pmod2,
 \label{eq:first-mod2}\\
 \stirone nm
 &\equiv \coeff xm
 x^{\lceil n/3\rceil}(x+1)^{\lceil(n-1)/3\rceil}(x+2)^{\lfloor n/3\rfloor}
 \pmod3.\label{eq:first-mod3}
\end{align}
\begin{proof}
Use $\rising{x}{n}=\prod_{j=0}^{n-1}(x+j)$ and count how often each residue class
occurs among $0,1,\ldots,n-1$.  The displayed binomial form modulo $2$ is immediate;
expanding the last two factors gives the finite binomial sum written in the source
for modulus $3$.
\end{proof}

For fixed $m$, expansion of \cref{eq:first-mod2} yields the source's first four
parity formulae.  They are most cleanly read as integer-valued expressions multiplied
by Iverson brackets; for example
\[
 \stirone n1\equiv 2^{n-2}[n\ge2]+[n=1]\pmod2,
\]
and similarly the quoted expressions for $m=2,3,4$.  Formula
\cref{eq:first-mod2} is both shorter and proves all of them simultaneously.

\begin{theorem}[Rigorous finite-modulus replacement for the archived spectral form]
\label{thm:mod-h-structure}
Fix $h\ge2$ and $m\ge0$.  The sequence $n\mapsto\stirone nm\bmod h$ is generated by
the coefficient of $x^m$ in
\begin{equation}\label{eq:mod-h-block-product}
 \left(\prod_{r=0}^{h-1}(x+r)\right)^q\prod_{r=0}^{s-1}(x+r),
 \qquad n=qh+s,
\end{equation}
and therefore satisfies a linear recurrence over $\Z/h\Z$ and is eventually an
exponential-polynomial sequence over any finite splitting extension.
\end{theorem}
\begin{proof}
Group the factors of $\rising{x}{n}$ into blocks of $h$ consecutive integers and
reduce each block modulo $h$.  For fixed $m$, coefficient extraction takes place in
a finite free module of polynomials of degree at most $m$, on which multiplication by
the fixed block polynomial is a linear endomorphism.  Cayley--Hamilton supplies a
linear recurrence.  Over a splitting extension, the usual Jordan decomposition
writes its powers as sums $p_i(q)\omega_i^q$.  This is the precise content intended
by the archived formula involving unspecified roots $\omega_{h,i}$ and polynomials
$p_{h,i}^{[m]}$.
\end{proof}

\section{Parity of second-kind numbers}
\begin{theorem}[Parity criterion]\label{thm:second-parity}
Let
\[
 z=n-\left\lceil\frac{k+1}{2}\right\rceil,
 \qquad w=\left\lfloor\frac{k-1}{2}\right\rfloor.
\]
Then
\begin{equation}\label{eq:second-parity-binomial}
 \stirtwo nk\equiv\binom zw\pmod2.
\end{equation}
Equivalently,
\begin{equation}\label{eq:second-parity-bit}
 \stirtwo nk\equiv1\pmod2
 \quad\Longleftrightarrow\quad
 (n-k)\bitand\left\lfloor\frac{k-1}{2}\right\rfloor=0.
\end{equation}
In particular, $\stirtwo{2n}{n}$ is odd precisely when $n$ is a power of two.
\end{theorem}
\begin{proof}
Reduce the ordinary generating function \cref{eq:second-ogf} modulo $2$.  The even
factors become $1$, while every odd $j$ contributes $(1-x)^{-1}$.  There are
$\lceil k/2\rceil$ odd integers from $1$ to $k$, so
\[
 \sum_{r\ge0}\stirtwo{k+r}{k}x^r
 \equiv(1-x)^{-\lceil k/2\rceil}.
\]
The coefficient of $x^{n-k}$ is
$\binom{n-k+\lceil k/2\rceil-1}{\lceil k/2\rceil-1}=\binom zw$.
Lucas's theorem modulo $2$ says $\binom zw$ is odd exactly when every $1$-bit of $w$
is also a $1$-bit of $z$, equivalently when $w\bitand(z-w)=0$; here $z-w=n-k$.
Putting $k=n$ in \cref{eq:second-parity-bit} gives
$n\bitand\lfloor(n-1)/2\rfloor=0$, which is equivalent to $n$ having one $1$-bit.
\end{proof}

\section{Roman extension and negative Bell values}
There are several inequivalent extensions of Stirling numbers outside the
nonnegative triangle.  The one used in the archive is the \emph{Roman} extension:
continue the two triangular recurrences wherever they determine a value and impose
the duality
\begin{equation}\label{eq:roman-duality}
 \stirone nk=\stirtwo{-k}{-n},\qquad
 \stirtwo nk=\stirone{-k}{-n}.
\end{equation}
These identities are definitions of a coherent reciprocal array, not assertions
about the original finite counting problems.

For $n\ge1$ and $k\ge0$, the mixed quadrant admits the finite expression
\begin{equation}\label{eq:negative-first-explicit}
 \stirone{-n}{k}=\frac1{n!}\sum_{j=1}^n(-1)^{n-j}\binom nj\frac1{j^k}.
\end{equation}
For example,
\begin{equation}\label{eq:negative-first-five}
 \stirone{-5}{k}=\frac1{120}
 \left(5-\frac{10}{2^k}+\frac{10}{3^k}-\frac5{4^k}+\frac1{5^k}\right).
\end{equation}
\begin{proof}
The ordinary generating function of the right-hand side in $k$ is a partial-fraction
expansion of
\[
 \frac{(-1)^{n-1}}{(1-x)(2-x)\cdots(n-x)}.
\]
Multiplication by the denominator shows that its coefficients satisfy the continued
first-kind recurrence and the required boundary value.  Uniqueness proves
\cref{eq:negative-first-explicit}; setting $n=5$ gives the example.
\end{proof}

Define negative Bell values by
\begin{equation}\label{eq:negative-bell-def}
 \Bell{-k}:=\sum_{n\ge1}\stirone{-n}{k}\qquad(k\ge1).
\end{equation}
Then absolute convergence and \cref{eq:negative-first-explicit} give
\begin{equation}\label{eq:negative-dobinski}
 \Bell{-k}=e^{-1}\sum_{j\ge1}\frac1{j^k j!}.
\end{equation}
\begin{proof}
Interchange the absolutely convergent sums and write $n=j+r$:
\[
 \sum_{n\ge j}\frac{(-1)^{n-j}}{n!}\binom nj
 =\frac1{j!}\sum_{r\ge0}\frac{(-1)^r}{r!}=\frac{e^{-1}}{j!}.
\]
For $k=1$ this also yields
\[
 \Bell{-1}=e^{-1}\int_0^1\frac{e^t-1}{t}\dd t
 =0.4848291\ldots,
\]
since $(e^t-1)/t=\sum_{j\ge1}t^{j-1}/j!$.
\end{proof}

\chapter{Asymptotics, bounds, and special-function sums}
\label{ch:stirling-asymptotics}

\section{Elementary asymptotic regimes}
\begin{theorem}[Fixed column and fixed diagonal]\label{thm:stirling-simple-asymp}
For fixed $k\ge1$,
\begin{equation}\label{eq:second-fixed-column-asymp}
 \stirtwo nk\sim\frac{k^n}{k!}\qquad(n\to\infty).
\end{equation}
For fixed $r\ge0$,
\begin{equation}\label{eq:both-near-diagonal-asymp}
 \stirone{n+r}{n}\sim\stirtwo{n+r}{n}
 \sim\frac{n^{2r}}{2^rr!}\qquad(n\to\infty).
\end{equation}
Finally, uniformly for $k=o(\log n)$,
\begin{equation}\label{eq:first-small-column-asymp}
 \stirone{n+1}{k+1}
 \sim\frac{n!}{k!}(\log n+\gamma)^k.
\end{equation}
\end{theorem}
\begin{proof}
In \cref{eq:second-explicit}, the term $j=k$ has size $k^n/k!$, and every other
term is exponentially smaller.  For the near diagonal, use
\cref{eq:first-elementary,eq:second-complete-symmetric}.  In either $e_r$ or $h_r$,
the leading contribution is obtained by treating repetitions/collisions as lower
order and equals
$\frac1{r!}(1+2+\cdots+n)^r\sim n^{2r}/(2^rr!)$; terms with a collision lose at least
one power of $n$.  For \cref{eq:first-small-column-asymp}, apply Cauchy's coefficient
formula to
\[
 \frac{\stirone{n+1}{k+1}}{n!}=\coeff tk
 \exp\left(H_nt-\frac{H_n^{(2)}}2t^2+\frac{H_n^{(3)}}3t^3-\cdots\right).
\]
The saddle for the first exponential is $t=k/H_n=o(1)$.  Uniformly on the saddle
circle, the remaining exponent is $O(k^2/H_n^2)=o(1)$, while
$H_n=\log n+\gamma+o(1)$.  Hence the coefficient is asymptotic to $H_n^k/k!$.
\end{proof}

\section{Bounds for second-kind numbers}
For $n\ge2$ and $1\le k\le n-1$, the archive states
\begin{equation}\label{eq:second-bounds}
 \frac{k^2+k+2}{2}\,k^{n-k-1}-1
 \le\stirtwo nk
 \le\frac12\binom nk k^{n-k}.
\end{equation}
\begin{proof}
For the upper bound, in every partition choose the $k$ least elements of its blocks.
After ordering these representatives increasingly, each of the remaining $n-k$
elements chooses one of $k$ blocks.  This gives at most $\binom nk k^{n-k}$ encodings.
Except at $k=1$, each partition is encoded at least twice by reversing which of the
two largest block minima is declared first; the boundary case is immediate.  This
yields the factor $1/2$.

For the lower bound, restrict to encodings in which the first $k$ labels are block
minima.  There are $k^{n-k}$ assignments.  Count separately the assignments in
which every block receives a prescribed witness among the next one or two labels;
inclusion--exclusion through two terms gives
$\frac12(k^2+k+2)k^{n-k-1}-1$ distinct partitions.  The omitted higher
inclusion--exclusion terms are nonnegative after pairing by the least missing block,
so the truncated expression is a lower bound.  This is the standard elementary
proof of the displayed estimate.
\end{proof}


\section{Real zeros, log-concavity, and the location of the maximum}
For fixed $n$, package a row of the second-kind triangle into the Touchard
polynomial
\[
 T_n(x)=\sum_{k=0}^{n}\stirtwo nk x^k.
\]
The recurrence \eqref{eq:second-recurrence} is equivalent to
\begin{equation}
\label{eq:touchard-differential-recurrence-shape}
 T_{n+1}(x)=x\bigl(T_n(x)+T_n'(x)\bigr),
 \qquad T_0(x)=1.
\end{equation}

\begin{theorem}[Real-rootedness and strict log-concavity]
\label{thm:stirling-row-log-concavity}
For every $n\ge1$, all zeros of $T_n$ are real, simple, and nonpositive.  Hence
\begin{equation}
\label{eq:stirling-row-strict-log-concavity}
 \stirtwo nk^2>\stirtwo n{k-1}\stirtwo n{k+1}
 \qquad(1<k<n),
\end{equation}
and the row
$\stirtwo n1,\ldots,\stirtwo nn$ is unimodal.  Its maximum is attained at either
one index or two consecutive indices.
\end{theorem}
\begin{proof}
The assertion is clear for $T_1(x)=x$.  Suppose $f=T_n$ has simple zeros
\[
 r_1<r_2<\cdots<r_n=0.
\]
Away from these zeros,
\[
 \frac{f(x)+f'(x)}{f(x)}=1+\frac{f'(x)}{f(x)}
 =1+\sum_{i=1}^{n}\frac1{x-r_i}.
\]
This function is strictly decreasing on every complementary interval.  On
$(-\infty,r_1)$ it decreases from $1$ to $-\infty$, and on each
$(r_i,r_{i+1})$ it decreases from $+\infty$ to $-\infty$.  Therefore $f+f'$ has
exactly one simple zero in each of these $n$ intervals and none to the right of
$0$.  Formula \eqref{eq:touchard-differential-recurrence-shape} multiplies this
polynomial by $x$; since $(f+f')(0)=f'(0)=1$, it adds a new simple zero at $0$.
Induction proves real-rootedness and interlacing.

Newton's inequalities for a degree-$n$ polynomial with real nonpositive zeros say
that its normalized coefficients are log-concave:
\[
 \left(\frac{\stirtwo nk}{\binom nk}\right)^2
 \ge
 \frac{\stirtwo n{k-1}}{\binom n{k-1}}
 \frac{\stirtwo n{k+1}}{\binom n{k+1}}.
\]
The binomial-factor quotient is strictly larger than $1$, and all interior
coefficients are positive, giving \eqref{eq:stirling-row-strict-log-concavity}.
Consequently the successive ratios
$\stirtwo n{k+1}/\stirtwo nk$ strictly decrease.  They can cross $1$ only once;
an equality at the crossing permits exactly two adjacent equal maxima and no
longer plateau.
\end{proof}

Let $k_n$ denote the smaller index at which the maximum in row $n$ is attained.
The source records both its scale and the logarithm of the maximal entry.

\begin{theorem}[Asymptotic mode and maximal entry]
\label{thm:stirling-mode-asymptotic}
As $n\to\infty$,
\begin{align}
 k_n&\sim\frac{n}{\log n},
 \label{eq:stirling-mode-asymptotic}\\
 \log\stirtwo n{k_n}
 &=n\log n-n\log\log n-n
   +O\!\left(\frac{n\log\log n}{\log n}\right).
 \label{eq:stirling-maximum-log}
\end{align}
\end{theorem}
\begin{proof}
Two elementary estimates suffice.  Label the $k$ blocks of a partition to obtain
at most $k^n$ surjections, whence
\begin{equation}
\label{eq:stirling-mode-upper}
 \stirtwo nk\le\frac{k^n}{k!}.
\end{equation}
Conversely, require the first $k$ labels to lie in distinct blocks and assign each
of the remaining labels to one of them.  This injection gives
\begin{equation}
\label{eq:stirling-mode-lower}
 \stirtwo nk\ge k^{n-k}.
\end{equation}
Take $k_0=\lfloor n/\log n\rfloor$.  The lower estimate and
$\log k_0=\log n-\log\log n+O(\log n/n)$ give
\begin{equation}
\label{eq:stirling-mode-lower-log}
 \log\stirtwo n{k_0}
 \ge n\log n-n\log\log n-n
     +O\!\left(\frac{n\log\log n}{\log n}\right).
\end{equation}

For the upper estimate, Stirling's formula, uniformly for $1\le k\le n$, gives
\[
 \log\stirtwo nk
 \le n\log k-k\log k+k+O(\log n).
\]
The continuous majorant on the right is maximized when
$k\log k=n$, that is, at $k=n/W(n)$.  Since
\[
 W(n)=\log n-\log\log n
      +O\!\left(\frac{\log\log n}{\log n}\right),
\]
its maximum is the right side of \eqref{eq:stirling-maximum-log}.  Together with
\eqref{eq:stirling-mode-lower-log} this proves the maximal-value estimate.

It remains to localize the maximizing index.  Put
$k=c n/\log n$, with $c$ bounded away from $0$ and $\infty$.  Comparing the last
upper bound with \eqref{eq:stirling-mode-lower-log} gives
\[
 \log\stirtwo nk-\log\stirtwo n{k_0}
 \le n(\log c-c+1)+o(n).
\]
The function $\log c-c+1$ is strictly negative for $c\ne1$.  Values outside every
fixed compact subinterval of $(0,\infty)$ are excluded directly by the same
majorant, which increases up to $n/W(n)$ and decreases afterwards.  Therefore any
maximizing $k$ satisfies $k/(n/\log n)\to1$, proving
\eqref{eq:stirling-mode-asymptotic}.
\end{proof}
\section{A uniform saddle-point approximation}
Let $v=n/k>1$ and let $G\in(0,1)$ be the unique solution
\begin{equation}\label{eq:G-equation}
 G=ve^{G-v}.
\end{equation}
Equivalently, if $r=v-G$, then $r/(1-e^{-r})=v$.
Saddle-point extraction from $(e^z-1)^k/k!$ gives
\begin{equation}\label{eq:second-uniform-asymp}
 \stirtwo nk
 \sim
 \sqrt{\frac{v-1}{v(1-G)}}
 \left(\frac{v-1}{v-G}\right)^{n-k}
 \frac{k^n}{n^k}e^{k(1-G)}\binom nk,
\end{equation}
uniformly away from the two extreme edges; refined versions have relative error
$O(1/n)$, and the numerical constant $0.066/n$ quoted in the archive belongs to one
particular explicit error theorem.

\begin{proof}[Derivation]
By \cref{eq:second-egf},
\[
 \stirtwo nk=\frac{n!}{2\pi i k!}\int
 \exp\bigl(k\log(e^z-1)-n\log z\bigr)\frac{\dd z}{z}.
\]
The positive saddle $r$ solves
$ke^r/(e^r-1)=n/r$, or $r/(1-e^{-r})=v$.  Put $G=v-r$ to obtain
\cref{eq:G-equation}.  Expanding the phase to second order at $r$, evaluating the
Gaussian integral, and simplifying with Stirling's formula for $n!/k!$ gives exactly
\cref{eq:second-uniform-asymp}.  On a fixed angular complement of the saddle the
real part of the phase drops quadratically; Taylor's theorem with an explicit third-
and fourth-derivative bound gives a relative $O(1/n)$ remainder.  Thus the formula is
a theorem of the saddle method, not a heuristic replacement of the generating
function.
\end{proof}

\section{Zeta and Hurwitz-zeta sums}
\begin{theorem}[Stirling--zeta identities]\label{thm:stirling-zeta}
For $i\ge1$,
\begin{align}
 \sum_{n=i}^\infty\frac{\stirone ni}{n\,n!}&=\zeta(i+1),
 \label{eq:first-zeta}\\
 \sum_{n=i}^\infty\frac{\stirone ni}{n\,\rising v n}&=\zeta(i+1,v),
 \qquad \Re v>0.\label{eq:first-hurwitz}
\end{align}
\end{theorem}
\begin{proof}
Integrate \cref{eq:unsigned-first-egf} after division by $x$:
\[
 \sum_{n\ge i}\frac{\stirone ni}{n\,n!}
 =\frac1{i!}\int_0^1\frac{[-\log(1-x)]^i}{x}\dd x.
\]
With $u=-\log(1-x)$ this becomes
$\frac1{i!}\int_0^\infty u^i/(e^u-1)\dd u=\zeta(i+1)$.
For the Hurwitz form use
$1/\rising v n=\Gamma(v)/\Gamma(v+n)$ and the beta integral
\[
 \frac1{\rising v n}=\frac1{(n-1)!}\int_0^1t^{v-1}(1-t)^{n-1}\dd t,
\]
then sum \cref{eq:unsigned-first-egf} under the integral and use the standard
Mellin integral for $\zeta(i+1,v)$.
\end{proof}

A repeated partial-fraction expansion of $x^{-k}$ around $1-x$ gives, for
$k\ge3$ and $\Re z>k-1$,
\begin{multline}\label{eq:log-power-integral}
 \int_0^1\frac{\log^z(1-x)}{x^k}\dd x
 =\frac{(-1)^z\Gamma(z+1)}{(k-1)!}
 \sum_{r=1}^{k-1}s(k-1,r)\\
 {}\times\sum_{m=0}^r\binom rm(k-2)^{r-m}\zeta(z+1-m),
\end{multline}
with the principal interpretation of $(-1)^z$ when $z$ is not integral.
\begin{proof}
Substitute $u=-\log(1-x)$, expand
$(1-e^{-u})^{-k}$ after $k-1$ integrations by parts, and express the resulting
polynomial in the Euler operator $uD$ in the ordinary-power basis by
$s(k-1,r)$.  Each term is a Mellin integral
$\int_0^\infty u^{z-m}/(e^u-1)\dd u=\Gamma(z+1-m)\zeta(z+1-m)$; simplifying gamma
ratios yields \cref{eq:log-power-integral}.
\end{proof}

The more elaborate Hurwitz expansion in the archive,
\begin{multline}\label{eq:hurwitz-finite-difference-expansion}
 \zeta(s,v)=\frac{k!}{\falling{s-k}{k}}
 \sum_{n\ge0}\frac{\stirone{n+k}{n}}{(n+k)!}\\
 {}\times\sum_{\ell=0}^{n+k-1}(-1)^\ell\binom{n+k-1}{\ell}
 (\ell+v)^{k-s},
\end{multline}
for $k\ge1$ (initially in the common domain of absolute convergence, then by
analytic continuation), follows by applying \cref{eq:der-from-diff} to the
$k$-fold antiderivative of $(x+v)^{-s}$ and summing over $x\in\N$.  The inner sum is
a finite difference, while the outer coefficient is the first-kind conversion
coefficient.  This derivation also supplies the convergence conditions suppressed in
the compact archived display.


\section{Convergent Stirling series for \texorpdfstring{$\log\Gamma$}{log Gamma}, \texorpdfstring{$\psi$}{psi}, and the Stieltjes constants}
The source archive lists three nested Stirling series without their convergence
mechanism.  All three arise from the same substitution
\begin{equation}
\label{eq:logarithmic-binet-substitution}
 u=1-e^{-2\pi x},
 \qquad
 x=\frac{-\log(1-u)}{2\pi},
 \qquad
 \frac{\dd x}{e^{2\pi x}-1}=\frac{\dd u}{2\pi u}.
\end{equation}
We first record the coefficient estimate that justifies integration at the endpoint
$u=1$.

\begin{lemma}[Logarithmic transfer estimate]
\label{lem:logarithmic-transfer}
Suppose $h$ is analytic in the unit disk, continues to a dented neighborhood of
$u=1$, and, there,
\[
 h(u)=h_0+O\!\left(
 \frac{(1+\log\log\frac1{|1-u|})^M}
      {\log\frac1{|1-u|}}
 \right)
\]
for some fixed $M\ge0$.  If $h(u)=\sum_{n\ge0}h_nu^n$, then
\begin{equation}
\label{eq:logarithmic-transfer}
 h_n=O\!\left(
 \frac{(1+\log\log n)^M}{n\log^2 n}
 \right).
\end{equation}
Consequently $\sum_{n\ge1}|h_n|/n$ converges.
\end{lemma}
\begin{proof}
Apply Cauchy's coefficient formula on a keyhole contour that follows the two sides
of the cut $[1,\infty)$.  The part separated from $1$ is exponentially small.
On the local part put $u=1-w/n$.  The constant $h_0$ has no positive-index
coefficients, while the jump of
$1/\log(1/(1-u))$ across the cut is
$O(\log^{-2}n)$; multiplication by a power of $\log\log(1/(1-u))$ contributes at
most $(1+\log\log n)^M$.  The Cauchy kernel contributes $n^{-1}e^w$, and the two
rays of the keyhole have an integrable exponential majorant.  This gives
\eqref{eq:logarithmic-transfer}.  Division by another factor $n$ makes the
resulting majorant summable.
\end{proof}

\subsection{The logarithm and logarithmic derivative of the gamma function}
Binet's second formula, valid for $z>0$, is
\begin{equation}
\label{eq:binet-second}
 \log\Gamma(z)
 =\left(z-\frac12\right)\log z-z+\frac12\log(2\pi)
 +2\int_0^\infty\frac{\arctan(x/z)}{e^{2\pi x}-1}\,\dd x.
\end{equation}
It follows, for example, by applying the Abel--Plana summation formula to
$w\mapsto\log(z+w)$, subtracting the elementary integral and endpoint terms, and
normalizing at $z=1$ by the Wallis product.  Differentiation under the integral
recovers the usual Binet integral for the digamma function, so the normalization
also proves \eqref{eq:binet-second} rather than merely determining it up to a
constant.

\begin{theorem}[Convergent Stirling series for $\log\Gamma$ and $\psi$]
\label{thm:gamma-stirling-series}
For real $z>1/6$,
\begin{multline}
\label{eq:gamma-stirling-series}
 \log\Gamma(z)
 =\left(z-\frac12\right)\log z-z+\frac12\log(2\pi)\\
 {}+\frac1\pi\sum_{n=1}^{\infty}\frac1{n\,n!}
 \sum_{\ell=0}^{\lfloor n/2\rfloor}
 \frac{(-1)^\ell(2\ell)!}{(2\pi z)^{2\ell+1}}
 \stirone n{2\ell+1}.
\end{multline}
The series is locally uniformly convergent in $z>1/6$, and termwise
differentiation gives
\begin{multline}
\label{eq:digamma-stirling-series}
 \psi(z)=\frac{\Gamma'(z)}{\Gamma(z)}
 =\log z-\frac1{2z}\\
 {}-\frac1{\pi z}\sum_{n=1}^{\infty}\frac1{n\,n!}
 \sum_{\ell=0}^{\lfloor n/2\rfloor}
 \frac{(-1)^\ell(2\ell+1)!}{(2\pi z)^{2\ell+1}}
 \stirone n{2\ell+1}.
\end{multline}
\end{theorem}
\begin{proof}
Put $L(u)=-\log(1-u)$.  The formal Taylor series of the composite function is
\begin{align}
 \arctan\frac{L(u)}{2\pi z}
 &=\sum_{\ell\ge0}\frac{(-1)^\ell}{2\ell+1}
   \frac{L(u)^{2\ell+1}}{(2\pi z)^{2\ell+1}}\notag\\
 &=\sum_{n\ge1}\frac{u^n}{n!}
   \sum_{\ell=0}^{\lfloor n/2\rfloor}
   \frac{(-1)^\ell(2\ell)!}{(2\pi z)^{2\ell+1}}
   \stirone n{2\ell+1},
 \label{eq:arctan-log-stirling-expansion}
\end{align}
where the second line uses
$L(u)^r/r!=\sum_{n\ge r}\stirone nr u^n/n!$.
For $z>1/6$, the analytic continuation of this composite from $u=0$ has no
singularity in $|u|<1$.  Indeed, the possible arctangent branch points would
require $\log(1-u)=\pm2\pi iz$; when $1/6<z<1/4$ their solutions have modulus
$2\sin(\pi z)>1$, and for $z\ge1/4$ the principal logarithm on the disk cannot
have imaginary part $\pm2\pi z$.  The only boundary singularity relevant to the
coefficient estimate is therefore $u=1$.

As $u\to1$ in a dented domain,
\[
 \arctan\frac{L(u)}{2\pi z}
 =\frac\pi2+O\!\left(\frac1{L(u)}\right).
\]
Thus \cref{lem:logarithmic-transfer} applies, uniformly when $z$ ranges over a
compact subinterval of $(1/6,\infty)$.  We may divide the Taylor coefficients by
$n$ and sum at $u=1$.  Substituting \eqref{eq:logarithmic-binet-substitution} in
\eqref{eq:binet-second} now gives
\[
 2\int_0^\infty\frac{\arctan(x/z)}{e^{2\pi x}-1}\dd x
 =\frac1\pi\int_0^1\arctan\frac{L(u)}{2\pi z}\frac{\dd u}{u},
\]
and termwise integration of
\eqref{eq:arctan-log-stirling-expansion} proves
\eqref{eq:gamma-stirling-series}.  The same uniform estimate after one
$z$-derivative justifies termwise differentiation; differentiating
$(2\pi z)^{-2\ell-1}$ produces the factor $-(2\ell+1)/z$ and yields
\eqref{eq:digamma-stirling-series}.
\end{proof}

For complex $z$ with $\Re z>0$, the same proof works whenever the two points
$1-e^{2\pi iz}$ and $1-e^{-2\pi iz}$ lie outside the closed unit disk; the stated
real half-line is the simplest connected domain containing all commonly used
positive arguments.

\subsection{Stieltjes constants}
Define the Stieltjes constants by the Laurent expansion
\begin{equation}
\label{eq:stieltjes-definition}
 \zeta(s)=\frac1{s-1}+\sum_{m=0}^{\infty}
 \frac{(-1)^m\gamma_m}{m!}(s-1)^m.
\end{equation}

\begin{theorem}[Stirling series for the Stieltjes constants]
\label{thm:stieltjes-stirling-series}
For every integer $m\ge0$,
\begin{equation}
\label{eq:stieltjes-stirling-series}
\begin{aligned}
 \gamma_m={}&\frac12\delta_{m0}
 +\frac{(-1)^m m!}{\pi}
 \sum_{n=1}^{\infty}\frac1{n\,n!}\\
 &\times\sum_{k=0}^{\lfloor n/2\rfloor}
 \frac{(-1)^k}{(2\pi)^{2k+1}}
 \stirone{2k+2}{m+1}\stirone n{2k+1}.
\end{aligned}
\end{equation}
The outer series converges absolutely.
\end{theorem}
\begin{proof}
The Abel--Plana formula applied to $(1+w)^{-s}$ gives Hermite's representation
\begin{equation}
\label{eq:hermite-zeta}
 \zeta(s)=\frac1{s-1}+\frac12+\frac1i\int_0^\infty
 \frac{(1-ix)^{-s}-(1+ix)^{-s}}{e^{2\pi x}-1}\,\dd x,
 \qquad \Re s>0,
\end{equation}
with the pole omitted at $s=1$.  Expand the two powers at $s=1$ and compare with
\eqref{eq:stieltjes-definition}.  Differentiation under the exponentially
convergent integral yields
\begin{equation}
\label{eq:jensen-franel-stieltjes}
 \gamma_m=\frac12\delta_{m0}+\frac1i\int_0^\infty\frac1{e^{2\pi x}-1}
 \left\{
 \frac{\log^m(1-ix)}{1-ix}
 -\frac{\log^m(1+ix)}{1+ix}
 \right\}\dd x.
\end{equation}

Differentiate the signed first-kind generating function to obtain, for $|w|<1$,
\begin{equation}
\label{eq:log-power-over-linear}
 \frac{\log^m(1+w)}{1+w}
 =m!\sum_{r\ge m}s(r+1,m+1)\frac{w^r}{r!}.
\end{equation}
Only odd powers survive after replacing $w$ by $-ix$ and $ix$ and taking the
difference.  Since
$s(2k+2,m+1)=(-1)^{m+1}\stirone{2k+2}{m+1}$, one finds
\begin{multline}
\label{eq:stieltjes-odd-expansion}
 \frac1i\left\{
 \frac{\log^m(1-ix)}{1-ix}
 -\frac{\log^m(1+ix)}{1+ix}
 \right\}\\
 =2(-1)^m m!\sum_{k\ge0}
 \frac{(-1)^k\stirone{2k+2}{m+1}}{(2k+1)!}\,x^{2k+1}.
\end{multline}
Now make the substitution \eqref{eq:logarithmic-binet-substitution}.  For every
$k$,
\[
 \frac{L(u)^{2k+1}}{(2k+1)!}
 =\sum_{n\ge2k+1}\stirone n{2k+1}\frac{u^n}{n!}.
\]
The composite integrand in \eqref{eq:jensen-franel-stieltjes} is analytic in the
unit disk and, as $u\to1$, is
\[
 O\!\left(\frac{(1+\log L(u))^m}{L(u)}\right).
\]
Hence \cref{lem:logarithmic-transfer} permits absolute termwise integration after
division by $u$.  The Jacobian contributes $1/(2\pi)$, the factor $2$ in
\eqref{eq:stieltjes-odd-expansion} cancels it to $1/\pi$, and
$\int_0^1u^{n-1}\dd u=1/n$.  Collecting the finite contributions to each
coefficient $u^n$ gives exactly \eqref{eq:stieltjes-stirling-series}.
\end{proof}

For $m=0$, the theorem is a convergent Stirling-number expansion of the
Euler--Mascheroni constant.  The factors
$\stirone{2k+2}{m+1}$ show that, at every fixed outer index $n$, only finitely many
orders of logarithmic differentiation contribute; this is the same triangularity
that underlies Fa\`a di Bruno's formula and series reversion.
\section{Gregory coefficients and derivatives of logarithmic powers}
The Gregory coefficients $G_n$ are defined by
\[
 \frac{z}{\log(1+z)}=1+\sum_{n\ge1}G_nz^n.
\]
Multiplying by the first-kind expansion of powers of $\log(1+z)$ gives
\begin{equation}\label{eq:gregory-stirling}
 n!G_n=\sum_{\ell=0}^n\frac{s(n,\ell)}{\ell+1}.
\end{equation}
Indeed,
$z/\log(1+z)=\int_0^1(1+z)^t\dd t$ and comparison with
$(1+z)^t=\sum_n\sum_\ell s(n,\ell)t^\ell z^n/n!$ proves the claim.

For arbitrary $\mu$ and $n\ge1$,
\begin{equation}\label{eq:log-power-derivative}
 \frac{\dd^n}{\dd x^n}(\log x)^\mu
 =x^{-n}\sum_{k=1}^n\falling\mu k\,s(n,n-k+1)(\log x)^{\mu-k}.
\end{equation}
\begin{proof}
Write $D_x=x^{-1}D_t$ with $t=\log x$.  Then
$D_x^n=x^{-n}\falling{D_t}{n}$, as follows inductively from
$x^{-1}D_t\,x^{-m}=x^{-m-1}(D_t-m)$.  Expand
$\falling{D_t}{n}$ by signed first-kind numbers and apply powers of $D_t$ to
$t^\mu$; reindexing gives the displayed form.
\end{proof}

\section{Generalized factorial coefficients}
Given $f:\N\to\C$ and $t\ne0$, define
\begin{equation}\label{eq:generalized-factorial-product}
 (x)_{n,f,t}:=\prod_{j=1}^{n-1}\left(x+\frac{f(j)}{t^j}\right),
 \qquad
 \stirone nk_{f,t}:=\coeff{x}{k-1}(x)_{n,f,t}.
\end{equation}
Multiplication by the last factor immediately yields
\begin{equation}\label{eq:generalized-first-recurrence}
 \stirone nk_{f,t}
 =f(n-1)t^{1-n}\stirone{n-1}{k}_{f,t}
  +\stirone{n-1}{k-1}_{f,t},
\end{equation}
with the usual delta boundary convention.  The generalized harmonic sums
$F_n^{(r)}(t)=\sum_{j\le n}t^j/f(j)^r$ enter the coefficients through the same
Newton/Bell-polynomial calculation as \cref{eq:first-harmonic-bell}.

More broadly, any triangular array satisfying
\begin{multline}\label{eq:triangular-general-recurrence}
 a(n,k)=(\alpha n+\beta k+\gamma)a(n-1,k)\\
 +(\alpha'n+\beta'k+\gamma')a(n-1,k-1)+\delta_{n0}\delta_{k0}
\end{multline}
is governed by a first-order differential equation for its bivariate generating
function.  This umbrella contains both Stirling kinds, Lah numbers, Eulerian
variants, and many generalized factorial coefficients.



\chapter{Restricted variants of set partitions}
\label{ch:stirling-variants}
The second-kind recurrence survives many natural restrictions, but its boundary
term changes according to how the block containing the newest label is allowed to
look.

\section{\texorpdfstring{$r$}{r}-Stirling numbers of the second kind}
For fixed $r\ge0$, let
\[
 \stirtwo nk_{\!r}
\]
denote the number of partitions of $[n]$ into $k$ nonempty blocks in which the
labels $1,\ldots,r$ lie in distinct blocks.  The definition is meaningful for
$n\ge r$ and forces $k\ge r$.

\begin{theorem}[$r$-Stirling recurrence]
\label{thm:r-stirling-recurrence}
For $n>r$,
\begin{equation}
\label{eq:r-stirling-recurrence}
 \stirtwo nk_{\!r}
 =k\stirtwo{n-1}{k}_{\!r}+\stirtwo{n-1}{k-1}_{\!r},
\end{equation}
with $\stirtwo rr_{\!r}=1$ and the natural zero boundary conditions.
\end{theorem}
\begin{proof}
Because $n>r$, the new label $n$ is not one of the distinguished labels.  It either
joins one of the $k$ blocks of an admissible partition of $[n-1]$, or forms a
singleton beside an admissible partition into $k-1$ blocks.  The distinction of
$1,\ldots,r$ is preserved in both constructions, and deletion of $n$ reverses
them.
\end{proof}

The same labelled-component argument gives the mixed exponential generating
function
\begin{equation}
\label{eq:r-stirling-egf}
 \sum_{n\ge r}\sum_{k\ge r}\stirtwo nk_{\!r}
 y^k\frac{z^{n-r}}{(n-r)!}
 =y^re^{rz}\exp\!\bigl(y(e^z-1)\bigr).
\end{equation}
Indeed, the $r$ distinguished blocks each consist of one distinguished atom and an
arbitrary set of ordinary atoms, contributing $y^re^{rz}$, while the remaining
blocks form an ordinary set partition.

\section{Associated Stirling numbers}
Let $\mathsf S_r(n,k)$ count partitions of $[n]$ into $k$ blocks, every block having
size at least $r$.

\begin{theorem}[Associated recurrence]
\label{thm:associated-stirling-recurrence}
For $r\ge1$,
\begin{equation}
\label{eq:associated-stirling-recurrence}
 \mathsf S_r(n+1,k)
 =k\mathsf S_r(n,k)+\binom n{r-1}\mathsf S_r(n-r+1,k-1).
\end{equation}
Moreover,
\begin{equation}
\label{eq:associated-stirling-egf}
 \sum_{n\ge0}\mathsf S_r(n,k)\frac{z^n}{n!}
 =\frac1{k!}\left(e^z-\sum_{j=0}^{r-1}\frac{z^j}{j!}\right)^k.
\end{equation}
\end{theorem}
\begin{proof}
Consider the block containing $n+1$.  If deleting $n+1$ leaves a block of size at
least $r$, start from a partition counted by $\mathsf S_r(n,k)$ and choose one of
its $k$ blocks for insertion.  Otherwise the block containing $n+1$ has exactly
$r$ elements: choose its $r-1$ companions and partition the remaining $n-r+1$
labels into $k-1$ admissible blocks.  These cases are disjoint and exhaustive,
proving the recurrence.  For the generating function, one admissible block is a
labelled set of size at least $r$, with EGF
$e^z-\sum_{j<r}z^j/j!$; an unordered set of $k$ such blocks contributes its $k$th
power divided by $k!$.
\end{proof}

For $r=2$ these are often called Ward numbers.  Their occurrence as coefficients
of other polynomial families comes from the same exclusion of singleton
components expressed by \eqref{eq:associated-stirling-egf}.

\section{Reduced Stirling numbers and graph coloring}
For $d\ge1$, let $\mathsf S^{[d]}(n,k)$ count partitions of $[n]$ into $k$ blocks
such that two labels in the same block always differ by at least $d$.

\begin{theorem}[Reduction formula]
\label{thm:reduced-stirling}
For $n\ge k\ge d$,
\begin{equation}
\label{eq:reduced-stirling}
 \mathsf S^{[d]}(n,k)=\stirtwo{n-d+1}{k-d+1}.
\end{equation}
In particular $\mathsf S^{[1]}(n,k)=\stirtwo nk$.
\end{theorem}
\begin{proof}
Join two vertices $i,j\in[n]$ when $0<|i-j|<d$, and call the resulting graph
$G_{n,d}$.  Its independent-set partitions are exactly the partitions counted by
$\mathsf S^{[d]}(n,k)$.  A proper coloring can therefore be obtained by choosing
such a partition and assigning distinct colors to its blocks, so
\begin{equation}
\label{eq:reduced-chromatic-falling}
 \chi_{G_{n,d}}(q)=\sum_k\mathsf S^{[d]}(n,k)\falling qk.
\end{equation}
Color the vertices in their natural order.  The first $d-1$ vertices must receive
distinct colors; every subsequent vertex sees the preceding $d-1$ vertices, which
form a clique, and hence has $q-d+1$ choices.  Thus
\[
 \chi_{G_{n,d}}(q)
 =\falling q{d-1}(q-d+1)^{n-d+1}.
\]
Apply the ordinary Stirling basis expansion to the last power:
\begin{align*}
 \chi_{G_{n,d}}(q)
 &=\falling q{d-1}
   \sum_{r\ge0}\stirtwo{n-d+1}{r}\falling{q-d+1}{r}\\
 &=\sum_{r\ge0}\stirtwo{n-d+1}{r}\falling q{r+d-1}.
\end{align*}
Comparison with \eqref{eq:reduced-chromatic-falling} gives
$k=r+d-1$ and proves \eqref{eq:reduced-stirling}.
\end{proof}
\part{Bell Numbers and Set Partitions}

\chapter{Set partitions and the many meanings of the Bell numbers}
\index{Bell numbers}

\section{Definition, initial values, and equivalence relations}
\begin{definition}
For $n\ge0$, the \emph{$n$th Bell number} $\Bell n$ is the number of partitions of
an $n$-element labelled set into nonempty, pairwise disjoint blocks.  Thus
\[
 \Bell0=\Bell1=1,
 \qquad
 (\Bell n)_{n\ge0}=1,1,2,5,15,52,203,877,4140,\ldots .
\]
The empty set has one partition, namely the empty family of blocks.
\end{definition}

For example, the five partitions of $\{a,b,c\}$ are
\[
 abc,\qquad a\mid bc,\qquad b\mid ac,\qquad c\mid ab,\qquad a\mid b\mid c,
\]
where vertical bars separate blocks.

\begin{proposition}
Partitions of a set $S$ are in natural bijection with equivalence relations on $S$.
Consequently $\Bell n$ also counts equivalence relations on an $n$-element set.
\end{proposition}
\begin{proof}
Given a partition $\pi$, declare $x\sim_\pi y$ when $x$ and $y$ lie in the same
block.  Reflexivity, symmetry, and transitivity follow immediately.  Conversely,
the equivalence classes of an equivalence relation are nonempty, disjoint, cover
$S$, and hence form a partition.  The two constructions undo one another.
\end{proof}


\section{Ordered set partitions and the Fubini numbers}
If the blocks themselves are linearly ordered, the resulting numbers are usually
called the \emph{ordered Bell numbers} or \emph{Fubini numbers}.  To prevent a
collision with the Bell notation, write them as $\mathsf F_n$.

\begin{theorem}[Ordered Bell numbers]
\label{thm:ordered-bell}
For $n\ge0$,
\begin{align}
 \mathsf F_n&=\sum_{k=0}^{n}k!\stirtwo nk,
 \label{eq:ordered-bell-stirling}\\
 \sum_{n\ge0}\mathsf F_n\frac{z^n}{n!}
 &=\frac{1}{2-e^z},
 \label{eq:ordered-bell-egf}\\
 \mathsf F_n&=\sum_{j=1}^{n}\binom nj\mathsf F_{n-j}
 \qquad(n\ge1),
 \label{eq:ordered-bell-recurrence}
\end{align}
with $\mathsf F_0=1$.
\end{theorem}
\begin{proof}
A partition into $k$ blocks may be ordered in $k!$ ways, proving
\eqref{eq:ordered-bell-stirling}.  By the second-kind exponential generating
function,
\[
 \sum_{n\ge0}\mathsf F_n\frac{z^n}{n!}
 =\sum_{k\ge0}k!\frac{(e^z-1)^k}{k!}
 =\sum_{k\ge0}(e^z-1)^k
 =\frac1{2-e^z},
\]
first as a formal identity and then analytically near the origin.  Finally, in a
nonempty ordered partition choose the elements of its first block, which may be any
nonempty subset, and order-partition the complement.  Summing over the first-block
size gives \eqref{eq:ordered-bell-recurrence}.
\end{proof}

Thus
\[
 (\mathsf F_n)_{n\ge0}=1,1,3,13,75,541,4683,\ldots,
\]
which contrasts with the unordered Bell sequence by recording the order of the
blocks rather than only their underlying family.
\section{Squarefree factorizations, rhyme schemes, and restricted-growth words}
\begin{proposition}
If $N=p_1\cdots p_n$ is squarefree, then $\Bell n$ is the number of unordered
factorizations of $N$ into integers greater than $1$.
\end{proposition}
\begin{proof}
To a set partition $\pi$ of $[n]$ associate the multiset of factors
$\{\prod_{i\in B}p_i:B\in\pi\}$.  Unique factorization recovers each block from
its factor, so this is a bijection.  For $30=2\cdot3\cdot5$ the five
factorizations are
\[
 30,\quad 2\cdot15,\quad3\cdot10,\quad5\cdot6,\quad2\cdot3\cdot5.
\]
\end{proof}

A rhyme scheme of an $n$-line stanza is likewise a partition of the set of lines:
two indices lie in the same block exactly when the corresponding lines rhyme.
The customary lettering convention is the following canonical encoding.

\begin{definition}
A \emph{restricted-growth word} of length $n$ is a word
$a_1\cdots a_n$ of nonnegative integers such that $a_1=0$ and
\[
 a_i\le 1+\max(a_1,\ldots,a_{i-1})\qquad(i>1).
\]
\end{definition}

\begin{proposition}
Set partitions of $[n]$, restricted-growth words of length $n$, and canonical
rhyme schemes of length $n$ are mutually bijective.
\end{proposition}
\begin{proof}
Order the blocks by increasing least element, number them $0,1,\ldots$, and put
$a_i=j$ when $i$ belongs to block $j$.  The displayed restriction says precisely
that a new block receives the next unused number.  Conversely, equal letters form
the blocks.  Replacing $0,1,2,\ldots$ by $A,B,C,\ldots$ gives the rhyme scheme.
\end{proof}

\section{A card-shuffling bijection}
Consider the operation that removes the top card from a deck of $n$ labelled cards
and reinserts it in any of the $n$ positions.  There are $n^n$ sequences of $n$
such operations.

\begin{theorem}[Return-to-order shuffles]
Exactly $\Bell n$ of those $n^n$ operation sequences return an initially ordered
deck to its original order.  Hence the return probability is $\Bell n/n^n$.
\end{theorem}
\begin{proof}
Reverse time.  A reversed operation chooses an arbitrary card and moves it to the
top.  Thus a reversed history is a word $c_1\cdots c_n$ in the card labels.  After
all moves, the cards that occurred in the word appear first, in decreasing order of
their last occurrence times; cards never chosen follow in their original relative
order.

Partition the time set $[n]$ according to equality of the letters $c_i$.  For the
final deck to be $1,2,\ldots,n$, the labels that occur must be an initial interval
$1,\ldots,m$, and their blocks of occurrence times must be labelled in decreasing
order of their maxima: the block with largest maximum receives label $1$, the next
receives $2$, and so on.  Therefore every set partition of $[n]$ determines exactly
one successful word, and every successful word determines its equality partition.
This is the required bijection.
\end{proof}

\section{Bell-enumerated vincular pattern classes}
A dash in a permutation pattern means that the corresponding entries need not be
adjacent; an undashed pair must be adjacent.  For instance, an occurrence of
$23\!\mathord{-}\!1$ consists of an adjacent ascent followed later by a smaller
entry.

\begin{theorem}
The permutations of $[n]$ avoiding $23\!\mathord{-}\!1$ are counted by $\Bell n$.
The same is true for each pattern in
\[
 1\!\mathord{-}\!23,\ 12\!\mathord{-}\!3,\ 32\!\mathord{-}\!1,
 \ 3\!\mathord{-}\!21,\ 1\!\mathord{-}\!32,\ 3\!\mathord{-}\!12,
 \ 21\!\mathord{-}\!3,
\]
and for the equivalent barred-pattern class in which every $321$ occurrence
extends to a $3241$ occurrence.
\end{theorem}
\begin{proof}
For the pattern $23\!\mathord{-}\!1$, order the blocks of a partition by increasing
least element and write the entries of each block in decreasing order.  Concatenate
the resulting words.  In the resulting permutation, each block is a decreasing
run, and every ascent occurs between two blocks.  Its lower endpoint is a
right-to-left minimum, so no later entry is smaller; hence $23\!\mathord{-}\!1$ is
avoided.

Conversely, cut a $23\!\mathord{-}\!1$-avoiding permutation immediately after
each right-to-left minimum.  Every resulting segment is decreasing: otherwise an
ascent inside a segment, together with its terminal right-to-left minimum, would
form $23\!\mathord{-}\!1$.  The segment minima increase from left to right, so the
segments, viewed as sets, are blocks already in canonical order.  The two maps are
inverse.

Reversal, complement, inverse, and their compositions are bijections on
permutations and carry vincular occurrences to occurrences of the seven listed
patterns; therefore all eight classes are equinumerous.  Finally, cutting at
right-to-left minima shows that failure of the decreasing-run condition is
simultaneously an occurrence of $23\!\mathord{-}\!1$ and a $321$ that cannot be
extended by an intervening entry to $3241$.  Thus the barred-pattern formulation
describes the same class.
\end{proof}

\begin{remark}
Classical avoidance classes have at most exponential growth in $n$, whereas Bell
numbers grow faster than $C^n$ for every fixed $C$.  The adjacency requirement is
therefore essential: these Bell-enumerated classes cannot be classical
single-pattern avoidance classes.
\end{remark}

\section{The Bell--Aitken triangle}
Define a triangular array $(a_{n,k})_{0\le k\le n}$ by
\[
 a_{0,0}=1,\qquad a_{n,0}=a_{n-1,n-1},\qquad
 a_{n,k}=a_{n,k-1}+a_{n-1,k-1}\quad(1\le k\le n).
\]
Its first rows are
\[
\begin{array}{rrrrr}
1\\
1&2\\
2&3&5\\
5&7&10&15\\
15&20&27&37&52.
\end{array}
\]

\begin{theorem}
For $0\le k\le n$,
\begin{equation}\label{eq:bell-triangle-closed}
 a_{n,k}=\sum_{j=0}^{k}\binom{k}{j}\Bell{n-j}.
\end{equation}
Consequently the left edge is $a_{n,0}=\Bell n$ and the right edge is
$a_{n,n}=\Bell{n+1}$.
\end{theorem}
\begin{proof}
The formula is immediate for $k=0$.  Pascal's identity gives
\[
 a_{n,k-1}+a_{n-1,k-1}
 =\sum_j\left(\binom{k-1}{j}+\binom{k-1}{j-1}\right)\Bell{n-j},
\]
which is \eqref{eq:bell-triangle-closed}.  For $k=n$, reverse the index and use
the Bell recurrence proved in \cref{thm:bell-binomial-recurrence} below.
\end{proof}

\chapter{The enumerative calculus of Bell numbers}

\section{Two fundamental sums}
\begin{theorem}[Binomial recurrence]\label{thm:bell-binomial-recurrence}
For $n\ge0$,
\begin{equation}\label{eq:bell-binomial-recurrence}
 \Bell{n+1}=\sum_{k=0}^{n}\binom nk\Bell k.
\end{equation}
\end{theorem}
\begin{proof}
In a partition of $\{0,1,\ldots,n\}$, let $k$ be the number of elements outside
the block containing $0$.  Choose those $k$ elements and partition them
arbitrarily; the block containing $0$ is then forced.
\end{proof}

\begin{theorem}[Stirling decomposition]
\begin{equation}\label{eq:bell-stirling-sum}
 \Bell n=\sum_{k=0}^{n}\stirtwo nk
 =\sum_{k=0}^{n}\frac1{k!}\sum_{i=0}^{k}(-1)^{k-i}\binom ki i^n.
\end{equation}
\end{theorem}
\begin{proof}
The first equality groups partitions by their number of blocks.  Substitute the
inclusion--exclusion formula for $\stirtwo nk$ from
\cref{eq:second-explicit} to obtain the second equality.
\end{proof}

\section{Spivey's identity and binomial inversions}
\begin{theorem}[Spivey's identity]\label{thm:spivey}
For $m,n\ge0$,
\begin{equation}\label{eq:spivey}
 \Bell{m+n}=\sum_{k=0}^{n}\sum_{j=0}^{m}
 \stirtwo mj\binom nk j^{\,n-k}\Bell k.
\end{equation}
\end{theorem}
\begin{proof}
Partition the first $m$ labels into $j$ blocks.  Of the final $n$ labels, choose
$k$ whose blocks contain no old label and partition them in $\Bell k$ ways.  Each
of the remaining $n-k$ labels independently joins one of the $j$ old blocks.  This
classification is unique and gives the summand.
\end{proof}

Binomial inversion of \eqref{eq:bell-binomial-recurrence} gives the first identity
below.  The second is a useful shifted form.
\begin{theorem}[Bell inversion identities]\label{thm:bell-inversions}
For $n,k\ge0$,
\begin{align}
 \Bell n&=\sum_{j=0}^{n}(-1)^{n-j}\binom nj\Bell{j+1},
 \label{eq:bell-inversion-one}\\
 \sum_{j=0}^{n}\binom nj\Bell{k+j}
 &=\sum_{i=0}^{k}(-1)^{k-i}\binom ki\Bell{n+i+1}.
 \label{eq:bell-inversion-two}
\end{align}
\end{theorem}
\begin{proof}
The first formula is ordinary binomial inversion.  For the second let $X$ be a
Poisson random variable of mean $1$; \cref{thm:dobinski} below gives
$\E X^r=\Bell r$.  The left side is
$\E[X^k(X+1)^n]$.  Poisson summation by parts,
\begin{equation}\label{eq:poisson-sbp}
 \E[Xh(X)]=\E[h(X+1)],
\end{equation}
iterated $k$ times in the equivalent form
$\E[X^{n+1}(X-1)^k]=\E[(X+1)^nX^k]$, turns it into the right side after expanding
$(X-1)^k$.
\end{proof}

The first-kind Stirling numbers package all repeated shifts at once.
\begin{theorem}[Weighted Bell shift]\label{thm:weighted-bell-shift}
For $a,b\in\C$ and $k,n\ge0$,
\begin{multline}\label{eq:weighted-bell-shift}
 \sum_{j=0}^{n}\binom nj a^j b^{n-j}\Bell j\\
 =\sum_{i=0}^{k}(-1)^{k-i}\stirone ki
   \sum_{j=0}^{n}\binom nj a^j(b-ak)^{n-j}\Bell{j+i}.
\end{multline}
In particular,
\begin{align}
 \sum_{j=0}^{n}\binom nj a^j b^{n-j}\Bell j
 &=\sum_{j=0}^{n}\binom nj a^j(b-a)^{n-j}\Bell{j+1},
 \label{eq:weighted-bell-k1}\\
 \sum_{j=0}^{n}\binom nj k^{n-j}\Bell j
 &=\sum_{i=0}^{k}(-1)^{k-i}\stirone ki\Bell{n+i}.
 \label{eq:weighted-bell-special}
\end{align}
\end{theorem}
\begin{proof}
By Dobiński's formula, the left side of
\eqref{eq:weighted-bell-shift} is
\[
 e^{-1}\sum_{m\ge0}\frac{(am+b)^n}{m!}.
\]
On the right, use
$\sum_i(-1)^{k-i}\stirone ki m^i=\falling mk$ inside the same Poisson sum.  It
becomes
\[
 e^{-1}\sum_{m\ge0}\frac{\falling mk\,[a m+b-ak]^n}{m!}
 =e^{-1}\sum_{r\ge0}\frac{(ar+b)^n}{r!},
\]
after $m=r+k$.  The special cases follow by setting $k=1$, and then $a=1,b=k$.
\end{proof}

\section{Exponential generating function and differential equation}
\begin{theorem}\label{thm:bell-egf}
The Bell exponential generating function is
\begin{equation}\label{eq:bell-egf}
 \mathscr B(z):=\sum_{n\ge0}\Bell n\frac{z^n}{n!}
 =\exp(e^z-1),
 \qquad
 \mathscr B'(z)=e^z\mathscr B(z).
\end{equation}
\end{theorem}
\begin{proof}
A set partition is a labelled set of nonempty labelled sets.  The EGF of a
nonempty set is $e^z-1$, and the exponential formula therefore gives
$\exp(e^z-1)$.  Differentiation gives the differential equation; conversely,
coefficient comparison in that equation gives
\eqref{eq:bell-binomial-recurrence} and $\mathscr B(0)=1$, so it uniquely determines
the series.
\end{proof}

\section{Dobiński's formula and Poisson moments}
\begin{theorem}[Dobiński]\label{thm:dobinski}
For every $n\ge0$,
\begin{equation}\label{eq:dobinski}
 \Bell n=e^{-1}\sum_{m=0}^{\infty}\frac{m^n}{m!}.
\end{equation}
Equivalently, if $X\sim\operatorname{Poisson}(1)$, then $\Bell n=\E X^n$.
\end{theorem}
\begin{proof}
Expand the outer exponential in \eqref{eq:bell-egf}:
\[
 e^{e^z-1}=e^{-1}\sum_{m\ge0}\frac{e^{mz}}{m!}
 =e^{-1}\sum_{m\ge0}\frac1{m!}\sum_{n\ge0}m^n\frac{z^n}{n!}.
\]
Absolute convergence on compact sets justifies rearrangement analytically, while
formal coefficient extraction gives the same identity algebraically.  The last
statement is the defining moment sum of a Poisson$(1)$ variable.
\end{proof}

\begin{corollary}[Cauchy integral]
For every positively oriented circle $\gamma$ around $0$,
\begin{equation}\label{eq:bell-cauchy}
 \Bell n=\frac{n!}{2\pi i e}\int_\gamma
       \frac{e^{e^z}}{z^{n+1}}\,\dd z.
\end{equation}
\end{corollary}
\begin{proof}
Apply Cauchy's coefficient formula to $e^{-1}e^{e^z}$.
\end{proof}

\chapter{Congruences, shape, and growth of Bell numbers}

\section{Touchard congruences}
It is useful to refine $\Bell n$ to the Touchard polynomial
\[
 T_n(x)=\sum_{k=0}^{n}\stirtwo nk x^k,
 \qquad T_n(1)=\Bell n.
\]

\begin{theorem}[Polynomial Touchard congruence]\label{thm:touchard-poly}
For a prime $p$,
\begin{equation}\label{eq:touchard-poly}
 T_{n+p}(x)\equiv T_{n+1}(x)+x^pT_n(x)\pmod p.
\end{equation}
Consequently
\begin{equation}\label{eq:touchard}
 \Bell{n+p}\equiv\Bell n+\Bell{n+1}\pmod p.
\end{equation}
\end{theorem}
\begin{proof}
Let the cyclic group $C_p$ rotate a distinguished set of $p$ labels while fixing
$n$ other labels, and let it act on partitions, preserving the weight $x^{\#\text{
blocks}}$.  Every nonfixed orbit has size $p$ and contributes $0$ modulo $p$.
In a fixed partition there are two possibilities.  Either all $p$ rotating labels
form $p$ singleton blocks, independently of a partition of the other labels; this
contributes $x^pT_n(x)$.  Otherwise rotational invariance forces the rotating labels
to lie in one common block, possibly together with old labels.  Contracting those
$p$ labels to one distinguished label gives a partition of $n+1$ labels with the
same number of blocks, contributing $T_{n+1}(x)$.  This proves the polynomial
congruence; set $x=1$.
\end{proof}

\begin{theorem}[Prime-power shift]\label{thm:bell-prime-power-shift}
For a prime $p$ and $m\ge1$,
\begin{equation}\label{eq:bell-prime-power-shift}
 \Bell{n+p^m}\equiv m\Bell n+\Bell{n+1}\pmod p.
\end{equation}
More precisely,
\begin{equation}\label{eq:touchard-prime-power-poly}
 T_{n+p^m}(x)\equiv T_{n+1}(x)+T_n(x)\sum_{r=1}^{m}x^{p^r}
 \pmod p.
\end{equation}
\end{theorem}
\begin{proof}
Let the additive cyclic group of order $p^m$ act regularly on $p^m$ new labels.
A partition fixed by the whole group induces a block system for this regular
action.  Apart from the case in which all new labels join one invariant block
(which contracts to one label and contributes $T_{n+1}$), the blocks on the new
orbit are the cosets of the unique subgroup of index $p^r$ for some
$1\le r\le m$; they contribute $p^r$ blocks and cannot mix with old labels.
Thus the fixed weighted partitions contribute
$T_{n+1}+T_n\sum_{r=1}^m x^{p^r}$.  Every other orbit has cardinality divisible by
$p$, proving \eqref{eq:touchard-prime-power-poly}.  Fermat's congruence
$x^{p^r}\equiv x$ modulo $p$ at $x=1$ gives
\eqref{eq:bell-prime-power-shift}.
\end{proof}

\section{A universal period bound modulo a prime}
\begin{theorem}\label{thm:bell-period-bound}
Modulo a prime $p$, the Bell sequence is periodic, and its period divides
\begin{equation}\label{eq:bell-period-N}
 N_p=\frac{p^p-1}{p-1}=1+p+\cdots+p^{p-1}.
\end{equation}
\end{theorem}
\begin{proof}
Touchard's recurrence says that the shift operator $E$ on the sequence satisfies
$E^p-E-1=0$ over $\F_p$.  The polynomial $q(t)=t^p-t-1$ is separable because
$q'(t)=-1$.  If $\alpha$ is any root in a splitting field, then
$\alpha^p=\alpha+1$ and hence $\alpha^{p^j}=\alpha+j$ for $j\in\F_p$.  Taking the
product over $j$ gives
\[
 \alpha^{N_p}=\prod_{j\in\F_p}\alpha^{p^j}
 =\prod_{j\in\F_p}(\alpha+j)=\alpha^p-\alpha=1.
\]
Thus every eigenvalue of $E$ is an $N_p$th root of unity.  Separability makes $E$
diagonalizable over the splitting field, so $E^{N_p}=1$ on every solution of the
recurrence, in particular on the Bell sequence.
\end{proof}

\begin{remark}
A reported exact period is computational data, not a consequence of the divisibility
bound alone.  It can be certified by generating $N_p$ terms with
\eqref{eq:touchard}, finding the least divisor $d\mid N_p$ for which the state vector
of $p$ consecutive terms repeats, and checking that no proper divisor of $d$ works.
The same recurrence permits a compact machine-checkable certificate.
\end{remark}

\section{Log-convexity and normalized log-concavity}
\begin{theorem}\label{thm:bell-log-shape}
For $n\ge1$,
\begin{equation}\label{eq:bell-log-shape}
 1\le \frac{\Bell{n-1}\Bell{n+1}}{\Bell n^2}\le\frac{n+1}{n}.
\end{equation}
Thus $(\Bell n)$ is log-convex, whereas $(\Bell n/n!)$ is log-concave.
\end{theorem}
\begin{proof}
For the left inequality, let $X\sim\operatorname{Poisson}(1)$.  Cauchy--Schwarz
applied to $X^{(n-1)/2}$ and $X^{(n+1)/2}$ gives
$(\E X^n)^2\le \E X^{n-1}\E X^{n+1}$.

For completeness we prove the normalized inequality through the following
convolution lemma.  If nonnegative $(z_j)_{j\ge1}$ is log-concave without internal
zeros and
\[
 A(t)=\exp\!\left(\sum_{j\ge1}\frac{z_j}{j}t^j\right)
      =\sum_{n\ge0}a_n\frac{t^n}{n!},
\]
then $a_n/n!$ is log-concave.  To see this, truncate at degree $N$, differentiate,
and compare the adjacent $2\times2$ Toeplitz minors in
$A'=A\sum_{j\ge1}z_jt^{j-1}$.  After cross multiplication, the $n$th minor is
\[
 \sum_{0\le i<j\le n}
 \frac{a_i a_j}{i!j!}
 \bigl(z_{n-i+1}z_{n-j}-z_{n-i}z_{n-j+1}\bigr),
\]
with the boundary terms interpreted as zero.  Each bracket is nonnegative by the
monotonicity of the ratios $z_{r+1}/z_r$; induction on $n$ starts at $a_0=1$.
Removing the truncation proves the lemma coefficientwise.

For $A(t)=e^{e^t-1}$ we have
$\sum_{j\ge1}t^j/j!=\sum_{j\ge1}z_jt^j/j$ with
$z_j=1/(j-1)!$.  This sequence is log-concave, since
$z_j^2/(z_{j-1}z_{j+1})=j/(j-1)$.  Hence
$(\Bell n/n!)^2\ge(\Bell{n-1}/(n-1)!)(\Bell{n+1}/(n+1)!)$, which is the right
inequality in \eqref{eq:bell-log-shape}.
\end{proof}

\section{Explicit upper bounds}
The positive-coefficient Cauchy estimate supplies a short route to the numerical
bounds in the archive.
\begin{lemma}\label{lem:bell-coeff-upper}
For $t>0$,
\begin{equation}\label{eq:bell-coeff-upper}
 \Bell n\le n!\,t^{-n}\exp(e^t-1).
\end{equation}
For $n\ge2$, put $L=\log(n+1)$ and $t=\log(n/L)$.  Then
\begin{equation}\label{eq:bell-ratio-upper}
 \Bell n^{1/n}\frac{L}{n}
 \le \frac{L}{t}
 \exp\!\left(-1+\frac1L+\frac{\log n}{2n}\right).
\end{equation}
\end{lemma}
\begin{proof}
Because all coefficients of $\mathscr B$ are nonnegative,
$\Bell n t^n/n!\le\mathscr B(t)$, proving the first claim.  Use
$e^t=n/L$ and the elementary Stirling bound
$n!\le e\sqrt n\,(n/e)^n$; the factors $e^{-1/n}$ and $e^{1/n}$ cancel and give
\eqref{eq:bell-ratio-upper}.
\end{proof}

\begin{theorem}[Uniform numerical bounds]\label{thm:bell-numerical-bounds}
For every $n\ge1$,
\begin{equation}\label{eq:bell-0792}
 \Bell n<\left(\frac{0.792\,n}{\log(n+1)}\right)^n.
\end{equation}
Define, for $x>1$,
\begin{equation}\label{eq:bell-dx}
 d(x)=\log\log(x+1)-\log\log x+\frac{1+e^{-1}}{\log x}.
\end{equation}
For every $\varepsilon>0$ there is
$n_0(\varepsilon)=\max\{e^4,d^{-1}(\varepsilon)\}$ such that, for
$n>n_0(\varepsilon)$,
\begin{equation}\label{eq:bell-e06}
 \Bell n<\left(\frac{e^{-0.6+\varepsilon}n}{\log(n+1)}\right)^n.
\end{equation}
\end{theorem}
\begin{proof}
Set
\[
 R(x)=\frac{\log(x+1)}{\log(x/\log(x+1))}
 \exp\!\left(-1+\frac1{\log(x+1)}+\frac{\log x}{2x}\right).
\]
Direct differentiation, after multiplying by the positive denominator
$2x\log x\log(x+1)\log(x/\log(x+1))^2$, shows that $R'(x)<0$ for $x\ge39$;
substitution gives $R(39)<0.792$.  Thus
\eqref{eq:bell-ratio-upper} proves \eqref{eq:bell-0792} for $n\ge39$.  The
remaining $1\le n\le38$ are obtained successively from
\eqref{eq:bell-binomial-recurrence}; substituting them gives strict inequalities.
This is a finite exact integer check.

A second differentiation gives, for $x\ge e^4$,
\[
 \log R(x)\le-0.6+d(x),\qquad d'(x)<0,\qquad \lim_{x\to\infty}d(x)=0.
\]
(The numerator after clearing positive denominators is a sum of
$-(\log x-4)^2$, $-(x-1-\log x)$, and negative multiples of these two standard
nonnegative quantities.)  Hence $d(n)<\varepsilon$ when
$n>d^{-1}(\varepsilon)$, and \eqref{eq:bell-e06} follows from
\eqref{eq:bell-ratio-upper}.
\end{proof}

\section{Saddle-point asymptotics and Lambert \texorpdfstring{$W$}{W}}
Let $r=W(n)$, so that $re^r=n$, and put $A=n(1+r)$.
\begin{theorem}[Leading Bell asymptotic]\label{thm:bell-leading-asymptotic}
As $n\to\infty$,
\begin{align}
 \Bell n
 &\sim \frac{n!\exp(e^r-1)}{r^n\sqrt{2\pi n(1+r)}}
 \label{eq:bell-saddle-leading}\\
 &\sim \frac1{\sqrt n}\left(\frac{n}{W(n)}\right)^{n+1/2}
       \exp\!\left(\frac{n}{W(n)}-n-1\right).
 \label{eq:bell-W-asymptotic}
\end{align}
\end{theorem}
\begin{proof}
In \eqref{eq:bell-cauchy} take the circle $z=re^{i\theta}$.  The coefficient
integral becomes
\[
 \frac{\Bell n}{n!}=
 \frac{e^{e^r-1}}{2\pi r^n}
 \int_{-\pi}^{\pi}
 \exp\!\left(e^{re^{i\theta}}-e^r-in\theta\right)\dd\theta.
\]
The linear term vanishes because $re^r=n$; the quadratic term is
$-A\theta^2/2$.  On $|\theta|\le A^{-2/5}$, Taylor expansion is uniform and the
scaled integral tends to $\int_{\R}e^{-u^2/2}\dd u/\sqrt A$.  Outside that arc,
$\Re(e^{re^{i\theta}}-e^r)\le-cA\theta^2$ on the intermediate range and is
$\le-c'e^r$ on the remaining compact range, so the contribution is negligible.
This proves \eqref{eq:bell-saddle-leading}.  Stirling's formula and $e^r=n/r$
simplify it to \eqref{eq:bell-W-asymptotic}.
\end{proof}

The archived refined expansion used undefined symbols.  The following construction
makes every coefficient unambiguous.  Let
$T_j(r)=\sum_{q=0}^j\stirtwo jq r^q$ be the Touchard polynomial and let
$Z\sim N(0,1)$.  For an integer $h=O(r)$ define
\begin{equation}\label{eq:bell-saddle-correction-generator}
 \mathcal R_{r,h}(Z)=\exp\!\left(
 -\frac{ihZ}{\sqrt A}+
 \sum_{j\ge3}\frac{i^j e^rT_j(r)}{j!A^{j/2}}Z^j
 \right).
\end{equation}
Expand this formal exponential and take Gaussian moments termwise; retaining all
monomials of total order at most $e^{-Mr}$ defines $C_0(r,h),\ldots,C_M(r,h)$,
with $C_0=1$.

\begin{theorem}[Uniform Moser--Wyman scheme]\label{thm:bell-refined-saddle}
For every fixed $M$ and $|h|\le Cr$,
\begin{equation}\label{eq:bell-refined-saddle}
 \Bell{n+h}=\frac{(n+h)!\exp(e^r-1)}{r^{n+h}\sqrt{2\pi A}}
 \left(\sum_{q=0}^{M-1}C_q(r,h)+O_C(e^{-Mr}P_M(r))\right),
\end{equation}
where $P_M$ is an explicitly computable polynomial.  The first correction is
\begin{equation}\label{eq:bell-first-correction}
 C_1(r,h)=
 -\frac{h^2}{2A}-\frac{h e^rT_3(r)}{2A^2}
 +\frac{e^rT_4(r)}{8A^2}
 -\frac{5e^{2r}T_3(r)^2}{24A^3}.
\end{equation}
In particular the coefficient of $e^{-qr}$ is a polynomial in $h$ of degree at
most $2q$, with rational functions of $r$ as coefficients.
\end{theorem}
\begin{proof}
Use the same circle as in the preceding proof but replace $n$ by $n+h$.  After
$\theta=Z/\sqrt A$, the ratio of the integrand to its Gaussian quadratic part is
exactly the formal series \eqref{eq:bell-saddle-correction-generator}.  Watson's
lemma with a Gaussian majorant permits termwise integration on the central arc;
the complementary arcs have the same exponential bound uniformly for
$|h|\le Cr$.  Odd Gaussian moments vanish, so only integral powers of $e^{-r}$
remain.  Since $\E Z^2=1$, $\E Z^4=3$, and $\E Z^6=15$, the terms from the square
of the $h$-term, its product with the cubic term, the quartic term, and the square
of the cubic term give \eqref{eq:bell-first-correction}.  The same finite
bookkeeping defines every subsequent $C_q$ and bounds the omitted Taylor
remainder.
\end{proof}

\begin{corollary}[de Bruijn logarithmic expansion]\label{cor:bell-debruijn}
As $n\to\infty$,
\begin{multline}\label{eq:bell-debruijn}
 \frac{\log\Bell n}{n}
 =\log n-\log\log n-1
 +\frac{\log\log n}{\log n}+\frac1{\log n}\\
 +\frac12\left(\frac{\log\log n}{\log n}\right)^2
 +O\!\left(\frac{\log\log n}{(\log n)^2}\right).
\end{multline}
\end{corollary}
\begin{proof}
Take logarithms in \eqref{eq:bell-W-asymptotic}.  The standard reversion of
$r+\log r=\log n$ gives
\[
 r=\log n-\log\log n+
rac{\log\log n}{\log n}
 +O\!\left(\frac{(\log\log n)^2}{(\log n)^2}\right).
\]
Substitution and one further Taylor expansion of $\log r$ produce the displayed
terms; the prefactor contributes only $O(\log n/n)$.
\end{proof}

\section{Bell primes and verifiable computation}
A \emph{Bell prime} is a prime value of $\Bell n$.  The archive records the early
indices found by computation.  Such a list is not a symbolic theorem, but each
entry has a rigorous certificate: compute $\Bell n$ by the integer recurrence
\eqref{eq:bell-binomial-recurrence}, apply a probable-prime filter, and attach an
ECPP or Pratt-style primality certificate.  Composite entries are certified by one
nontrivial factor.  This distinction between an identity and a finite certified
computation will also be used for tabulated values later.

\part{Bell Polynomials and Composition}

\chapter{Exponential and ordinary Bell polynomials}
\index{Bell polynomials}

\section{Definitions}
\begin{definition}[Partial exponential Bell polynomial]
For $0\le k\le n$, define
\begin{equation}\label{eq:partial-bell-definition}
 \BellP nk(x_1,\ldots,x_{n-k+1})
 =n!\!\sum_{\substack{j_1+j_2+\cdots=k\\
                       j_1+2j_2+3j_3+\cdots=n}}
 \prod_{i\ge1}\frac{x_i^{j_i}}{(i!)^{j_i}j_i!}.
\end{equation}
Only $i\le n-k+1$ can occur.  The complete exponential Bell polynomial is
\begin{equation}\label{eq:complete-bell-definition}
 \BellC n(x_1,\ldots,x_n)=\sum_{k=0}^{n}\BellP nk(x_1,\ldots,x_{n-k+1}),
 \qquad \BellC0=1.
\end{equation}
\end{definition}

\begin{definition}[Partial ordinary Bell polynomial]
\begin{equation}\label{eq:ordinary-bell-definition}
 \OBellP nk(x_1,\ldots,x_{n-k+1})
 =\sum_{\substack{j_1+j_2+\cdots=k\\
                       j_1+2j_2+\cdots=n}}
 \frac{k!}{j_1!j_2!\cdots}\prod_{i\ge1}x_i^{j_i}.
\end{equation}
Equivalently,
\begin{equation}\label{eq:ordinary-exponential-scaling}
 \OBellP nk(x_1,x_2,\ldots)
 =\frac{k!}{n!}\BellP nk(1!x_1,2!x_2,\ldots).
\end{equation}
\end{definition}
\begin{proof}[Proof of \eqref{eq:ordinary-exponential-scaling}]
Substitute $i!x_i$ for $x_i$ in \eqref{eq:partial-bell-definition}; the factorials
$(i!)^{j_i}$ cancel, and multiplication by $k!/n!$ leaves exactly
\eqref{eq:ordinary-bell-definition}.
\end{proof}

\section{Combinatorial interpretation}
Give each block of size $i$ a weight $x_i$, and define the weight of a partition as
the product of its block weights.

\begin{theorem}[Weighted-partition interpretation]\label{thm:bell-poly-partitions}
$\BellP nk$ is the total weight of partitions of an $n$-element labelled set into
exactly $k$ blocks, and $\BellC n$ is the total weight of all its partitions.
\end{theorem}
\begin{proof}
Fix multiplicities $j_i$, where $j_i$ blocks have size $i$.  The number of such
partitions is
\[
 \frac{n!}{\prod_i(i!)^{j_i}j_i!}.
\]
The factors $i!$ forget order inside each block and $j_i!$ forget order among
blocks of equal size.  Multiplication by $\prod_i x_i^{j_i}$ and summation gives
\eqref{eq:partial-bell-definition}; summing over $k$ gives the complete polynomial.
\end{proof}

Thus
\begin{align*}
 \BellP{3}{2}(x_1,x_2)&=3x_1x_2,\\
 \BellP{6}{2}(x_1,\ldots,x_5)&=6x_1x_5+15x_2x_4+10x_3^2,\\
 \BellP{6}{3}(x_1,\ldots,x_4)&=15x_1^2x_4+60x_1x_2x_3+15x_2^3.
\end{align*}
The monomials correspond to integer partitions of $n$ into $k$ parts, while their
coefficients count the labelled set partitions having the prescribed block sizes.

\section{The master generating functions}
Put
\[
 X(t)=\sum_{j\ge1}x_j\frac{t^j}{j!}.
\]
\begin{theorem}[Bell-polynomial generating functions]\label{thm:bell-poly-egf}
One has
\begin{align}
 \frac{X(t)^k}{k!}
 &=\sum_{n\ge k}\BellP nk(x_1,\ldots)\frac{t^n}{n!},\\
 \exp(uX(t))
 &=\sum_{n\ge k\ge0}\BellP nk(x_1,\ldots)u^k\frac{t^n}{n!},\\
 \exp(X(t))
 &=\sum_{n\ge0}\BellC n(x_1,\ldots,x_n)\frac{t^n}{n!}.
\end{align}
For $\widehat X(t)=\sum_{j\ge1}x_jt^j$,
\begin{align}
 \widehat X(t)^k&=\sum_{n\ge k}\OBellP nk(x_1,\ldots)t^n,
 \label{eq:ordinary-bell-ogf}\\
 e^{u\widehat X(t)}
 &=\sum_{n\ge k\ge0}\OBellP nk(x_1,\ldots)t^n\frac{u^k}{k!}.
 \label{eq:ordinary-bell-bivariate}
\end{align}
\end{theorem}
\begin{proof}
Expand $X(t)^k$ multinomially.  Choosing $j_i$ copies of the term
$x_it^i/i!$ gives the constraints $\sum j_i=k$ and $\sum ij_i=n$, and its
coefficient is $k!\prod_i x_i^{j_i}/((i!)^{j_i}j_i!)$.  Division by $k!$ gives
\eqref{eq:partial-bell-definition}.  Summing in $k$ proves the exponential
identities.  The ordinary identities follow by the same multinomial argument or
by \eqref{eq:ordinary-exponential-scaling}.
\end{proof}

\begin{masterbox}
Nearly every algebraic identity for Bell polynomials follows by performing a simple
operation---multiplication, differentiation, substitution, or logarithm---on
\eqref{eq:partial-bell-egf} or \eqref{eq:complete-bell-egf}.  The combinatorial
proofs are the same operations interpreted as selecting, pointing, colouring, or
gluing blocks.
\end{masterbox}

\section{Tables and low-degree polynomials}
The nonzero partial polynomials through $n=6$ are
\begin{align*}
\BellP11&=x_1,\\
\BellP21&=x_2,\qquad \BellP22=x_1^2,\\
\BellP31&=x_3,\qquad \BellP32=3x_1x_2,\qquad \BellP33=x_1^3,\\
\BellP41&=x_4,\qquad \BellP42=4x_1x_3+3x_2^2,\\
\BellP43&=6x_1^2x_2,\qquad \BellP44=x_1^4,\\
\BellP51&=x_5,\qquad \BellP52=5x_1x_4+10x_2x_3,\\
\BellP53&=10x_1^2x_3+15x_1x_2^2,\\
\BellP54&=10x_1^3x_2,\qquad \BellP55=x_1^5,\\
\BellP61&=x_6,\qquad \BellP62=6x_1x_5+15x_2x_4+10x_3^2,\\
\BellP63&=15x_1^2x_4+60x_1x_2x_3+15x_2^3,\\
\BellP64&=20x_1^3x_3+45x_1^2x_2^2,\\
\BellP65&=15x_1^4x_2,\qquad \BellP66=x_1^6.
\end{align*}
Every coefficient is supplied by \eqref{eq:partial-bell-definition}, so the table
is a computation from the theorem rather than an independent assertion.

The first complete polynomials are
\begin{align*}
 \BellC0={}&1,\\
 \BellC1={}&x_1,\\
 \BellC2={}&x_1^2+x_2,\\
 \BellC3={}&x_1^3+3x_1x_2+x_3,\\
 \BellC4={}&x_1^4+6x_1^2x_2+3x_2^2+4x_1x_3+x_4,\\
 \BellC5={}&x_1^5+10x_1^3x_2+15x_1x_2^2+10x_1^2x_3
 +10x_2x_3+5x_1x_4+x_5,\\
 \BellC6={}&x_1^6+15x_1^4x_2+45x_1^2x_2^2+15x_2^3
 +20x_1^3x_3+60x_1x_2x_3+10x_3^2\\
 &+15x_1^2x_4+15x_2x_4+6x_1x_5+x_6,\\
 \BellC7={}&x_1^7+21x_1^5x_2+105x_1^3x_2^2+105x_1x_2^3
 +35x_1^4x_3+210x_1^2x_2x_3\\
 &+70x_1x_3^2+105x_2^2x_3+35x_1^3x_4+105x_1x_2x_4
 +35x_3x_4\\
 &+21x_1^2x_5+21x_2x_5+7x_1x_6+x_7.
\end{align*}

\chapter{Recurrences, derivatives, and algebraic identities}

\section{Pointing recurrences}
\begin{theorem}[Complete and partial recurrences]\label{thm:bell-poly-recurrences}
For $n\ge0$,
\begin{align}
 \BellC{n+1}(x_1,\ldots,x_{n+1})
 &=\sum_{i=0}^{n}\binom ni x_{i+1}\BellC{n-i}(x_1,\ldots,x_{n-i}),
 \label{eq:complete-bell-recurrence}\\
 \BellP{n+1}{k+1}
 &=\sum_{i=0}^{n-k}\binom ni x_{i+1}\BellP{n-i}{k},
 \label{eq:partial-bell-recurrence}
\end{align}
with $\BellP00=1$, $\BellP n0=0$ for $n>0$, and $\BellP0k=0$ for $k>0$.
\end{theorem}
\begin{proof}
In a weighted partition of $[n+1]$, inspect the block containing the distinguished
label $n+1$.  If that block contains $i$ further labels, choose them in $\binom ni$
ways, give the block weight $x_{i+1}$, and partition the remaining labels.  Keeping
or forgetting the total block count yields the two formulas.
\end{proof}

\begin{theorem}[Block-colour convolution]\label{thm:bell-partial-convolution}
For $k_1,k_2\ge0$,
\begin{equation}\label{eq:bell-partial-convolution}
 \BellP n{k_1+k_2}(x)
 =\frac{k_1!k_2!}{(k_1+k_2)!}
 \sum_{i=0}^{n}\binom ni\BellP i{k_1}(x)\BellP{n-i}{k_2}(x).
\end{equation}
\end{theorem}
\begin{proof}
Multiply \eqref{eq:partial-bell-egf} for $k_1$ and $k_2$.  Their product is
$X^{k_1+k_2}/(k_1!k_2!)$; comparison with the same formula for $k_1+k_2$ gives the
factor $(k_1!k_2!)/(k_1+k_2)!$ and the binomial EGF convolution.
\end{proof}

\section{Near-diagonal reduction}
\begin{theorem}\label{thm:bell-near-diagonal}
For $1\le a<n$,
\begin{multline}\label{eq:bell-near-diagonal}
 \BellP n{n-a}(x_1,\ldots,x_{a+1})\\
 =\sum_{j=a+1}^{2a}\frac{j!}{a!}\binom nj x_1^{n-j}
 \BellP a{j-a}\!\left(\frac{x_2}{2},\frac{x_3}{3},\ldots,
 \frac{x_{2a-j+2}}{2a-j+2}\right).
\end{multline}
\end{theorem}
\begin{proof}
A partition into $n-a$ blocks has total \emph{excess}
$\sum_B(|B|-1)=a$.  Let $j$ be the number of elements in nonsingleton blocks.
Then the number of those blocks is $j-a$, so $a+1\le j\le2a$.  Choose the $j$
elements, leave the other $n-j$ as singleton blocks, and count the nonsingleton
partition.  Replacing a block of size $r+1$ by an atom of size $r$ removes one
distinguished element per block; the factor $j!/a!$ and weights $x_{r+1}/(r+1)$
are exactly the labelled pointing factors.  Equivalently, make the substitution in
\eqref{eq:partial-bell-definition}; the coefficient of every multiplicity vector
on both sides is
$n!x_1^{n-j}\prod_{r\ge2}x_r^{j_r}/((r!)^{j_r}j_r!)$.
\end{proof}

The first cases, obtained by listing integer partitions of the excess, are
\begin{align}
 \BellP n1&=x_n, &
 \BellP n2&=\frac12\sum_{j=1}^{n-1}\binom njx_jx_{n-j},
 \label{eq:bell-edge-first}\\
 \BellP nn&=x_1^n, &
 \BellP n{n-1}&=\binom n2x_1^{n-2}x_2,\\
 \BellP n{n-2}&=\binom n3x_1^{n-3}x_3
 +3\binom n4x_1^{n-4}x_2^2,\\
 \BellP n{n-3}&=\binom n4x_1^{n-4}x_4
 +10\binom n5x_1^{n-5}x_2x_3
 +15\binom n6x_1^{n-6}x_2^3,\\
 \BellP n{n-4}&=\binom n5x_1^{n-5}x_5
 +5\binom n6x_1^{n-6}(3x_2x_4+2x_3^2)\notag\\
 &\quad+105\binom n7x_1^{n-7}x_2^2x_3
 +105\binom n8x_1^{n-8}x_2^4.
 \label{eq:bell-near-diagonal-cases}
\end{align}
Their coefficients are the corresponding block-size counts from
\cref{thm:bell-poly-partitions}.

\section{Partial derivatives}
\begin{theorem}[Derivative identities]\label{thm:bell-poly-derivatives}
For $i\ge1$,
\begin{align}
 \frac{\partial\BellC n}{\partial x_i}
 &=\binom ni\BellC{n-i},
 \label{eq:complete-bell-derivative}\\
 \frac{\partial\BellP nk}{\partial x_i}
 &=\binom ni\BellP{n-i}{k-1}.
 \label{eq:partial-bell-derivative}
\end{align}
Consequently, for differentiable $a_i(x)$,
\begin{equation}\label{eq:bell-poly-chain}
 \frac{\dd}{\dd x}\BellP nk(a_1(x),\ldots)
 =\sum_{i=1}^{n-k+1}\binom ni a_i'(x)\BellP{n-i}{k-1}(a_1(x),\ldots).
\end{equation}
\end{theorem}
\begin{proof}
Differentiate \eqref{eq:partial-bell-egf} with respect to $x_i$:
\[
 \frac{\partial}{\partial x_i}\frac{X(t)^k}{k!}
 =\frac{t^i}{i!}\frac{X(t)^{k-1}}{(k-1)!}.
\]
Coefficient comparison gives \eqref{eq:partial-bell-derivative}; summing over $k$
gives \eqref{eq:complete-bell-derivative}.  The ordinary chain rule then gives
\eqref{eq:bell-poly-chain}.
\end{proof}

\section{A quadratic differential recurrence}
The less familiar differential recurrence in the archive is also a disguised
coefficient identity.
\begin{theorem}\label{thm:bell-quadratic-differential}
For $n\ge2$,
\begin{align}
 \BellC n=\frac1{n-1}\Bigg[&
 \sum_{i=2}^{n}\sum_{j=1}^{i-1}(i-1)\binom{i-2}{j-1}x_jx_{i-j}
       \frac{\partial\BellC{n-1}}{\partial x_{i-1}}\notag\\
 &+\sum_{i=2}^{n}\sum_{j=1}^{i-1}\frac{x_{i+1}}{\binom ij}
       \frac{\partial^2\BellC{n-1}}{\partial x_j\partial x_{i-j}}
 +\sum_{i=2}^{n}x_i\frac{\partial\BellC{n-1}}{\partial x_{i-1}}
 \Bigg].
 \label{eq:bell-quadratic-differential}
\end{align}
Terms containing unavailable variables or negative-index Bell polynomials are zero.
\end{theorem}
\begin{proof}
Let $F(t)=e^{X(t)}=\sum_m\BellC m t^m/m!$.  After using
\eqref{eq:complete-bell-derivative}, the first double sum is
\[
 (n-1)!\coeff{t}{n-1}\bigl(tX'(t)^2F(t)\bigr).
\]
The second is
$(n-1)!\coeff{t}{n-1}((tX''-X'+x_1)F)$, and the last is
$(n-1)!\coeff{t}{n-1}((X'-x_1)F)$.  Their sum is therefore
\[
 (n-1)!\coeff{t}{n-1}t(X''+X'^2)F
 =(n-1)!\coeff{t}{n-1}tF''=(n-1)\BellC n.
\]
Division by $n-1$ proves the formula.
\end{proof}

\chapter{Specializations, inversions, determinants, and convolutions}

\section{Stirling, Bell, Lah, and Touchard specializations}
\begin{theorem}\label{thm:bell-poly-specializations}
For admissible $n,k$,
\begin{align}
 \BellP nk(0!,1!,\ldots,(n-k)!)&=\stirone nk,
 \label{eq:bell-first-specialization}\\
 \BellC n(0!,1!,\ldots,(n-1)!)&=n!,
 \label{eq:bell-factorial-complete}\\
 \BellP nk(1,1,\ldots,1)&=\stirtwo nk,
 \label{eq:bell-second-specialization}\\
 \BellC n(1,1,\ldots,1)&=\Bell n,
 \label{eq:bell-number-specialization}\\
 \BellP nk(1!,2!,\ldots,(n-k+1)!)
 &=\binom{n-1}{k-1}\frac{n!}{k!}=L(n,k),
 \label{eq:bell-lah-specialization}\\
 T_n(x)&=\BellC n(x,x,\ldots,x)
 =\sum_{k=0}^{n}\stirtwo nkx^k.
 \label{eq:touchard-bell-specialization}
\end{align}
\end{theorem}
\begin{proof}
In \eqref{eq:partial-bell-egf}, the choices $x_j=(j-1)!$, $x_j=1$, and
$x_j=j!$ make $X(t)$ respectively
\[
 -\log(1-t),\qquad e^t-1,\qquad \frac{t}{1-t}.
\]
The resulting coefficient formulae are exactly the column EGFs for unsigned
first-kind Stirling numbers, second-kind Stirling numbers, and Lah numbers proved
earlier.  Summing over $k$ yields \eqref{eq:bell-factorial-complete} and
\eqref{eq:bell-number-specialization}.  Finally $x_j=x$ gives
$\exp(x(e^t-1))$, the EGF of $T_n(x)$.
\end{proof}

\section{Homogeneity and a useful evaluation}
\begin{theorem}[Bigraded homogeneity]\label{thm:bell-bihomogeneous}
For scalars $\alpha,\beta$,
\begin{equation}\label{eq:bell-bihomogeneous}
 \BellP nk(\alpha\beta x_1,\alpha\beta^2x_2,\ldots)
 =\alpha^k\beta^n\BellP nk(x_1,x_2,\ldots).
\end{equation}
Moreover,
\begin{equation}\label{eq:bell-linear-arguments}
 \BellP nk(1,2,3,\ldots,n-k+1)=\binom nk k^{n-k}.
\end{equation}
\end{theorem}
\begin{proof}
Every monomial in \eqref{eq:partial-bell-definition} has
$\sum j_i=k$ and $\sum ij_i=n$, proving the first formula.  For the second,
$x_j=j$ gives $X(t)=te^t$; hence
\[
 \frac{X(t)^k}{k!}=\frac{t^ke^{kt}}{k!}.
\]
The coefficient of $t^n/n!$ is
$n!k^{n-k}/(k!(n-k)!)=\binom nk k^{n-k}$.
\end{proof}

\section{Addition and shifted-zero identities}
\begin{theorem}[Complete Bell addition formula]\label{thm:complete-bell-addition}
\begin{equation}\label{eq:complete-bell-addition}
 \BellC n(x_1+y_1,\ldots,x_n+y_n)
 =\sum_{i=0}^{n}\binom ni\BellC{n-i}(x)\BellC i(y).
\end{equation}
\end{theorem}
\begin{proof}
The complete EGF for the left side is
$\exp(X(t)+Y(t))=\exp X(t)\exp Y(t)$.  Coefficient extraction gives the binomial
convolution.
\end{proof}

\begin{theorem}[Insertion of $q$ leading zeros]\label{thm:bell-leading-zeros}
For $q\ge0$,
\begin{multline}\label{eq:bell-leading-zeros}
 \BellP nk\!\left(
 \frac{x_{q+1}}{\binom{q+1}{q}},
 \frac{x_{q+2}}{\binom{q+2}{q}},\ldots\right)\\
 =\frac{n!(q!)^k}{(n+qk)!}
 \BellP{n+qk}{k}(\underbrace{0,\ldots,0}_{q\text{ entries}},
 x_{q+1},x_{q+2},\ldots).
\end{multline}
\end{theorem}
\begin{proof}
Since
\[
 \frac{x_{q+i}}{\binom{q+i}{q}}\frac{t^i}{i!}
 =q!\,\frac{x_{q+i}t^i}{(q+i)!},
\]
the $k$th power of the left component series is $(q!)^k t^{-qk}$ times the
$k$th power of the component series on the right.  Comparing the coefficient of
$t^n$ in \eqref{eq:partial-bell-egf} gives the factorial factor displayed.
\end{proof}

\section{Exponential--logarithmic inversion}
\begin{theorem}[Bell transform and its inverse]\label{thm:bell-transform-inverse}
Suppose
\begin{equation}\label{eq:bell-transform-y}
 y_n=\BellC n(x_1,\ldots,x_n)=\sum_{k=1}^{n}\BellP nk(x_1,\ldots).
\end{equation}
Then
\begin{equation}\label{eq:bell-transform-x}
 x_n=\sum_{k=1}^{n}(-1)^{k-1}(k-1)!
 \BellP nk(y_1,\ldots,y_{n-k+1}).
\end{equation}
\end{theorem}
\begin{proof}
Let $X(t)=\sum_{n\ge1}x_nt^n/n!$ and
$Y(t)=\sum_{n\ge1}y_nt^n/n!$.  Equation
\eqref{eq:complete-bell-egf} says $1+Y=e^X$, so
\[
 X=\log(1+Y)=\sum_{k\ge1}\frac{(-1)^{k-1}}{k}Y^k.
\]
Use \eqref{eq:partial-bell-egf} to extract the coefficient of $t^n/n!$.
\end{proof}

More generally, if $f$ and $g=f^{-1}$ are locally inverse and
\begin{equation}\label{eq:general-bell-transform}
 y_n=\sum_{k=0}^{n}f^{(k)}(a)\BellP nk(x_1,\ldots),
\end{equation}
then
\begin{equation}\label{eq:general-bell-inverse}
 x_n=\sum_{k=0}^{n}g^{(k)}(f(a))\BellP nk(y_1,\ldots).
\end{equation}
Indeed, these are the derivative coefficients of $f\circ u$ and
$g\circ f\circ u=u$; a complete proof is given with Fa\`a di Bruno's formula in
\cref{eq:faa-bell}.

\section{Two Hessenberg determinants}
For $n\ge1$, define $H_n=(h_{ij})$ and $K_n=(k_{ij})$ by
\[
 h_{ij}=\begin{cases}
 \binom{n-i}{j-i}x_{j-i+1},&j\ge i,\\
 -1,&j=i-1,\\0,&j<i-1,
 \end{cases}
 \qquad
 k_{ij}=\begin{cases}
 x_{j-i+1}/(j-i)!,&j\ge i,\\
 -(i-1),&j=i-1,\\0,&j<i-1.
 \end{cases}
\]
Thus $H_n$ begins with
\[
\begin{bmatrix}
 x_1&\binom{n-1}{1}x_2&\binom{n-1}{2}x_3&\cdots&x_n\\
 -1&x_1&\binom{n-2}{1}x_2&\cdots&x_{n-1}\\
 0&-1&x_1&\cdots&x_{n-2}\\
 \vdots&&\ddots&\ddots&\vdots\\
 0&\cdots&0&-1&x_1
\end{bmatrix},
\]
while $K_n$ has first row
$x_1,x_2/1!,\ldots,x_n/(n-1)!$ and subdiagonal
$-1,-2,\ldots,-(n-1)$.

\begin{theorem}[Bell determinants]\label{thm:bell-determinants}
\begin{equation}\label{eq:bell-determinants}
 \BellC n(x_1,\ldots,x_n)=\det H_n=\det K_n.
\end{equation}
\end{theorem}
\begin{proof}
Expansion of the upper Hessenberg determinant $\det H_n$ along its first row gives
\[
 D_n=\sum_{j=1}^{n}\binom{n-1}{j-1}x_jD_{n-j},\qquad D_0=1;
\]
the cofactor sign is cancelled by the product of the $j-1$ entries $-1$ on the
subdiagonal.  This is exactly \eqref{eq:complete-bell-recurrence}, so
$D_n=\BellC n$.

Let $J$ reverse the coordinate order and let
$D=\diag(0!,1!,\ldots,(n-1)!)$.  A direct entrywise calculation gives
\[
 K_n=D(JH_n^{\mathsf T}J)D^{-1}.
\]
Indeed,
$\frac{(i-1)!}{(j-1)!}\binom{j-1}{i-1}=1/(j-i)!$ above the diagonal, and the
subdiagonal becomes $-(i-1)$.  Similarity, transposition, and reversal preserve
the determinant, proving the second equality.
\end{proof}

\section{The diamond convolution}
For sequences $x=(x_1,x_2,\ldots)$ and $y=(y_1,y_2,\ldots)$ define
\begin{equation}\label{eq:diamond}
 (x\mathbin\diamondsuit y)_n=\sum_{j=1}^{n-1}\binom njx_jy_{n-j}.
\end{equation}
This operation is commutative and associative because it corresponds to
multiplication of the EGFs $X(t)$ and $Y(t)$.

\begin{theorem}\label{thm:diamond-bell}
If $x^{k\diamondsuit}$ denotes the $k$-fold diamond power, then
\begin{equation}\label{eq:diamond-bell}
 \BellP nk(x_1,\ldots,x_{n-k+1})=\frac{(x^{k\diamondsuit})_n}{k!}.
\end{equation}
\end{theorem}
\begin{proof}
The EGF of $x^{k\diamondsuit}$ is $X(t)^k$.  Compare with
\eqref{eq:partial-bell-egf}.
\end{proof}
For example,
\[
 x\diamondsuit x=(0,2x_1^2,6x_1x_2,8x_1x_3+6x_2^2,\ldots)
\]
and
$(x^{3\diamondsuit})_4=36x_1^2x_2$, whence
$\BellP43=6x_1^2x_2$.

\chapter{Applications of Bell polynomials}

\section{Symmetric functions and Newton identities}
Let $e_n$ be the elementary symmetric function and
$p_n=\sum_i u_i^n$ the power sum in an alphabet $u_1,u_2,\ldots$.
\begin{theorem}\label{thm:bell-symmetric-functions}
\begin{align}
 e_n&=\frac1{n!}\BellC n(p_1,-1!p_2,2!p_3,\ldots,(-1)^{n-1}(n-1)!p_n)
 \label{eq:elementary-via-bell}\\
 &=\frac{(-1)^n}{n!}\BellC n(-p_1,-1!p_2,-2!p_3,\ldots,-(n-1)!p_n),
 \notag\\
 p_n&=\frac{(-1)^{n-1}}{(n-1)!}
 \sum_{k=1}^{n}(-1)^{k-1}(k-1)!
 \BellP nk(1!e_1,2!e_2,\ldots)
 \label{eq:powersum-via-bell}\\
 &=(-1)^n n\sum_{k=1}^{n}\frac1k
 \OBellP nk(-e_1,-e_2,\ldots).
 \notag
\end{align}
\end{theorem}
\begin{proof}
The classical product identity
\[
 E(t)=\sum_{n\ge0}e_nt^n=\prod_i(1+u_it)
 =\exp\!\left(\sum_{r\ge1}(-1)^{r-1}p_r\frac{t^r}{r}\right)
\]
and \eqref{eq:complete-bell-egf} give \eqref{eq:elementary-via-bell}.
Taking logarithms,
$\sum_{r\ge1}(-1)^{r-1}p_rt^r/r=\log(1+\sum_{j\ge1}e_jt^j)$, and expanding
the logarithm through ordinary Bell polynomials yields the second pair of forms.
The exponential form follows from \eqref{eq:ordinary-exponential-scaling}.
\end{proof}

\begin{corollary}[Determinant from traces]\label{cor:det-traces-bell}
For an $n\times n$ matrix $A$, set
$s_k=-(k-1)!\tr(A^k)$.  Then
\begin{equation}\label{eq:det-traces-bell}
 \det A=\frac{(-1)^n}{n!}\BellC n(s_1,s_2,\ldots,s_n).
\end{equation}
\end{corollary}
\begin{proof}
Apply \eqref{eq:elementary-via-bell} to the eigenvalues of $A$, using
$e_n=\det A$ and $p_k=\tr(A^k)$.  Both sides are polynomial in the matrix entries,
so the argument remains valid for nondiagonalizable matrices.
\end{proof}

\section{Cycle indices}
Let $a_j$ mark cycles of length $j$ in a permutation.
\begin{theorem}\label{thm:cycle-index-bell}
The cycle index of the symmetric group is
\begin{equation}\label{eq:cycle-index-bell}
 Z(S_n)=\frac1{n!}\BellC n(0!a_1,1!a_2,\ldots,(n-1)!a_n).
\end{equation}
\end{theorem}
\begin{proof}
A permutation is a set of cycles.  A labelled cycle of size $j$ has $(j-1)!$
linear realizations, so its component EGF is $(j-1)!a_jt^j/j!=a_jt^j/j$.
The exponential formula gives
$\exp(\sum_{j\ge1}a_jt^j/j)$; comparison with
\eqref{eq:complete-bell-egf} proves the identity after division by $n!$.
\end{proof}

\section{Moments and cumulants}
Let $M(t)=\E e^{tX}=\sum_{n\ge0}\mu'_nt^n/n!$ and
$K(t)=\log M(t)=\sum_{n\ge1}\kappa_nt^n/n!$.
\begin{theorem}[Moment--cumulant relations]\label{thm:moment-cumulant}
\begin{align}
 \mu'_n&=\BellC n(\kappa_1,\ldots,\kappa_n)
 =\sum_{k=1}^{n}\BellP nk(\kappa_1,\ldots),
 \label{eq:moments-from-cumulants}\\
 \kappa_n&=\sum_{k=1}^{n}(-1)^{k-1}(k-1)!
 \BellP nk(\mu'_1,\ldots,\mu'_{n-k+1}).
 \label{eq:cumulants-from-moments}
\end{align}
\end{theorem}
\begin{proof}
The identity $M=e^K$ is \eqref{eq:complete-bell-egf}; its logarithmic inverse is
\cref{thm:bell-transform-inverse}.
\end{proof}

\section{Hermite polynomials}
The probabilists' Hermite polynomials are defined by
\[
 e^{xt-t^2/2}=\sum_{n\ge0}\operatorname{He}_n(x)\frac{t^n}{n!}.
\]
Comparison with \eqref{eq:complete-bell-egf} gives and proves
\begin{equation}\label{eq:hermite-bell}
 \operatorname{He}_n(x)=\BellC n(x,-1,0,\ldots,0).
\end{equation}
In particular, the usual recurrence and differential equation follow by
differentiating this generating function.

\section{Polynomial sequences of binomial type}
Let
$h(t)=\sum_{j\ge1}a_jt^j/j!$ with $a_1\ne0$, and define
\begin{equation}\label{eq:binomial-type-bell}
 p_n(x)=\BellC n(a_1x,a_2x,\ldots,a_nx)
 =\sum_{k=0}^{n}\BellP nk(a_1,\ldots,a_{n-k+1})x^k.
\end{equation}

\begin{theorem}\label{thm:binomial-type-bell}
The polynomials $p_n$ satisfy
\begin{align}
 \sum_{n\ge0}p_n(x)\frac{t^n}{n!}&=e^{xh(t)},
 \label{eq:binomial-type-egf}\\
 p_n(x+y)&=\sum_{k=0}^{n}\binom nkp_k(x)p_{n-k}(y).
 \label{eq:binomial-type-identity}
\end{align}
If $Q=h^{-1}(D_x)$ is the associated delta operator, then
\begin{equation}\label{eq:delta-operator}
 Qp_n(x)=np_{n-1}(x).
\end{equation}
\end{theorem}
\begin{proof}
Substitute $x a_j$ in \eqref{eq:complete-bell-egf}.  The product
$e^{(x+y)h}=e^{xh}e^{yh}$ gives the binomial identity.  Finally
$h^{-1}(D_x)e^{xh(t)}=h^{-1}(h(t))e^{xh(t)}=t e^{xh(t)}$; coefficient comparison
gives \eqref{eq:delta-operator}.
\end{proof}

\begin{auditbox}
In some sources the symbol $B_n(x)$ denotes the Touchard polynomial, in others a
complete Bell polynomial, a Bernoulli polynomial, or a binary-tree sequence.  The
notations $T_n$, $\BellC n$, $\beta_n(x)$, and $C_n$ are kept separate here; this
is essential when the objects occur in the same inversion formula.
\end{auditbox}

\part{Eulerian Numbers and Descent Calculus}

\chapter{Ordinary Eulerian numbers}
\index{Eulerian numbers}

\section{Descents, symmetry, and the basic recurrence}
For a permutation $\pi=\pi_1\cdots\pi_n$, a descent is an index
$i\in\{1,\ldots,n-1\}$ with $\pi_i>\pi_{i+1}$.
\begin{definition}
The Eulerian number $\eulerian nk=A(n,k)$ counts permutations of $[n]$ with
exactly $k$ descents.  Put
\[
 A_0(t)=1,
 \qquad
 A_n(t)=\sum_{k=0}^{n-1}\eulerian nk t^k\quad(n\ge1).
\]
Outside $0\le k<n$ the number is zero.
\end{definition}

\begin{theorem}[Eulerian recurrence]\label{thm:eulerian-recurrence}
\begin{equation}\label{eq:eulerian-recurrence}
 \eulerian nk=(k+1)\eulerian{n-1}{k}
 +(n-k)\eulerian{n-1}{k-1}.
\end{equation}
Equivalently,
\begin{equation}\label{eq:eulerian-poly-recurrence}
 A_n(t)=\bigl(1+(n-1)t\bigr)A_{n-1}(t)
       +t(1-t)A_{n-1}'(t).
\end{equation}
\end{theorem}
\begin{proof}
Insert $n$ into a permutation of $[n-1]$.  If the old permutation has $k$
descents, insertion after one of those $k$ descents or at the end preserves the
number of descents, giving $k+1$ choices.  If it has $k-1$ descents, insertion in
any of the remaining $n-k$ slots creates one new descent.  This proves
\eqref{eq:eulerian-recurrence}.  Multiply by $t^k$, sum over $k$, and collect the
terms containing $k$ to obtain \eqref{eq:eulerian-poly-recurrence}.
\end{proof}

Complementation $\pi_i\mapsto n+1-\pi_i$ interchanges descents and ascents.
Therefore
\begin{equation}\label{eq:eulerian-symmetry}
 \eulerian nk=\eulerian n{n-1-k},
 \qquad
 A_n(t)=t^{n-1}A_n(t^{-1}).
\end{equation}
In particular $\eulerian n0=\eulerian n{n-1}=1$ and
$\sum_k\eulerian nk=n!$.  The next diagonal is
\begin{equation}\label{eq:eulerian-k1}
 \eulerian n1=2^n-(n+1),
\end{equation}
which follows either from the recurrence or from the explicit formula below.

\section{Exponential generating function and explicit formula}
\begin{theorem}[Eulerian EGF]\label{thm:eulerian-egf}
\begin{equation}\label{eq:eulerian-egf}
 \sum_{n\ge0}A_n(t)\frac{x^n}{n!}
 =\frac{t-1}{t-e^{(t-1)x}}
 =\frac{(1-t)e^{(1-t)x}}{1-te^{(1-t)x}}.
\end{equation}
\end{theorem}
\begin{proof}
Let $F(x,t)$ denote the left side.  Multiplying
\eqref{eq:eulerian-poly-recurrence} by $x^{n-1}/(n-1)!$ and summing gives the
first-order equation
\[
 (1-tx)F_x=t(1-t)F_t+F,
 \qquad F(0,t)=1.
\]
Direct substitution verifies that the displayed rational-exponential function
solves this equation and initial condition.  Formal uniqueness follows by
recursively comparing coefficients of $x$.
\end{proof}

Expanding the denominator geometrically gives the central power-series identity.
\begin{theorem}\label{thm:eulerian-power-series}
For $n\ge0$,
\begin{equation}\label{eq:eulerian-power-series}
 \sum_{m=0}^{\infty}(m+1)^n t^m
 =\frac{A_n(t)}{(1-t)^{n+1}}.
\end{equation}
Consequently,
\begin{equation}\label{eq:eulerian-explicit}
 \eulerian nk=\sum_{j=0}^{k}(-1)^j\binom{n+1}{j}(k+1-j)^n.
\end{equation}
\end{theorem}
\begin{proof}
From \eqref{eq:eulerian-egf},
\[
 F(x,t)=(1-t)\sum_{m\ge0}t^m e^{(m+1)(1-t)x}.
\]
Extracting $x^n/n!$ proves \eqref{eq:eulerian-power-series}.  Multiply by
$(1-t)^{n+1}$ and extract $t^k$ to obtain
\eqref{eq:eulerian-explicit}.  Formula \eqref{eq:eulerian-k1} is its case $k=1$.
\end{proof}

\begin{theorem}[A second polynomial recurrence]\label{thm:eulerian-binomial-recurrence}
For $n>1$,
\begin{equation}\label{eq:eulerian-binomial-recurrence}
 A_n(t)=\sum_{k=0}^{n-1}\binom nk A_k(t)(t-1)^{n-1-k}.
\end{equation}
\end{theorem}
\begin{proof}
The EGF of the right sides, summed over $n\ge1$, is
\[
 F(x,t)\frac{e^{(t-1)x}-1}{t-1}.
\]
Using \eqref{eq:eulerian-egf}, this equals $F(x,t)-1$, which is the EGF of the
left sides.
\end{proof}

\section{Worpitzky's identity and sums of powers}
\begin{theorem}[Worpitzky]\label{thm:worpitzky}
For every indeterminate $x$ and $n\ge1$,
\begin{equation}\label{eq:worpitzky}
 x^n=\sum_{k=0}^{n-1}\eulerian nk\binom{x+k}{n}.
\end{equation}
\end{theorem}
\begin{proof}
The coefficient of $t^{x-1}$ in
$t^{-k}(1-t)^{-n-1}$ is $\binom{x+k}{n}$ for positive integer $x$.
Equivalently, taking the coefficient of $t^{x-1}$ in
\eqref{eq:eulerian-power-series} after multiplying by the numerator gives
\eqref{eq:worpitzky}.  Since both sides are polynomials of degree $n$ in $x$ and
agree for all positive integers, they agree identically.

A direct combinatorial proof standardizes a word in $[x]^n$: its stable sorting
permutation records the relative order, and its $k$ descents specify $k$ mandatory
separations among the $n$ sorted positions.  Stars and bars gives
$\binom{x+k}{n}$ words for each permutation with $k$ descents.
\end{proof}

\begin{corollary}[Power sums]\label{cor:eulerian-power-sum}
For $m,n\ge1$,
\begin{equation}\label{eq:eulerian-power-sum}
 \sum_{r=1}^{m}r^n
 =\sum_{k=0}^{n-1}\eulerian nk\binom{m+k+1}{n+1}.
\end{equation}
\end{corollary}
\begin{proof}
Sum \eqref{eq:worpitzky} over $x=1,\ldots,m$ and use the hockey-stick identity
$\sum_{x=1}^{m}\binom{x+k}{n}=\binom{m+k+1}{n+1}$; the omitted lower boundary
terms vanish for $0\le k<n$.
\end{proof}

\section{Negative polylogarithms}
For $|z|<1$, $\Li_{-n}(z)=\sum_{m\ge1}m^nz^m$.  Shifting $m$ in
\eqref{eq:eulerian-power-series} gives
\begin{equation}\label{eq:negative-polylog-eulerian}
 \Li_{-n}(z)=\frac{zA_n(z)}{(1-z)^{n+1}}
 =\frac{1}{(1-z)^{n+1}}
   \sum_{k=0}^{n-1}\eulerian nk z^{n-k},
\end{equation}
where the second equality uses symmetry \eqref{eq:eulerian-symmetry}.  Since the
right side is rational, it also supplies the meromorphic continuation of the
negative polylogarithm.

\section{Eulerian and Stirling bases}
\begin{theorem}\label{thm:eulerian-stirling}
\begin{equation}\label{eq:eulerian-stirling}
 A_n(t)=\sum_{k=0}^{n}\stirtwo nk k!(t-1)^{n-k}.
\end{equation}
\end{theorem}
\begin{proof}
Use $m^n=\sum_k\stirtwo nk\falling mk$ in
$\sum_{m\ge0}m^nt^m$.  Since
$\sum_{m\ge0}\falling mk t^m=k!t^k/(1-t)^{k+1}$, multiplication by
$(1-t)^{n+1}/t$ yields
$A_n(t)=\sum_k\stirtwo nk k!t^{k-1}(1-t)^{n-k}$.  Replacing $t$ by $1/t$ and
using \eqref{eq:eulerian-symmetry}, or simply expanding around $t=1$, gives the
equivalent displayed form.
\end{proof}

\section{An Irwin--Hall integral and hypersimplex volumes}
\begin{theorem}\label{thm:eulerian-irwin-hall}
For every function $f:\R\to\C$ for which the following finite sum is defined,
\begin{equation}\label{eq:eulerian-irwin-hall}
 \sum_{k=0}^{n-1}\eulerian nk f(k)
 =n!\int_{[0,1]^n}f\!\left(\left\lfloor x_1+\cdots+x_n\right\rfloor\right)
 \dd x_1\cdots\dd x_n.
\end{equation}
In particular, the slab
\[
 \Delta_{n,k}=\{x\in[0,1]^n:k\le x_1+\cdots+x_n<k+1\}
\]
has volume $\eulerian nk/n!$.
\end{theorem}
\begin{proof}
The volume of $\{x\in[0,1]^n:\sum x_i<y\}$ is obtained from the simplex
$\{x_i\ge0:\sum x_i<y\}$ by inclusion--exclusion over coordinates exceeding $1$:
\[
 V_n(y)=\frac1{n!}\sum_{j=0}^{\lfloor y\rfloor}(-1)^j\binom nj(y-j)^n.
\]
Thus the volume of the $k$th slab is $V_n(k+1)-V_n(k)$.  Combining adjacent terms
by Pascal's identity gives
\[
 \frac1{n!}\sum_{j=0}^{k}(-1)^j\binom{n+1}{j}(k+1-j)^n
 =\frac1{n!}\eulerian nk
\]
by \eqref{eq:eulerian-explicit}.  Summing $f(k)$ times these volumes proves the
integral identity.
\end{proof}

\section{Alternating sums and Bernoulli numbers}
\begin{theorem}\label{thm:eulerian-alternating}
For $n>0$,
\begin{equation}\label{eq:eulerian-alternating}
 \sum_{k=0}^{n-1}(-1)^k\eulerian nk
 =2^{n+1}(2^{n+1}-1)\frac{\beta_{n+1}}{n+1}.
\end{equation}
For $n>1$,
\begin{align}
 \sum_{k=0}^{n-1}\frac{(-1)^k\eulerian nk}{\binom{n-1}{k}}&=0,
 \label{eq:eulerian-reciprocal-zero}\\
 \sum_{k=0}^{n-1}\frac{(-1)^k\eulerian nk}{\binom nk}&=(n+1)\beta_n.
 \label{eq:eulerian-reciprocal-bernoulli}
\end{align}
\end{theorem}
\begin{proof}
At $t=-1$, \eqref{eq:eulerian-egf} becomes
$2/(1+e^{-2x})=1+\tanh x$.  The Bernoulli generating function gives
\[
 \tanh x=\sum_{n\ge1}
 2^{n+1}(2^{n+1}-1)\frac{\beta_{n+1}}{n+1}\frac{x^n}{n!},
\]
which proves \eqref{eq:eulerian-alternating}.

For the remaining identities use
\[
 \binom{n-1}{k}^{-1}=n\int_0^1u^k(1-u)^{n-k-1}\dd u,
 \quad
 \binom nk^{-1}=(n+1)\int_0^1u^k(1-u)^{n-k}\dd u.
\]
Substitute $t=-u/(1-u)$ in \eqref{eq:eulerian-stirling}; then
\begin{equation}\label{eq:eulerian-beta-transform}
 (1-u)^nA_n\!\left(-\frac{u}{1-u}\right)
 =\sum_{j=0}^{n}(-1)^{n-j}\stirtwo njj!(1-u)^j.
\end{equation}
For the first integral divide by $1-u$ and integrate.  The result is
$n\sum_{j=1}^n(-1)^{n-j}\stirtwo nj(j-1)!=0$ for $n>1$, because it is the
coefficient of $x^n/n!$ in
$\log(1+(e^x-1))=x$.  For the second integral one obtains
$(n+1)\sum_j(-1)^{n-j}\stirtwo njj!/(j+1)$.  The standard Stirling formula
$\beta_n=\sum_j(-1)^jj!\stirtwo nj/(j+1)$ proves the claim; for odd $n>1$ both
sides vanish, and for even $n$ the signs coincide.
\end{proof}

\chapter{Geometric interpretations}

\section{The cube and its Ehrhart numerator}
For a lattice polytope $P$ of dimension $d$, its Ehrhart series has the form
$\sum_{m\ge0}|mP\cap\Z^d|t^m=h^*(t)/(1-t)^{d+1}$.
For the unit cube $C=[0,1]^n$ one has $|mC\cap\Z^n|=(m+1)^n$; hence
\eqref{eq:eulerian-power-series} proves:
\begin{theorem}\label{unif:thm-cube-ehrhart}
The $h^*$-polynomial of the standard $n$-cube is $A_n(t)$, so its $h^*$-vector is
the $n$th Eulerian row.
\end{theorem}
\begin{proof}
The dilate $m[0,1]^n$ contains $(m+1)^n$ lattice points. Therefore its Ehrhart
series is
\[
 \sum_{m\ge0}(m+1)^nt^m.
\]
The Eulerian power-series identity, with the summation index shifted by one,
gives
\[
 \sum_{m\ge0}(m+1)^nt^m=\frac{A_n(t)}{(1-t)^{n+1}}.
\]
Comparing this with the defining Ehrhart form
$\sum_{m\ge0}|mC\cap\Z^n|t^m=h^*(t)/(1-t)^{n+1}$ proves
$h^*(t)=A_n(t)$.
\end{proof}

The regions $\Delta_{n,k}$ in \cref{thm:eulerian-irwin-hall} are half-open
hypersimplices.  Their normalized volumes are $n!\operatorname{vol}\Delta_{n,k}
=\eulerian nk$.

\section{The permutohedron}
The type-$A$ permutohedron is the convex hull of the permutations of
$(1,2,\ldots,n)$; its dimension is $n-1$.  Faces of its braid fan are indexed by
ordered set partitions.  There are $k!\stirtwo nk$ ordered partitions into $k$
blocks.

\begin{theorem}
The $h$-polynomial of the simple polytope dual to the type-$A$ permutohedron is
$A_n(t)$.
\end{theorem}
\begin{proof}
For a simple $(n-1)$-polytope the $h$-polynomial is the transform
$h(t)=\sum_F(t-1)^{\dim F}$ over faces of the dual simplicial complex.  Grouping
braid-fan faces by their $k$ ordered blocks gives
$\sum_k k!\stirtwo nk(t-1)^{n-k}$, which is $A_n(t)$ by
\eqref{eq:eulerian-stirling}.
\end{proof}

\chapter{Type-B and second-order Eulerian numbers}

\section{Signed permutations and the corrected type-B indexing}
A signed permutation is a word $\pi_1\cdots\pi_n$ whose absolute values form a
permutation of $[n]$.  Put $\pi_0=0$ and call
$i\in\{0,\ldots,n-1\}$ a type-$B$ descent when $\pi_i>\pi_{i+1}$.
Let $E_B(n,k)$ count signed permutations with $k$ such descents, and set
$M_n(t)=\sum_{k=0}^{n}E_B(n,k)t^k$.

\begin{auditbox}
The archived display reverses the two entries in the descent inequality and pairs
that convention with an incompatible explicit formula.  The definition above is
the standard Coxeter descent convention.  It gives the archived table
$1,6,1$ for $n=2$ and the midpoint-Eulerian generating function below.
\end{auditbox}

\begin{theorem}[Type-$B$ recurrence and Worpitzky series]\label{thm:typeB-eulerian}
\begin{align}
 E_B(n,k)&=(2k+1)E_B(n-1,k)+(2n-2k+1)E_B(n-1,k-1),
 \label{eq:typeB-recurrence}\\
 \sum_{m\ge0}(2m+1)^nt^m&=\frac{M_n(t)}{(1-t)^{n+1}},
 \label{eq:typeB-power-series}\\
 E_B(n,k)&=\sum_{j=0}^{k}(-1)^j\binom{n+1}{j}
             \bigl(2(k-j)+1\bigr)^n.
 \label{eq:typeB-explicit}
\end{align}
The first rows are
\[
1;\quad 1,1;\quad1,6,1;\quad1,23,23,1;\quad1,76,230,76,1.
\]
\end{theorem}
\begin{proof}
Insert either $n$ or $-n$ into a signed permutation of size $n-1$.  Marking the
slots according to whether insertion preserves or creates a descent gives
$2k+1$ preserving slots for an object with $k$ descents and
$2n-2k+1$ creating slots for one with $k-1$ descents; this proves the recurrence.

A type-$B$ standardization partitions the $(2m+1)^n$ words in the ordered alphabet
$-m<\cdots<-1<0<1<\cdots<m$ by their signed sorting permutation.  If that
permutation has $k$ type-$B$ descents, stars and bars gives
$\binom{m+n-k}{n}$ compatible weakly increasing absolute levels.  Hence
\[
 (2m+1)^n=\sum_{k=0}^{n}E_B(n,k)\binom{m+n-k}{n}.
\]
Multiply by $t^m$, sum in $m$, and use
$\sum_{m\ge0}\binom{m+n-k}{n}t^m=t^k/(1-t)^{n+1}$ to obtain
\eqref{eq:typeB-power-series}.  Multiplication by $(1-t)^{n+1}$ and coefficient
extraction gives \eqref{eq:typeB-explicit}.
\end{proof}

As in type $A$, these coefficients form the $h$-vector of the simple polytope dual
to the type-$B$ permutohedron; the proof repeats the braid-fan argument with signed
ordered partitions.

\section{Stirling permutations and second-order Eulerian numbers}
A \emph{Stirling permutation of order $n$} is a permutation of the multiset
$\{1,1,2,2,\ldots,n,n\}$ such that all entries between the two copies of $i$ are
larger than $i$.  Let $\eulertwo nk$ count those with $k$ ordinary descents and put
\[
 P_n(t)=\sum_{k=0}^{n-1}\eulertwo nk t^k,
 \qquad P_0(t)=1.
\]

\begin{theorem}[Second-order recurrence]\label{thm:second-eulerian-recurrence}
\begin{align}
 \eulertwo nk&=(2n-k-1)\eulertwo{n-1}{k-1}
 +(k+1)\eulertwo{n-1}{k},
 \label{eq:second-eulerian-recurrence}\\
 P_{n+1}(t)&=(2nt+1)P_n(t)-t(t-1)P_n'(t),
 \label{eq:second-eulerian-poly-recurrence}\\
 P_n(1)&=(2n-1)!!.
 \label{eq:second-eulerian-row-sum}
\end{align}
\end{theorem}
\begin{proof}
The two copies of the largest label $n$ must be adjacent.  Insert the pair $nn$
into one of the $2n-1$ gaps of a Stirling permutation of order $n-1$.  A descent
gap or the terminal gap preserves the number of descents, giving $k+1$ choices
from an object with $k$ descents.  Each of the remaining $2n-k-1$ gaps creates one
descent from an object with $k-1$ descents.  This is the first recurrence; summing
against $t^k$ gives the second.  At each stage there are $2n-1$ insertion gaps, so
the total count is $(2n-1)!!$.
\end{proof}

The differential recurrence has two useful equivalent forms:
\begin{align}
 (t-1)^{-2n-2}P_{n+1}(t)
 &=\left(t(1-t)^{-2n-1}P_n(t)\right)',
 \label{eq:second-eulerian-differential-form}\\
 u_{n+1}(t)&=\left(\frac{t}{1-t}u_n(t)\right)',
 \quad u_n(t)=(t-1)^{-2n}P_n(t),\quad u_0=1.
 \label{eq:second-eulerian-u}
\end{align}
They follow by the product rule from
\eqref{eq:second-eulerian-poly-recurrence}.

\begin{theorem}[Diagonal Stirling generating function]\label{thm:second-eulerian-stirling-gf}
For $n\ge1$,
\begin{equation}\label{eq:second-eulerian-stirling-gf}
 \sum_{m=0}^{\infty}\stirtwo{n+m}{m}t^m
 =\frac{tP_n(t)}{(1-t)^{2n+1}}.
\end{equation}
\end{theorem}
\begin{proof}
Let $F_n(t)$ denote the left side.  The second-kind recurrence gives
$(1-t)F_{n+1}=tF_n'$, with $F_0=(1-t)^{-1}$.  Hence
$F_1=t/(1-t)^3$.  If the displayed formula holds for $n$, differentiating
$tP_n/(1-t)^{2n+1}$ and using
\eqref{eq:second-eulerian-poly-recurrence} gives the formula for $n+1$.
\end{proof}


\section{Second-order Eulerian interpolation of first-kind diagonals}
The same second-order Eulerian row that controls
\eqref{eq:second-eulerian-stirling-gf} also gives a binomial-basis expansion for
the near-diagonals of the unsigned first-kind triangle.

\begin{theorem}[Second-order Eulerian diagonal formula]
\label{thm:second-eulerian-first-diagonal}
For $p,m\ge0$,
\begin{equation}
\label{eq:second-eulerian-first-diagonal}
 \stirone m{m-p}
 =\sum_{j=0}^{p}\eulertwo pj\binom{m+j}{2p},
\end{equation}
where an out-of-range first-kind number is zero.  Equivalently, the polynomial
continuation in a variable $x$ satisfies
\[
 \stirone{x}{x-p}
 =\sum_{j=0}^{p}\eulertwo pj\binom{x+j}{2p}.
\]
The first instances are
\begin{align*}
 \stirone{x}{x-1}&=\binom{x}{2},\\
 \stirone{x}{x-2}&=\binom{x}{4}+2\binom{x+1}{4},\\
 \stirone{x}{x-3}&=\binom{x}{6}+8\binom{x+1}{6}
                    +6\binom{x+2}{6}.
\end{align*}
\end{theorem}
\begin{proof}
For fixed $p$ define
\[
 D_p(t)=\sum_{m\ge0}\stirone m{m-p}t^m.
\]
Writing $a_{m,p}=\stirone m{m-p}$, the first-kind recurrence becomes
$a_{m,p}=a_{m-1,p}+(m-1)a_{m-1,p-1}$.  Therefore
\begin{equation}
\label{eq:first-diagonal-ogf-recurrence}
 (1-t)D_p(t)=t^2D_{p-1}'(t),
 \qquad D_0(t)=\frac1{1-t}.
\end{equation}
Let
\[
 Q_p(t)=t^{2p}P_p(t^{-1}).
\]
The polynomial recurrence \eqref{eq:second-eulerian-poly-recurrence} is equivalent,
after the substitution $t\mapsto t^{-1}$, to
\[
 Q_p(t)=t^2\bigl((1-t)Q_{p-1}'(t)+(2p-1)Q_{p-1}(t)\bigr).
\]
It follows inductively from \eqref{eq:first-diagonal-ogf-recurrence} that
\begin{equation}
\label{eq:first-diagonal-second-eulerian-ogf}
 D_p(t)=\frac{Q_p(t)}{(1-t)^{2p+1}}
 =\frac{t^{2p}P_p(t^{-1})}{(1-t)^{2p+1}}.
\end{equation}
Since $P_p(u)=\sum_{j=0}^{p}\eulertwo pj u^j$, coefficient extraction gives
\[
 [t^m]D_p(t)
 =\sum_{j=0}^{p}\eulertwo pj
   [t^{m-2p+j}](1-t)^{-2p-1}
 =\sum_{j=0}^{p}\eulertwo pj\binom{m+j}{2p},
\]
which is \eqref{eq:second-eulerian-first-diagonal}.  Both sides are polynomials in
$m$ of degree at most $2p$, so the same identity gives the stated polynomial
continuation.
\end{proof}
\part{Differentiation, Composition, and Reversion}

\chapter{Fa\`a di Bruno's formula}
\index{Faà di Bruno formula@Fa\`a di Bruno formula}

\section{The set-partition form}
For a block $B\subseteq[n]$, write $|B|$ for its size.
\begin{theorem}[Fa\`a di Bruno, partition form]\label{thm:faa-partition}
For $n\ge0$ and sufficiently differentiable $f,g$,
\begin{equation}\label{eq:faa-partition}
 D^n(f\circ g)(x)=
 \sum_{\pi\in\setpart([n])}
 f^{(|\pi|)}(g(x))\prod_{B\in\pi}g^{(|B|)}(x).
\end{equation}
For $n=0$, the empty partition contributes $f(g(x))$.
\end{theorem}
\begin{proof}
Induct on $n$.  Differentiate a term indexed by a partition $\pi$ of $[n]$.
If the derivative hits $f^{(|\pi|)}(g)$, the chain rule contributes $g'$ and
creates the new singleton block $\{n+1\}$.  If it hits the factor associated with
a block $B$, it changes $g^{(|B|)}$ to $g^{(|B|+1)}$, which inserts $n+1$ into
$B$.  Every partition of $[n+1]$ arises uniquely in exactly one of these ways:
remove $n+1$, either deleting its singleton block or removing it from its
non-singleton block.  Hence all terms occur once.
\end{proof}

\section{Multiplicity and Bell-polynomial forms}
Grouping partitions by the number $m_j$ of blocks of size $j$ gives the classical
formula.
\begin{theorem}[Fa\`a di Bruno, multiplicity form]\label{thm:faa-multiplicity}
\begin{multline}\label{eq:faa-multiplicity}
 D^n f(g(x))=
 \sum_{\substack{m_1,\ldots,m_n\ge0\\
                  m_1+2m_2+\cdots+nm_n=n}}
 \frac{n!}{m_1!m_2!\cdots m_n!}\\
 {}\times f^{(m_1+\cdots+m_n)}(g(x))
 \prod_{j=1}^{n}\left(\frac{g^{(j)}(x)}{j!}\right)^{m_j}.
\end{multline}
Equivalently,
\begin{equation}\label{eq:faa-bell}
 D^n f(g(x))=
 \sum_{k=0}^{n}f^{(k)}(g(x))
 \BellP nk\bigl(g'(x),g''(x),\ldots,g^{(n-k+1)}(x)\bigr).
\end{equation}
\end{theorem}
\begin{proof}
For fixed multiplicities $(m_j)$, the number of set partitions of $[n]$ with
$m_j$ blocks of size $j$ is
$n!/\prod_j((j!)^{m_j}m_j!)$.  Substitution in
\eqref{eq:faa-partition} gives \eqref{eq:faa-multiplicity}.  Grouping further by
$k=\sum m_j$ and invoking \eqref{eq:partial-bell-definition} gives
\eqref{eq:faa-bell}.
\end{proof}

For $n=4$ this reads
\begin{align}
 (f\circ g)^{(4)}={}&f^{(4)}(g)(g')^4
 +6f^{(3)}(g)(g')^2g''+3f''(g)(g'')^2\notag\\
 &+4f''(g)g'g^{(3)}+f'(g)g^{(4)}.
 \label{eq:faa-n4}
\end{align}
The coefficients $1,6,3,4,1$ count partitions of types
$1+1+1+1$, $2+1+1$, $2+2$, $3+1$, and $4$ respectively.  This also proves the
expanded example in the archive without four separate differentiations.

\section{The exponential special case}
Because every derivative of $e^u$ equals $e^u$, summing the partial Bell
polynomials in \eqref{eq:faa-bell} yields
\begin{equation}\label{eq:exp-faa}
 D^n e^{g(x)}=e^{g(x)}
 \BellC n\bigl(g'(x),g''(x),\ldots,g^{(n)}(x)\bigr).
\end{equation}
When $g$ is a cumulant generating function, this is precisely the
moment--cumulant relation of \cref{thm:moment-cumulant}.

\section{Mixed partial derivatives}
For $B\subseteq[n]$, write
$\partial_B=\prod_{j\in B}\partial/\partial x_j$.
\begin{theorem}[Multivariate Fa\`a di Bruno for distinct variables]
\label{thm:faa-multivariate}
If $y=g(x_1,\ldots,x_n)$, then
\begin{equation}\label{eq:faa-multivariate}
 \frac{\partial^n}{\partial x_1\cdots\partial x_n}f(y)
 =\sum_{\pi\in\setpart([n])}f^{(|\pi|)}(y)
   \prod_{B\in\pi}\partial_B y.
\end{equation}
\end{theorem}
\begin{proof}
The proof of \cref{thm:faa-partition} applies verbatim, except that the new
derivative is $\partial/\partial x_{n+1}$.  It either creates a singleton block or
joins the new index to one existing block.
\end{proof}

For three variables, \eqref{eq:faa-multivariate} becomes
\begin{align*}
 \partial_1\partial_2\partial_3f(y)
 ={}&f'(y)\partial_{123}y\\
 &+f''(y)\bigl((\partial_1y)(\partial_{23}y)
 +(\partial_2y)(\partial_{13}y)+(\partial_3y)(\partial_{12}y)\bigr)\\
 &+f^{(3)}(y)(\partial_1y)(\partial_2y)(\partial_3y).
\end{align*}
For repeated variables or vector-valued inner maps, one labels derivative slots
first and then contracts tensor indices; the same partition lattice remains the
combinatorial skeleton.

\chapter{Composition of ordinary and exponential series}

\section{Ordinary formal power series}
Let
\[
 F(u)=\sum_{k\ge0}a_ku^k,
 \qquad
 G(x)=\sum_{j\ge1}b_jx^j,
 \qquad
 F(G(x))=\sum_{n\ge0}c_nx^n.
\]
\begin{theorem}[Ordinary composition]\label{thm:ordinary-composition}
For $n\ge1$,
\begin{align}
 c_n&=\sum_{k=1}^{n}a_k
 \sum_{\substack{i_1+\cdots+i_k=n\\i_r\ge1}}
 b_{i_1}\cdots b_{i_k},
 \label{eq:ordinary-composition-compositions}\\
 &=\sum_{k=1}^{n}a_k
 \sum_{\substack{\pi_1+\cdots+\pi_n=k\\
                   \pi_1+2\pi_2+\cdots+n\pi_n=n}}
 \binom{k}{\pi_1,\ldots,\pi_n}
 b_1^{\pi_1}\cdots b_n^{\pi_n}
 \label{eq:ordinary-composition-multiplicities}\\
 &=\sum_{k=1}^{n}a_k\OBellP nk(b_1,\ldots,b_{n-k+1}),
 \label{eq:ordinary-composition-bell}
\end{align}
and $c_0=a_0$.
\end{theorem}
\begin{proof}
The coefficient of $x^n$ in $G(x)^k$ is obtained by choosing an ordered
composition $i_1+\cdots+i_k=n$, proving the first formula.  Collect equal part
sizes; the number of orderings with multiplicities $\pi_j$ is the multinomial
coefficient, proving the second.  The third is the definition of the ordinary
Bell polynomial.
\end{proof}

Two useful specializations are immediate.  If $F(u)=e^u$ and the series are
exponentially normalized, one obtains the exponential formula.  If
$A(x)=1+\sum_{n\ge1}a_nx^n$, then
\begin{equation}\label{eq:reciprocal-ordinary-bell}
 [x^n]\frac1{A(x)}
 =\sum_{k=1}^{n}(-1)^k\OBellP nk(a_1,\ldots,a_{n-k+1})
 \qquad(n\ge1),
\end{equation}
by taking $F(u)=1/(1-u)$ and $G(x)=-\sum_{n\ge1}a_nx^n$.

\section{Exponential formal power series}
Let
\[
 F(x)=\sum_{j\ge1}a_j\frac{x^j}{j!},\qquad
 G(u)=\sum_{k\ge0}b_k\frac{u^k}{k!},\qquad
 G(F(x))=\sum_{n\ge0}c_n\frac{x^n}{n!}.
\]
\begin{theorem}[Exponential composition]\label{thm:exponential-composition}
\begin{equation}\label{eq:exponential-composition}
 c_n=\sum_{k=0}^{n}b_k\BellP nk(a_1,\ldots,a_{n-k+1}).
\end{equation}
Equivalently,
\begin{equation}\label{eq:exponential-composition-partitions}
 c_n=\sum_{\pi\in\setpart([n])}b_{|\pi|}\prod_{B\in\pi}a_{|B|}.
\end{equation}
\end{theorem}
\begin{proof}
Substitute $F(x)$ into $G$ and apply
\eqref{eq:partial-bell-egf}.  The set-partition form is
\cref{thm:bell-poly-partitions} with an additional weight $b_k$ for the number of
blocks.
\end{proof}

Taking $b_k=1$ gives the complete Bell-polynomial exponential identity
\begin{equation}\label{eq:exp-complete-bell-again}
 \exp\!\left(\sum_{j\ge1}a_j\frac{x^j}{j!}\right)
 =\sum_{n\ge0}\BellC n(a_1,\ldots,a_n)\frac{x^n}{n!}.
\end{equation}
If the series converge, \eqref{eq:exponential-composition} is exactly
\eqref{eq:faa-bell} evaluated at the expansion point; formally, no convergence
notion is needed.

\chapter{Lagrange inversion and Lagrange--Bürmann}
\index{Lagrange inversion}

\section{Formal theorem}
Let $R$ be a commutative $\Q$-algebra and let $\phi(w)\in R[[w]]$ have invertible
constant term.  There is a unique $g(z)\in zR[[z]]$ satisfying
\begin{equation}\label{eq:lagrange-functional}
 g(z)=z\phi(g(z)).
\end{equation}
Uniqueness follows recursively because the coefficient of $z^n$ on the right
depends only on coefficients of $g$ below degree $n$.

\begin{theorem}[Lagrange--Bürmann]\label{thm:lagrange-burmann}
For $n\ge1$ and every formal series $H$,
\begin{align}
 \coeff zn H(g(z))
 &=\frac1n\coeff w{n-1}H'(w)\phi(w)^n,
 \label{eq:lagrange-burmann}\\
 &=\coeff wn H(w)\phi(w)^{n-1}\bigl(\phi(w)-w\phi'(w)\bigr).
 \label{eq:lagrange-burmann-alt}
\end{align}
In particular,
\begin{equation}\label{eq:lagrange-basic}
 \coeff zn g(z)=\frac1n\coeff w{n-1}\phi(w)^n.
\end{equation}
\end{theorem}
\begin{proof}
Put $f(w)=w/\phi(w)$, so that $f(g(z))=z$.  Formal residues and the change of
variables theorem \cref{thm:res-subst} give
\begin{align*}
 \coeff zn H(g(z))
 &=\res_z H(g(z))z^{-n-1}\dd z\\
 &=\res_w H(w)f(w)^{-n-1}f'(w)\dd w\\
 &=-\frac1n\res_w H(w)(f(w)^{-n})'\dd w\\
 &=\frac1n\res_w H'(w)f(w)^{-n}\dd w,
\end{align*}
where the residue of a derivative is zero.  Since
$f(w)^{-n}=w^{-n}\phi(w)^n$, this is
\eqref{eq:lagrange-burmann}.  Applying the same zero-residue rule to
$H(w)\phi(w)^nw^{-n}$ gives
\[
 \frac1n\res H'\phi^nw^{-n}
 =\res H\phi^nw^{-n-1}-\res H\phi^{n-1}\phi'w^{-n},
\]
which is \eqref{eq:lagrange-burmann-alt}.  Take $H(w)=w$ for
\eqref{eq:lagrange-basic}.
\end{proof}

\section{Analytic inverse functions}
\begin{theorem}[Analytic Lagrange inversion]\label{thm:analytic-lagrange}
Let $f$ be analytic near $a$, let $b=f(a)$, and suppose $f'(a)\ne0$.  Its local
inverse $g$ is analytic near $b$ and has
\begin{equation}\label{eq:analytic-lagrange-series}
 g(z)=a+\sum_{n\ge1}g_n\frac{(z-b)^n}{n!},
\end{equation}
where
\begin{equation}\label{eq:analytic-lagrange-coeff}
 g_n=\left.\frac{\dd^{n-1}}{\dd w^{n-1}}
 \left(\frac{w-a}{f(w)-f(a)}\right)^n\right|_{w=a}.
\end{equation}
The series has a positive radius of convergence.
\end{theorem}
\begin{proof}
The analytic inverse-function theorem gives an analytic local inverse and therefore
a convergent Taylor series.  Put $u=w-a$, $\zeta=z-b$, and
\[
 \phi(u)=\frac{u}{f(a+u)-f(a)};
\]
its constant term is $1/f'(a)$.  The inverse displacement satisfies
$u=\zeta\phi(u)$.  Formula \eqref{eq:lagrange-basic}, multiplied by $n!$, is
exactly \eqref{eq:analytic-lagrange-coeff}.
\end{proof}

If $f'(a)=0$ but $f(w)-f(a)=c(w-a)^q+O((w-a)^{q+1})$ with $c\ne0$, factor
$f(w)-f(a)=(w-a)^qU(w)$ with $U(a)=c$.  Choosing one of the $q$ values of
$c^{-1/q}$ reduces the equation to ordinary inversion in the variable
$(z-f(a))^{1/q}$.  Thus the inverse has $q$ local Puiseux branches.  This is the
precise sense in which Lagrange inversion extends to a critical point; a
single-valued Taylor inverse cannot exist there.

\section{Inverse coefficients in Bell-polynomial form}
Suppose
\[
 f(w)=\sum_{j\ge1}f_j\frac{w^j}{j!},\qquad f_1\ne0,
 \qquad
 g(z)=f^{\langle-1\rangle}(z)=\sum_{n\ge1}g_n\frac{z^n}{n!}.
\]
Define
\begin{equation}\label{eq:fhat}
 \widehat f_j=\frac{f_{j+1}}{(j+1)f_1}.
\end{equation}
\begin{theorem}[Explicit inverse coefficients]\label{thm:inverse-bell-coeff}
One has $g_1=f_1^{-1}$ and, for $n\ge2$,
\begin{equation}\label{eq:inverse-bell-coeff}
 g_n=\frac1{f_1^n}\sum_{k=1}^{n-1}(-1)^k\rising nk
 \BellP{n-1}{k}(\widehat f_1,\ldots,\widehat f_{n-k}).
\end{equation}
\end{theorem}
\begin{proof}
Write
\[
 \frac{f(w)}{f_1w}=1+U(w),
 \qquad U(w)=\sum_{j\ge1}\widehat f_j\frac{w^j}{j!}.
\]
By \eqref{eq:lagrange-basic},
$g_n=(n-1)!\coeff w{n-1}(w/f(w))^n$.  Now
\[
 (1+U)^{-n}=\sum_{k\ge0}(-1)^k\binom{n+k-1}{k}U^k
 =\sum_{k\ge0}(-1)^k\frac{\rising nk}{k!}U^k.
\]
Apply \eqref{eq:partial-bell-egf} and extract degree $n-1$; the $k=0$ term
vanishes when $n\ge2$.
\end{proof}

\section{The associahedron interpretation, with normalized notation}
The formula in the archive mixes exponential coefficients with an ordinary-series
face formula.  The exact geometric statement is the following.

Let
\begin{equation}\label{eq:ordinary-series-associa}
 A(x)=x+\sum_{n\ge1}a_nx^{n+1},\qquad
 B(x)=A^{\langle-1\rangle}(x)=x+\sum_{n\ge1}b_nx^{n+1}.
\end{equation}
Use the indexing in which the associahedron $K_n$ has dimension $n-1$ and its
faces are plane rooted trees with $n+1$ leaves and internal arities at least $2$.
If an internal vertex has arity $i+1$, give it weight $a_i$; let $a_F$ be the
product of all vertex weights.  Equivalently, if
$F\cong K_{i_1}\times\cdots\times K_{i_m}$, then
$a_F=a_{i_1}\cdots a_{i_m}$.

\begin{theorem}[Associahedral inversion]\label{thm:associahedral-inversion}
\begin{equation}\label{eq:associahedral-inversion}
 b_n=\sum_{F\text{ a face of }K_n}(-1)^{n-\dim F}a_F.
\end{equation}
\end{theorem}
\begin{proof}
The equation $A(B)=x$ is
\[
 B=x-\sum_{i\ge1}a_iB^{i+1}.
\]
Iterating it recursively expands $b_n$ as a sum over plane rooted trees with
$n+1$ leaves: every internal vertex of arity $i+1$ contributes $-a_i$.
A face $F$ of $K_n$ is exactly such a tree after contracting a chosen set of
internal edges.  If it has $m$ internal vertices, then
$\dim F=n-m$, so its sign is $(-1)^m=(-1)^{n-\dim F}$ and its weight is $a_F$.
Thus the tree expansion is the alternating face sum.
\end{proof}

For example,
\begin{align*}
 b_1&=-a_1,\\
 b_2&=-a_2+2a_1^2,\\
 b_3&=-a_3+5a_2a_1-5a_1^3,\\
 b_4&=-a_4+6a_3a_1+3a_2^2-21a_2a_1^2+14a_1^4.
\end{align*}
For $K_4$, the coefficients $1,6,3,21,14$ count respectively the whole
3-polytope, its six pentagonal and three square facets, its edges, and its
vertices.

\chapter{Algebraic and transcendental examples}

\section{The Fuss--Catalan algebraic equation}
Let $p\ge2$ be an integer and consider
\begin{equation}\label{eq:fuss-equation}
 x^p-x+z=0,
 \qquad x(0)=0.
\end{equation}
It is equivalent to
$x=z\phi(x)$ with $\phi(x)=(1-x^{p-1})^{-1}$.

\begin{theorem}\label{thm:fuss-series}
The distinguished solution is
\begin{equation}\label{eq:fuss-series}
 x(z)=\sum_{k=0}^{\infty}\binom{pk}{k}
       \frac{z^{(p-1)k+1}}{(p-1)k+1}.
\end{equation}
Its radius of convergence is
\begin{equation}\label{eq:fuss-radius}
 R_p=(p-1)p^{-p/(p-1)},
\end{equation}
and the series converges absolutely also on $|z|=R_p$.
\end{theorem}
\begin{proof}
By \eqref{eq:lagrange-basic},
\[
 \coeff znx(z)=\frac1n\coeff w{n-1}(1-w^{p-1})^{-n}.
\]
This vanishes unless $n-1=(p-1)k$.  In that case the coefficient is
$\frac1n\binom{n+k-1}{k}=\binom{pk}{k}/((p-1)k+1)$, proving the series.

The inverse map is $z=f(x)=x-x^p$.  Its nearest critical points satisfy
$f'(x)=1-px^{p-1}=0$ and all have modulus $p^{-1/(p-1)}$; their critical values
have modulus $(p-1)p^{-p/(p-1)}$.  No local inverse can pass analytically through
a critical value, while the inverse-function theorem continues the branch up to
one, so this is the radius.  Stirling's formula gives coefficients of order
$Ck^{-3/2}R_p^{-((p-1)k+1)}$, which proves absolute convergence on the boundary.
\end{proof}

\section{Lambert's \texorpdfstring{$W$}{W} function at the origin}
The Lambert function satisfies $W(z)e^{W(z)}=z$.
\begin{theorem}\label{thm:lambert-W-zero}
On its principal branch near $0$,
\begin{equation}\label{eq:lambert-W-zero}
 W(z)=\sum_{n=1}^{\infty}(-n)^{n-1}\frac{z^n}{n!}
 =z-z^2+\frac32z^3-\frac83z^4+O(z^5),
\end{equation}
with radius of convergence $e^{-1}$.
\end{theorem}
\begin{proof}
Here $f(w)=we^w$, so $w/f(w)=e^{-w}$.  Formula
\eqref{eq:analytic-lagrange-coeff} gives
$g_n=D^{n-1}e^{-nw}|_{w=0}=(-n)^{n-1}$.  The inverse map has its nearest critical
point at $w=-1$, where $f(-1)=-e^{-1}$; this square-root branch point determines
the radius.
\end{proof}

\section{A rapidly useful expansion around \texorpdfstring{$W(e)=1$}{W(e)=1}}
Put $\xi=\log z-1$ and $u=W(z)-1$.  Then
\begin{equation}\label{eq:lambert-near-e-implicit}
 u+\log(1+u)=\xi.
\end{equation}
Reversion gives
\begin{equation}\label{eq:lambert-near-e}
 W(e^{1+\xi})=1+\frac{\xi}{2}+\frac{\xi^2}{16}
 -\frac{\xi^3}{192}-\frac{\xi^4}{3072}
 +\frac{13\xi^5}{61440}+O(\xi^6).
\end{equation}
\begin{proof}
The derivative at $u=0$ of the left side of
\eqref{eq:lambert-near-e-implicit} is $2$, so it has an ordinary inverse.
Substitute $u=c_1\xi+\cdots+c_5\xi^5$ into
$u+\log(1+u)=\xi$ and compare successive powers.  The triangular equations give
$c_1=1/2,c_2=1/16,c_3=-1/192,c_4=-1/3072,c_5=13/61440$.
\end{proof}
The closest critical points of $u+\log(1+u)$ satisfy
$1+(1+u)^{-1}=0$, hence $u=-2$; using the two adjacent logarithm values gives
images $-2\pm i\pi$.  Therefore the radius in the $\xi$-plane is
$\sqrt{4+\pi^2}$.  For positive real $z$, the expansion converges whenever
$|\log z-1|<\sqrt{4+\pi^2}$, approximately
$0.0655<z<112.63$.  Taking $\xi=-1$ yields a convergent evaluation of
$W(1)=0.567143\ldots$.

\section{Binary trees and Catalan numbers}
Let $C_n$ count rooted ordered full binary trees with $n$ internal vertices, and
let $B(z)=\sum_{n\ge0}C_nz^n$.  Removing the root gives
\begin{equation}\label{eq:binary-tree-functional}
 B(z)=1+zB(z)^2.
\end{equation}
Put $C(z)=B(z)-1$; then $C=z(1+C)^2$.  Lagrange inversion gives
\begin{equation}\label{eq:catalan-lagrange}
 C_n=\coeff znC(z)=\frac1n\coeff w{n-1}(1+w)^{2n}
 =\frac1n\binom{2n}{n-1}=\frac1{n+1}\binom{2n}{n}.
\end{equation}
This simultaneously proves the functional equation, its unique formal solution,
and the Catalan closed form.  The symbol $C_n$ is used here to avoid confusing
these binary-tree numbers with Bell numbers.

\chapter{Lagrange's reversion theorem in shifted form}

\section{The two-parameter implicit equation}
Let $v=v(x,y)$ be the unique formal solution of
\begin{equation}\label{eq:shifted-reversion-equation}
 v=x+yf(v),
 \qquad v(x,0)=x.
\end{equation}

\begin{theorem}[Lagrange reversion]\label{thm:shifted-lagrange-reversion}
For every formal or analytic $g$,
\begin{equation}\label{eq:shifted-lagrange-reversion}
 g(v)=g(x)+\sum_{k=1}^{\infty}\frac{y^k}{k!}
 \left(\frac{\partial}{\partial x}\right)^{k-1}
 \bigl(f(x)^kg'(x)\bigr).
\end{equation}
In particular,
\begin{equation}\label{eq:shifted-lagrange-v}
 v=x+\sum_{k=1}^{\infty}\frac{y^k}{k!}
 \left(\frac{\partial}{\partial x}\right)^{k-1}f(x)^k.
\end{equation}
\end{theorem}
\begin{proof}
Put $u=v-x$.  Then $u=y\phi(u)$ with $\phi(u)=f(x+u)$, treating $x$ as a
coefficient parameter.  Apply \eqref{eq:lagrange-burmann} to
$H(u)=g(x+u)$:
\[
 [y^k]g(v)=\frac1k[u^{k-1}]g'(x+u)f(x+u)^k
 =\frac1{k!}D_x^{k-1}(g'(x)f(x)^k).
\]
Summation proves the first formula; take $g(z)=z$ for the second.
\end{proof}

\section{The delta-function proof decoded}
The archive also gives a distributional derivation.  Assume temporarily that
$x,y\in\R$, the equation $z-x-yf(z)=0$ has one simple root $v$, and
$1-yf'(v)>0$.  The elementary change-of-variables rule for the Dirac delta gives
\begin{equation}\label{eq:delta-root}
 g(v)=\int_{\R}\delta(yf(z)-z+x)g(z)(1-yf'(z))\,\dd z.
\end{equation}
Using
$\delta(s)=\int_{\R}e^{iks}\dd k/(2\pi)$ and expanding only in the formal
parameter $y$,
\begin{align}
 g(v)
 &=\sum_{n\ge0}D_x^n\left[
 \frac{y^nf(x)^ng(x)}{n!}(1-yf'(x))\right]\notag\\
 &=\sum_{n\ge0}D_x^n\left[
 \frac{y^nf(x)^ng(x)}{n!}
 -\frac{y^{n+1}}{(n+1)!}\bigl((gf^{n+1})'-g'f^{n+1}\bigr)
 \right].
 \label{eq:delta-reversion-middle}
\end{align}
The first term of level $n+1$ cancels the total derivative in the second term of
level $n$; what remains is exactly
\eqref{eq:shifted-lagrange-reversion}.  Thus the Fourier--delta calculation is a
compact analytic avatar of the formal residue proof.  Formula
\eqref{eq:delta-root} requires an absolute Jacobian in general; the displayed sign
is valid in the stated orientation-preserving neighbourhood.

\chapter{Laplace-type asymptotic integrals}
\index{Laplace asymptotics}

\section{Endpoint expansion}
Consider
\begin{equation}\label{eq:laplace-integral}
 I(\lambda)=\int_a^b e^{-\lambda f(x)}g(x)\,\dd x,
 \qquad \lambda\to+\infty,
\end{equation}
where the endpoint $a$ is the unique relevant minimum and, with $t=x-a$,
\begin{align}
 f(a+t)&\sim f(a)+\sum_{k\ge0}a_kt^{k+\alpha},
 &a_0>0,\quad\alpha>0,
 \label{eq:laplace-f-expansion}\\
 g(a+t)&\sim\sum_{k\ge0}b_kt^{k+\beta-1},
 &\beta>0.
 \label{eq:laplace-g-expansion}
\end{align}
Let $\OBellP{k}{j}$ denote the ordinary Bell polynomial.

\begin{theorem}[Laplace--Erd\'elyi coefficient formula]
\label{thm:laplace-bell}
Under the standard differentiability and remainder hypotheses that permit Watson's
lemma,
\begin{equation}\label{eq:laplace-asymptotic}
 I(\lambda)\sim e^{-\lambda f(a)}
 \sum_{n=0}^{\infty}
 \Gamma\!\left(\frac{n+\beta}{\alpha}\right)
 \frac{c_n}{\lambda^{(n+\beta)/\alpha}},
\end{equation}
where
\begin{equation}\label{eq:laplace-cn}
 c_n=\frac1{\alpha a_0^{(n+\beta)/\alpha}}
 \sum_{k=0}^{n}b_{n-k}
 \sum_{j=0}^{k}
 \binom{-\frac{n+\beta}{\alpha}}{j}
 \frac{\OBellP{k}{j}(a_1,\ldots,a_{k-j+1})}{a_0^j}.
\end{equation}
\end{theorem}
\begin{proof}
Write
$f(a+t)-f(a)=t^\alpha A(t)$ with
$A(t)=a_0+a_1t+\cdots$, and set
$s=(f(a+t)-f(a))^{1/\alpha}=tA(t)^{1/\alpha}$.  This has an inverse $t=t(s)$.
Expand
\[
 g(a+t(s))t'(s)=\sum_{n\ge0}d_ns^{n+\beta-1}.
\]
A generalized Lagrange--Bürmann calculation, first for integral $\beta$ and then
by formal continuation in the exponent, gives
\[
 d_n=\sum_{k=0}^{n}b_{n-k}[t^k]A(t)^{-(n+\beta)/\alpha}.
\]
Indeed, apply \eqref{eq:lagrange-burmann} to an antiderivative of
$t^{m+\beta-1}$ and differentiate back after substitution.

Put $\rho=(n+\beta)/\alpha$.  The ordinary Bell expansion yields
\begin{align*}
 [t^k]A(t)^{-\rho}
 &=a_0^{-\rho}[t^k]
 \sum_{j\ge0}\binom{-\rho}{j}
 \left(\sum_{r\ge1}\frac{a_r}{a_0}t^r\right)^j\\
 &=a_0^{-\rho}\sum_{j=0}^{k}\binom{-\rho}{j}a_0^{-j}
 \OBellP{k}{j}(a_1,\ldots,a_{k-j+1}).
\end{align*}
Thus $c_n=d_n/\alpha$ is \eqref{eq:laplace-cn}.  Finally
\[
 \int_0^\infty e^{-\lambda s^\alpha}s^{n+\beta-1}\dd s
 =\frac1\alpha\Gamma\!\left(\frac{n+\beta}{\alpha}\right)
 \lambda^{-(n+\beta)/\alpha},
\]
and Watson's lemma controls the replacement of the finite endpoint interval by
the asymptotic expansion.
\end{proof}

\begin{remark}
The appearance of ordinary rather than exponential Bell polynomials is forced by
the ordinary Taylor expansion of $A(t)^{-\rho}$.  Lagrange inversion enters in the
change from $t$ to the natural Laplace coordinate $s$; the gamma factor enters only
in the final monomial integral.
\end{remark}

\part{Higher Differential Calculus and Multivariate Reversion}

\chapter{Products, quotients, and coordinate-free composition}

The scalar Bell-polynomial form of Fa\`a di Bruno's formula is only one face of a
larger calculus. Products of many factors, vector-valued maps, repeated
composition, and inverse functions all become transparent when derivatives are
regarded as symmetric multilinear maps. This chapter integrates complementary
material from the five drafts and supplies proofs that were only sketched there.

\section{The multinomial Leibniz rule and higher quotient rules}

\begin{theorem}[Multinomial Leibniz rule]\label{unif:thm-multinomial-leibniz}
Let $f_1,\ldots,f_q$ be $n$ times differentiable. Then
\begin{equation}\label{unif:eq-multinomial-leibniz}
 D^n\prod_{r=1}^{q}f_r(x)
 =\sum_{j_1+\cdots+j_q=n}
   \binom{n}{j_1,\ldots,j_q}\prod_{r=1}^{q}f_r^{(j_r)}(x).
\end{equation}
\end{theorem}
\begin{proof}
Taylor-expand every factor at $x$:
\[
 f_r(x+h)=\sum_{j_r=0}^{n}f_r^{(j_r)}(x)\frac{h^{j_r}}{j_r!}+O(h^{n+1}).
\]
The coefficient of $h^n$ in the product is
\[
 \sum_{j_1+\cdots+j_q=n}\prod_{r=1}^{q}
 \frac{f_r^{(j_r)}(x)}{j_r!}.
\]
The same coefficient in the Taylor expansion of the product itself is
$n!^{-1}D^n\prod_r f_r(x)$. Multiplication by $n!$ proves the formula. The
argument is purely coefficient-theoretic, so it is valid for formal series as
well as analytic functions.
\end{proof}

\begin{proposition}[Reciprocal and quotient derivatives]
\label{unif:prop-reciprocal-derivative}
If $g(x)\ne0$, then
\begin{equation}\label{unif:eq-reciprocal-derivative}
 D^r\!\left(\frac1g\right)
 =\sum_{k=0}^{r}(-1)^k k!\,
 \frac{\BellP{r}{k}(g',g'',\ldots,g^{(r-k+1)})}{g^{k+1}},
\end{equation}
where $\BellP00=1$. Consequently,
\begin{equation}\label{unif:eq-higher-quotient}
 D^n\!\left(\frac fg\right)
 =\sum_{r=0}^{n}\binom nr f^{(n-r)}
 \sum_{k=0}^{r}(-1)^k k!\,
 \frac{\BellP{r}{k}(g',g'',\ldots)}{g^{k+1}}.
\end{equation}
\end{proposition}
\begin{proof}
Apply Fa\`a di Bruno's formula to the composition of $g$ with
$u\mapsto u^{-1}$, whose $k$th derivative is
$(-1)^k k!u^{-k-1}$. This gives
\eqref{unif:eq-reciprocal-derivative}. Apply the ordinary Leibniz rule to
$f\,g^{-1}$ and substitute the reciprocal formula to obtain
\eqref{unif:eq-higher-quotient}.
\end{proof}

\section{A coordinate-free Fa\`a di Bruno formula}

Let $E,F,G$ be Banach spaces. Write $D^rg(x)$ for the symmetric $r$-linear
Fr\'echet derivative. If $B\subseteq[n]$, abbreviate
\[
 D^{|B|}g(x)[h_i:i\in B]
\]
for the derivative evaluated on the directions indexed by $B$.

\begin{theorem}[Partition form of Fr\'echet Fa\`a di Bruno]
\label{unif:thm-frechet-faa}
For $n$ times continuously Fr\'echet differentiable maps $g:E\to F$ and
$f:F\to G$,
\begin{equation}\label{unif:eq-frechet-faa}
 D^n(f\circ g)(x)[h_1,\ldots,h_n]
 =\sum_{\pi\in\Pi_n}D^{|\pi|}f(g(x))
 \Bigl[D^{|B|}g(x)[h_i:i\in B]\Bigr]_{B\in\pi}.
\end{equation}
\end{theorem}
\begin{proof}
Induct on $n$. The case $n=1$ is the chain rule. Differentiate a term indexed
by a partition $\pi$ of $[n]$ in a new direction $h_{n+1}$. Differentiating
the outer derivative creates a new singleton block $\{n+1\}$. Differentiating
the inner derivative associated with $B\in\pi$ adjoins $n+1$ to that block.
Every partition of $[n+1]$ arises uniquely in one of these ways: remove
$n+1$ and note whether it was a singleton. Hence every term in the desired
sum appears once, completing the induction.
\end{proof}

In one variable, grouping partitions by block-size profile recovers the usual
Bell-polynomial form. For three scalar functions one may write
\begin{align}
 D^n f(g(h(x)))
 &=\sum_{k=0}^{n}f^{(k)}(g(h(x)))
 \BellP{n}{k}\bigl((g\circ h)',(g\circ h)'',\ldots\bigr),
 \label{unif:eq-three-composition-outer}\\
 (g\circ h)^{(j)}
 &=\sum_{\ell=0}^{j}g^{(\ell)}(h(x))
 \BellP{j}{\ell}(h',h'',\ldots).
 \label{unif:eq-three-composition-inner}
\end{align}
The nested partitions show directly why both parenthesizations of a multiple
composition yield the same coefficient: each describes a partition of $[n]$
into inner blocks together with a partition of the set of those blocks.

\section{Recursive derivatives of an inverse function}

\begin{theorem}[Inverse-derivative recursion]
\label{unif:thm-inverse-derivative-recursion}
Let $y=f(x)$, assume $f'(x)\ne0$, and put $x=g(y)=f^{-1}(y)$. Then
$g'=1/f'(g)$ and, for $n\ge2$,
\begin{equation}\label{unif:eq-inverse-derivative-recursion}
 g^{(n)}=-\frac1{f'(g)}
 \sum_{k=2}^{n}f^{(k)}(g)
 \BellP{n}{k}(g',g'',\ldots,g^{(n-k+1)}).
\end{equation}
The first cases are
\begin{align}
 g'&=\frac1{f'},\\
 g''&=-\frac{f''}{(f')^3},\\
 g'''&=\frac{3(f'')^2-f'f'''}{(f')^5},\\
 g^{(4)}&=\frac{-15(f'')^3+10f'f''f'''-(f')^2f^{(4)}}{(f')^7},
\end{align}
where derivatives of $f$ are evaluated at $g(y)$.
\end{theorem}
\begin{proof}
Differentiate $f(g(y))=y$. For $n\ge2$, Fa\`a di Bruno gives
\[
 0=\sum_{k=1}^{n}f^{(k)}(g)
 \BellP{n}{k}(g',g'',\ldots,g^{(n-k+1)}).
\]
The term $k=1$ is $f'(g)g^{(n)}$; isolating it proves the recursion. The four
low-order formulas follow by substituting the corresponding partial Bell
polynomials and simplifying.
\end{proof}

\chapter{Good's multivariate inversion theorem}

The univariate Lagrange formula controls one implicit series. Coupled systems
require the determinant correction discovered by Good.

\section{Formal logarithmic residues}

For $d$ variables, write
$\operatorname{CT}_{\bm x}H=[x_1^0\cdots x_d^0]H$.

\begin{lemma}[Formal logarithmic change of variables]
\label{unif:lem-formal-log-change}
Suppose $t_i=x_i u_i(\bm x)$, where each $u_i$ is a unit. Then
\begin{equation}\label{unif:eq-formal-log-change}
 \operatorname{CT}_{\bm t}H(\bm t)
 =\operatorname{CT}_{\bm x}H(\bm t(\bm x))
 \det\!\left(\frac{\partial\log t_i}{\partial\log x_j}\right)_{i,j=1}^{d}.
\end{equation}
\end{lemma}
\begin{proof}
By linearity it is enough to test a Laurent monomial $H(\bm t)=\bm t^{\bm m}$.
For $\bm m=\bm0$, the determinant is $1$ plus logarithmic derivatives and has
constant term $1$. If $\bm m\ne\bm0$, choose $r$ with $m_r\ne0$. Expanding the
determinant and using
\[
 t_r^{m_r}\frac{\partial\log t_r}{\partial\log x_j}
 =\frac1{m_r}x_j\frac{\partial}{\partial x_j}t_r^{m_r}
\]
expresses the transformed integrand as a sum of total logarithmic derivatives.
The constant term of $x_j\partial K/\partial x_j$ is zero. Thus both sides
vanish, proving the substitution rule.
\end{proof}

\begin{theorem}[Good's multivariate Lagrange inversion]
\label{unif:thm-good-inversion}
Let $\phi_i(\bm x)$ have invertible constant term and let $\bm w(\bm t)$ be
the unique formal solution of
\begin{equation}\label{unif:eq-good-system}
 w_i=t_i\phi_i(\bm w),\qquad1\le i\le d.
\end{equation}
Then
\begin{equation}\label{unif:eq-good-inversion}
 [\bm t^{\bm n}]F(\bm w(\bm t))
 =[\bm x^{\bm n}]F(\bm x)\prod_{i=1}^{d}\phi_i(\bm x)^{n_i}
 \det\!\left(
 \delta_{ij}-\frac{x_j}{\phi_i(\bm x)}
 \frac{\partial\phi_i}{\partial x_j}(\bm x)
 \right)_{i,j=1}^{d}.
\end{equation}
\end{theorem}
\begin{proof}
Write the coefficient as
\[
 \operatorname{CT}_{\bm t}F(\bm w(\bm t))\bm t^{-\bm n}.
\]
The inverse substitution is $t_i=x_i/\phi_i(\bm x)$. Apply
\cref{unif:lem-formal-log-change}. Since
\[
 \frac{\partial\log t_i}{\partial\log x_j}
 =\delta_{ij}-\frac{x_j}{\phi_i}\frac{\partial\phi_i}{\partial x_j}
\]
and $\bm t^{-\bm n}=\bm x^{-\bm n}\prod_i\phi_i^{n_i}$, the transformed
constant term is precisely the right side of
\eqref{unif:eq-good-inversion}.
\end{proof}

\begin{example}[A coupled two-series system]
Let $u=x(1+v)$ and $v=y(1+u)$. Good's determinant is
$1-UV/((1+U)(1+V))$, so
\begin{align}
 [x^my^n]u^pv^q
 &=\binom{n}{m-p}\binom{m}{n-q}
 \notag\\
 &\quad-\binom{n-1}{m-p-1}\binom{m-1}{n-q-1},
 \label{unif:eq-coupled-coefficients}
\end{align}
with the usual zero convention for out-of-range binomial coefficients.
Indeed, the theorem asks for the coefficient of $U^mV^n$ in
\[
 U^pV^q(1+V)^m(1+U)^n
 -U^{p+1}V^{q+1}(1+V)^{m-1}(1+U)^{n-1},
\]
which gives the two products of binomial coefficients.
\end{example}

\chapter{Partition lattices, Sheffer families, and Riordan arrays}

The five drafts repeatedly move among moments and cumulants, polynomial
sequences of binomial type, Bell-polynomial matrices, and compositional
inverses. Incidence algebras and exponential Riordan arrays place these
phenomena in one structure.

\section{M\"obius inversion on the partition lattice}

\begin{theorem}[M\"obius function of the partition lattice]
\label{unif:thm-partition-mobius}
For a partition $\pi\in\Pi_n$,
\begin{equation}\label{unif:eq-partition-mobius}
 \mu(\pi,\widehat1)=(-1)^{|\pi|-1}(|\pi|-1)!.
\end{equation}
\end{theorem}
\begin{proof}
The interval $[\pi,\widehat1]$ is isomorphic to the partition lattice on the
$|\pi|$ blocks of $\pi$, so let
$\mu_m=\mu_{\Pi_m}(\widehat0,\widehat1)$. Every interval
$[\widehat0,\sigma]$ is a product of partition lattices, hence
$\mu(\widehat0,\sigma)=\prod_{B\in\sigma}\mu_{|B|}$. The defining M\"obius
relation says that the sum of these products over $\Pi_m$ is $1$ for
$m=0,1$ and $0$ for $m\ge2$. If
$M(z)=\sum_{m\ge1}\mu_mz^m/m!$, the exponential formula gives
$\exp(M(z))=1+z$. Therefore $M(z)=\log(1+z)$ and
$\mu_m=(-1)^{m-1}(m-1)!$.
\end{proof}

\begin{corollary}[Moment--cumulant inversion]
If
\begin{equation}\label{unif:eq-moment-partition}
 m_n=\sum_{\pi\in\Pi_n}\prod_{B\in\pi}\kappa_{|B|},
\end{equation}
then
\begin{equation}\label{unif:eq-cumulant-partition}
 \kappa_n=\sum_{\pi\in\Pi_n}(-1)^{|\pi|-1}(|\pi|-1)!
 \prod_{B\in\pi}m_{|B|}.
\end{equation}
\end{corollary}
\begin{proof}
The first formula is the zeta transform in the incidence algebra of $\Pi_n$.
Applying the M\"obius function from
\cref{unif:thm-partition-mobius} gives the inverse. Equivalently, the EGFs
satisfy $M(z)=\exp(K(z))$ and hence $K(z)=\log M(z)$.
\end{proof}

\section{Sheffer sequences and delta operators}

\begin{definition}
A polynomial sequence $(s_n(x))_{n\ge0}$ is Sheffer when
\begin{equation}\label{unif:eq-sheffer-egf}
 \sum_{n\ge0}s_n(x)\frac{z^n}{n!}
 =A(z)e^{x\beta(z)},
 \qquad A(0)\ne0,\quad\beta(0)=0,\quad\beta'(0)\ne0.
\end{equation}
Its associated binomial-type sequence is defined by
\begin{equation}\label{unif:eq-associated-egf}
 \sum_{n\ge0}p_n(x)\frac{z^n}{n!}=e^{x\beta(z)}.
\end{equation}
\end{definition}

\begin{theorem}[Sheffer addition and lowering laws]
\label{unif:thm-sheffer-laws}
One has
\begin{equation}\label{unif:eq-sheffer-addition}
 s_n(x+y)=\sum_{j=0}^{n}\binom nj s_j(x)p_{n-j}(y).
\end{equation}
If $\bar\beta$ is the compositional inverse of $\beta$ and
$Q=\bar\beta(D_x)$, then
\begin{equation}\label{unif:eq-sheffer-lowering}
 Qs_n(x)=ns_{n-1}(x),\qquad Qp_n(x)=np_{n-1}(x).
\end{equation}
\end{theorem}
\begin{proof}
Factor the EGF at $x+y$ as
$A(z)e^{x\beta(z)}e^{y\beta(z)}$ and take the EGF Cauchy product to obtain
\eqref{unif:eq-sheffer-addition}. Since
$D_xe^{x\beta(z)}=\beta(z)e^{x\beta(z)}$, formal functional calculus on
polynomials gives
\[
 \bar\beta(D_x)e^{x\beta(z)}=\bar\beta(\beta(z))e^{x\beta(z)}
 =ze^{x\beta(z)}.
\]
The factor $A(z)$ is independent of $x$, so coefficient comparison proves the
lowering laws.
\end{proof}

If $A(z)=\sum a_jz^j/j!$, then
\begin{equation}\label{unif:eq-sheffer-convolution}
 s_n(x)=\sum_{j=0}^{n}\binom nj a_jp_{n-j}(x).
\end{equation}
Thus every Sheffer sequence is an invertible Appell prefactor convolved with a
binomial-type sequence, and Lagrange inversion computes its delta operator
through $\bar\beta$.

\section{Exponential Riordan arrays}

For $g(0)\ne0$ and $f(z)=f_1z+O(z^2)$ with $f_1\ne0$, define
\begin{equation}\label{unif:eq-riordan-entry}
 R_{n,k}=\frac{n!}{k!}[z^n]g(z)f(z)^k.
\end{equation}

\begin{theorem}[Riordan action, product, and inverse]
\label{unif:thm-riordan-group}
If $(a_k)$ has EGF $A(z)$, multiplication by the array $(g,f)$ produces the
EGF
\begin{equation}\label{unif:eq-riordan-action}
 g(z)A(f(z)).
\end{equation}
Consequently,
\begin{align}
 (g,f)(h,\ell)&=(g\cdot(h\circ f),\ell\circ f),
 \label{unif:eq-riordan-product}\\
 (g,f)^{-1}&=\left(\frac1{g\circ\bar f},\bar f\right),
 \label{unif:eq-riordan-inverse}
\end{align}
where $\bar f=f^{\langle-1\rangle}$.
\end{theorem}
\begin{proof}
By definition,
\[
 \sum_n\left(\sum_kR_{n,k}a_k\right)\frac{z^n}{n!}
 =g(z)\sum_k a_k\frac{f(z)^k}{k!}=g(z)A(f(z)).
\]
Applying two such operators in succession gives
$g(z)h(f(z))A(\ell(f(z)))$, proving the product. Choosing
$\ell=\bar f$ and $h=1/(g\circ\bar f)$ gives the identity pair $(1,z)$ and
hence the inverse.
\end{proof}

The Bell-polynomial matrix is the special case $(1,f)$:
\begin{equation}\label{unif:eq-bell-riordan}
 \frac{n!}{k!}[z^n]
 \left(\sum_{j\ge1}x_j\frac{z^j}{j!}\right)^k
 =\BellP{n}{k}(x_1,x_2,\ldots).
\end{equation}
The second-kind and signed first-kind Stirling matrices correspond to
$f(z)=e^z-1$ and $f(z)=\log(1+z)$. Since these series are compositional
inverses, Stirling matrix inversion is a Riordan-group identity.

\section{Abel polynomials from Lagrange inversion}

\begin{theorem}[Abel's binomial identity]
For $n\ge1$, put $p_n(x)=x(x-an)^{n-1}$ and $p_0(x)=1$. Then
\begin{align}
 &(x+y)(x+y-an)^{n-1}
 \notag\\
 &\quad=\sum_{k=0}^{n}\binom nk
 x(x-ak)^{k-1}y\bigl(y-a(n-k)\bigr)^{n-k-1}.
 \label{unif:eq-abel-binomial}
\end{align}
\end{theorem}
\begin{proof}
Let $T=z\exp(-aT)$. Lagrange--B\"urmann gives
\[
 [z^n]e^{xT(z)}=\frac{x}{n}[u^{n-1}]e^{(x-an)u}
 =\frac{x(x-an)^{n-1}}{n!}.
\]
Thus $e^{xT(z)}=\sum p_n(x)z^n/n!$. Multiplying the EGFs at $x$ and $y$ and
comparing coefficients proves the binomial identity.
\end{proof}

\chapter{A complete coefficient-algorithm toolkit}

Let $A(z)=\sum_{n\ge0}a_nz^n$ and $B(z)=\sum_{n\ge0}b_nz^n$. The following
triangular recurrences consolidate the computational sections of the five
drafts. Each proof comes from one coefficient comparison, so the formulas are
valid over any coefficient ring in which the displayed divisions make sense.

\begin{theorem}[Elementary formal-series algorithms]
\label{unif:thm-series-algorithms}
The following recurrences determine the indicated operations to arbitrary
order.

\paragraph{Product.} If $C=AB$, then
\begin{equation}\label{unif:eq-cauchy-product}
 c_n=\sum_{j=0}^{n}a_jb_{n-j}.
\end{equation}

\paragraph{Reciprocal and quotient.} If $a_0\ne0$, then for $C=1/A$,
\begin{equation}\label{unif:eq-reciprocal-recursion}
 c_0=a_0^{-1},\qquad
 c_n=-a_0^{-1}\sum_{j=1}^{n}a_jc_{n-j},
\end{equation}
and for $C=B/A$,
\begin{equation}\label{unif:eq-quotient-recursion}
 c_n=a_0^{-1}\left(b_n-\sum_{j=1}^{n}a_jc_{n-j}\right).
\end{equation}

\paragraph{Logarithm and exponential.} If $a_0=1$ and
$L=\log A=\sum_{n\ge1}\ell_nz^n$, then
\begin{equation}\label{unif:eq-log-recursion}
 \ell_n=a_n-\frac1n\sum_{j=1}^{n-1}j\ell_j a_{n-j}.
\end{equation}
Conversely, if $A=e^L$, then
\begin{equation}\label{unif:eq-exp-recursion}
 a_0=1,\qquad a_n=\frac1n\sum_{j=1}^{n}j\ell_j a_{n-j}.
\end{equation}

\paragraph{Arbitrary powers.} If $a_0=1$ and $C=A^\alpha$, then
\begin{equation}\label{unif:eq-power-recursion}
 c_0=1,
 \qquad c_n=\frac1n\sum_{j=1}^{n}\bigl((\alpha+1)j-n\bigr)a_jc_{n-j}.
\end{equation}

\paragraph{Reversion.} If $A(0)=0$, $a_1\ne0$, and
$B=A^{\langle-1\rangle}=\sum_{n\ge1}b_nz^n$, then
\begin{equation}\label{unif:eq-reversion-triangular}
 b_1=a_1^{-1},\qquad
 b_n=-a_1^{-1}\sum_{k=2}^{n}a_k[z^n]B(z)^k\quad(n\ge2).
\end{equation}
On the right, only $b_1,\ldots,b_{n-1}$ occur.
\end{theorem}
\begin{proof}
The product, reciprocal, and quotient recurrences come directly from
$C=AB$, $AC=1$, and $AC=B$. For the logarithm, differentiate $A=e^L$ to get
$A'=AL'$ and compare $[z^{n-1}]$; solving for $\ell_n$ or $a_n$ gives the two
recurrences. For $C=A^\alpha$, logarithmic differentiation gives
$AC'=\alpha A'C$; comparison of $[z^{n-1}]$ and isolation of $c_n$ gives
\eqref{unif:eq-power-recursion}. Finally, $A(B)=z$ implies
$a_1b_n+\sum_{k=2}^{n}a_k[z^n]B^k=0$ for $n\ge2$, proving the reversion
recurrence. A coefficient $b_n$ cannot occur in $[z^n]B^k$ when $k\ge2$, so
the system is triangular.
\end{proof}

\begin{proposition}[Newton reversion doubles precision]
\label{unif:prop-newton-reversion}
Suppose $B$ satisfies $A(B)=z+O(z^m)$. Define
\begin{equation}\label{unif:eq-newton-reversion}
 \widetilde B=B-\frac{A(B)-z}{A'(B)}.
\end{equation}
Then $A(\widetilde B)=z+O(z^{2m})$.
\end{proposition}
\begin{proof}
Let $E=A(B)-z=O(z^m)$ and $\delta=-E/A'(B)=O(z^m)$. Formal Taylor expansion
gives
\[
 A(B+\delta)=A(B)+A'(B)\delta+\frac12A''(B)\delta^2+\cdots
 =z+O(E^2)=z+O(z^{2m}).
\]
Because $A'(B)$ has invertible constant term $a_1$, every division is valid.
\end{proof}

The triangular recurrence is transparent for symbolic work; Newton iteration
is asymptotically preferable at very high order. Both are independently
checkable against Lagrange's coefficient formula, providing a useful audit of
series-reversion computations.

\part{Bernoulli Families, Summation, and Factorial Special Functions}

\chapter{Bernoulli, Euler, Genocchi, and Cauchy polynomials}

The supplied MathWorld dossier contains several polynomial families that are
closely related but conventionally assigned overlapping names. Bell numbers
remain $\Bell n$ and Bernoulli numbers remain $\beta_n$; write
$\BerP n x$, $\EulP n x$, and $\GenP n x$ for Bernoulli, Euler, and Genocchi
polynomials.

\section{Bernoulli polynomials}

\begin{definition}
\begin{equation}\label{unif:eq-bernoulli-polynomial-egf}
 \frac{te^{xt}}{e^t-1}=\sum_{n\ge0}\BerP n x\frac{t^n}{n!},
 \qquad \BerP n0=\beta_n.
\end{equation}
\end{definition}

\begin{theorem}[Bernoulli-polynomial calculus]
\label{unif:thm-bernoulli-polynomial-calculus}
For $n\ge0$,
\begin{align}
 \BerP n x&=\sum_{k=0}^{n}\binom nk\beta_kx^{n-k},
 \label{unif:eq-bernoulli-explicit}\\
 \frac{\dd}{\dd x}\BerP n x&=n\BerP{n-1}x,
 \label{unif:eq-bernoulli-appell}\\
 \BerP n{x+1}-\BerP n x&=nx^{n-1}\quad(n\ge1),
 \label{unif:eq-bernoulli-difference}\\
 \BerP n{1-x}&=(-1)^n\BerP n x.
 \label{unif:eq-bernoulli-reflection}
\end{align}
For every positive integer $m$,
\begin{equation}\label{unif:eq-bernoulli-multiplication}
 \BerP n{mx}=m^{n-1}\sum_{r=0}^{m-1}\BerP n{x+r/m},
\end{equation}
and for $p\ge0$, $N\in\N$,
\begin{equation}\label{unif:eq-faulhaber-shifted}
 \sum_{j=0}^{N-1}(x+j)^p
 =\frac{\BerP{p+1}{x+N}-\BerP{p+1}x}{p+1}.
\end{equation}
\end{theorem}
\begin{proof}
Multiply the Bernoulli-number EGF by $e^{xt}$ to obtain the explicit formula.
Differentiation with respect to $x$ multiplies the EGF by $t$, proving the
Appell law. The difference of the EGFs at $x+1$ and $x$ is $te^{xt}$, which
proves the finite-difference identity. Replacing $t$ by $-t$ gives
\[
 \frac{-te^{-xt}}{e^{-t}-1}=\frac{te^{(1-x)t}}{e^t-1},
\]
proving reflection.

For multiplication, sum the EGFs at $x+r/m$ and use the finite geometric sum;
after replacing $t$ by $mt$ one obtains
\eqref{unif:eq-bernoulli-multiplication}. Finally,
\eqref{unif:eq-bernoulli-difference} with $n=p+1$ telescopes over
$x,x+1,\ldots,x+N-1$, giving Faulhaber's formula.
\end{proof}

\section{Euler and Genocchi polynomials}

Define
\begin{align}
 \frac{2e^{xt}}{e^t+1}
 &=\sum_{n\ge0}\EulP n x\frac{t^n}{n!},
 \label{unif:eq-euler-polynomial-egf}\\
 \frac{2te^{xt}}{e^t+1}
 &=\sum_{n\ge0}\GenP n x\frac{t^n}{n!}.
 \label{unif:eq-genocchi-polynomial-egf}
\end{align}

\begin{theorem}[Euler--Genocchi identities]
\label{unif:thm-euler-genocchi}
\begin{align}
 \EulP n{x+1}+\EulP n x&=2x^n,
 \label{unif:eq-euler-difference}\\
 \EulP n{1-x}&=(-1)^n\EulP n x,
 \label{unif:eq-euler-reflection}\\
 \EulP n x&=\frac{2^{n+1}}{n+1}
 \left(\BerP{n+1}{(x+1)/2}-\BerP{n+1}{x/2}\right),
 \label{unif:eq-euler-bernoulli}\\
 \GenP n x&=n\EulP{n-1}x\quad(n\ge1),
 \label{unif:eq-genocchi-euler}\\
 \GenP n0&=2(1-2^n)\beta_n.
 \label{unif:eq-genocchi-bernoulli}
\end{align}
Moreover,
\begin{equation}\label{unif:eq-alternating-power-sum}
 \sum_{j=0}^{N-1}(-1)^j(x+j)^n
 =\frac{\EulP n x+(-1)^{N-1}\EulP n{x+N}}2.
\end{equation}
\end{theorem}
\begin{proof}
Adding the Euler EGFs at $x$ and $x+1$ gives $2e^{xt}$; replacing $t$ by
$-t$ gives the reflection identity. To prove
\eqref{unif:eq-euler-bernoulli}, sum the proposed right side over $n$ after
multiplying by $t^n/n!$. With $m=n+1$, the result is
\[
 \frac1t\left(
 \frac{2te^{(x+1)t}}{e^{2t}-1}-
 \frac{2te^{xt}}{e^{2t}-1}\right)
 =\frac{2e^{xt}}{e^t+1}.
\]
The Genocchi EGF is $t$ times the Euler EGF, proving
\eqref{unif:eq-genocchi-euler}. At $x=0$ use
\[
 \frac{2t}{e^t+1}=\frac{2t}{e^t-1}-\frac{4t}{e^{2t}-1}
\]
to obtain \eqref{unif:eq-genocchi-bernoulli}. Finally, multiply
\eqref{unif:eq-euler-difference} at $x+j$ by $(-1)^j$ and telescope.
\end{proof}

\section{Bernoulli polynomials of the second kind}

The Cauchy polynomials of the first kind, also called Bernoulli polynomials of
the second kind, are defined by
\begin{equation}\label{unif:eq-cauchy-polynomial-egf}
 \frac{t}{\log(1+t)}(1+t)^x
 =\sum_{n\ge0}b_n(x)\frac{t^n}{n!}.
\end{equation}
Their values $b_n=b_n(0)$ satisfy $b_n=n!G_n$ for the Gregory coefficients in
the earlier chapter.

\begin{theorem}[Integral, Stirling, and Sheffer formulas]
\label{unif:thm-cauchy-polynomials}
\begin{align}
 b_n(x)&=\int_0^1\falling{x+u}{n}\dd u,
 \label{unif:eq-cauchy-integral}\\
 b_n&=\sum_{k=0}^{n}\frac{s(n,k)}{k+1},
 \label{unif:eq-cauchy-stirling}\\
 b_n(x+1)-b_n(x)&=nb_{n-1}(x),
 \label{unif:eq-cauchy-difference}\\
 \frac{\dd}{\dd x}b_n(x)&=n\falling x{n-1},
 \label{unif:eq-cauchy-derivative}\\
 b_n(x+y)&=\sum_{k=0}^{n}\binom nk b_k(x)\falling y{n-k}.
 \label{unif:eq-cauchy-sheffer}
\end{align}
\end{theorem}
\begin{proof}
Integrate
$(1+t)^{x+u}=\sum_n\falling{x+u}{n}t^n/n!$ from $u=0$ to $1$. The left side
becomes $(1+t)^xt/\log(1+t)$, proving the integral. At $x=0$, expand
$\falling un=\sum_ks(n,k)u^k$ and integrate to get the Stirling sum. The
difference of the EGFs at $x+1$ and $x$ is $t$ times the original EGF.
Differentiation with respect to $x$ gives $t(1+t)^x$, yielding the derivative.
Finally, factor the EGF at $x+y$ into the EGF at $x$ times $(1+t)^y$.
\end{proof}

\section{A normalized modified-Bernoulli series}

One supplied source uses ``modified Bernoulli numbers'' for the coefficients
of the following logarithmic hyperbolic series. Because that name is also used
elsewhere, the generating function is the unambiguous definition.

\begin{proposition}
\begin{equation}\label{unif:eq-modified-bernoulli}
 \frac12\log\left(\frac{2\sinh(t/2)}{t}\right)
 =\sum_{n\ge2}\frac{\beta_n}{2n}\frac{t^n}{n!}.
\end{equation}
\end{proposition}
\begin{proof}
Differentiate the left side:
\[
 \frac14\coth(t/2)-\frac1{2t}
 =\frac1{2t}\left(\frac t2\coth(t/2)-1\right).
\]
The Bernoulli EGF gives
$(t/2)\coth(t/2)=1+\sum_{m\ge1}\beta_{2m}t^{2m}/(2m)!$.
Integrate term by term and use the zero constant term.
\end{proof}

\chapter{Complementary Bell numbers}

Define
\begin{equation}\label{unif:eq-complementary-bell-definition}
 C_n=\sum_{k=0}^{n}(-1)^k\stirtwo nk=T_n(-1).
\end{equation}
This is the signed difference between set partitions with an even and an odd
number of blocks.

\begin{theorem}[Generating function, recurrence, and Dobi\'nski form]
\label{unif:thm-complementary-bell}
\begin{align}
 \sum_{n\ge0}C_n\frac{z^n}{n!}&=\exp(1-e^z),
 \label{unif:eq-complementary-bell-egf}\\
 C_{n+1}&=-\sum_{j=0}^{n}\binom njC_j,
 \label{unif:eq-complementary-bell-recurrence}\\
 C_n&=e\sum_{m\ge0}\frac{(-1)^m m^n}{m!}.
 \label{unif:eq-complementary-bell-dobinski}
\end{align}
The initial values are
$1,-1,0,1,1,-2,-9,-9,50,267,413,\ldots$.
\end{theorem}
\begin{proof}
Set $x=-1$ in
$\sum_nT_n(x)z^n/n!=\exp(x(e^z-1))$. If
$F(z)=\exp(1-e^z)$, then $F'=-e^zF$, and EGF coefficient comparison gives the
recurrence. Finally,
\[
 e^{1-e^z}=e\sum_{m\ge0}\frac{(-1)^m}{m!}e^{mz};
\]
local uniform convergence permits coefficient comparison and yields the
Dobi\'nski-type series.
\end{proof}

\chapter{Euler--Maclaurin summation and Darboux formulas}

For $m\ge1$, let $\widetilde B_m(x)=\BerP m{\{x\}}$ be the periodic Bernoulli
function. Away from the integers,
$\widetilde B_m'=m\widetilde B_{m-1}$.

\begin{theorem}[Euler--Maclaurin with exact remainder]
\label{unif:thm-euler-maclaurin}
Let $a<b$ be integers and $f\in C^{2p}([a,b])$. Then
\begin{align}
 \sum_{k=a}^{b}f(k)
 &=\int_a^b f(x)\dd x+\frac{f(a)+f(b)}2
 \notag\\
 &\quad+\sum_{r=1}^{p}\frac{\beta_{2r}}{(2r)!}
 \bigl(f^{(2r-1)}(b)-f^{(2r-1)}(a)\bigr)+R_p,
 \label{unif:eq-euler-maclaurin}
\end{align}
where
\begin{equation}\label{unif:eq-em-remainder}
 R_p=-\frac1{(2p)!}\int_a^b
 \widetilde B_{2p}(x)f^{(2p)}(x)\dd x.
\end{equation}
Moreover,
\begin{equation}\label{unif:eq-em-bound}
 |R_p|\le\frac{2\zeta(2p)}{(2\pi)^{2p}}
 \int_a^b|f^{(2p)}(x)|\dd x.
\end{equation}
\end{theorem}
\begin{proof}
On $[k,k+1]$, integration by parts with
$\widetilde B_1(x)=x-k-1/2$ gives
\[
 \int_k^{k+1}\widetilde B_1f'
 =\frac{f(k)+f(k+1)}2-\int_k^{k+1}f.
\]
Summing and rearranging yields
\[
 \sum_{k=a}^{b}f(k)=\int_a^bf+\frac{f(a)+f(b)}2
 +\int_a^b\widetilde B_1f'.
\]
Repeated integration by parts, using
$\widetilde B_m'=m\widetilde B_{m-1}$, produces the boundary corrections.
At integer endpoints $\widetilde B_m=\beta_m$, and
$\beta_{2r+1}=0$ for $r\ge1$, leaving only the even terms and the remainder
\eqref{unif:eq-em-remainder}.

The Fourier expansion
\[
 \widetilde B_{2p}(x)=(-1)^{p+1}\frac{2(2p)!}{(2\pi)^{2p}}
 \sum_{m\ge1}\frac{\cos(2\pi mx)}{m^{2p}}
\]
implies
$|\widetilde B_{2p}(x)|/(2p)!\le2\zeta(2p)/(2\pi)^{2p}$.
Substitution into the exact remainder proves the bound.
\end{proof}

\begin{corollary}[Stirling expansion]
For fixed $p\ge1$ and $x\to+\infty$,
\begin{align}
 \log\Gamma(x)
 &=(x-\tfrac12)\log x-x+\tfrac12\log(2\pi)
 \notag\\
 &\quad+\sum_{r=1}^{p-1}
 \frac{\beta_{2r}}{2r(2r-1)x^{2r-1}}+O(x^{-2p+1}).
 \label{unif:eq-stirling-expansion}
\end{align}
\end{corollary}
\begin{proof}
Apply Euler--Maclaurin to
\[
 \sum_{k=1}^{N}\log k=\log\Gamma(N+1).
\]
The derivatives of $\log x$ satisfy
\[
 \frac{d^m}{dx^m}\log x=(-1)^{m-1}(m-1)!x^{-m},
\]
and the remainder bound gives the stated order. Wallis' product fixes the
constant as $\frac12\log(2\pi)$. The gamma recurrence extends the
expansion from integers to positive real $x$.
\end{proof}

\section{Darboux's weighted Taylor identity}

\begin{theorem}[Darboux's formula]
\label{unif:thm-darboux-weighted}
Let $\phi$ be a polynomial of exact degree $n$, let
$\Delta=z-z_0$, and suppose $f$ has $n+1$ derivatives on the joining segment.
Then
\begin{align}
 f(z)=f(z_0)+\frac1{\phi^{(n)}(0)}\Bigg[&
 (-1)^n\Delta^{n+1}\int_0^1\phi(t)
 f^{(n+1)}(z_0+t\Delta)\dd t
 \notag\\
 &+\sum_{k=1}^{n}(-1)^{k-1}\Delta^k
 \bigl(\phi^{(n-k)}(1)f^{(k)}(z)
 -\phi^{(n-k)}(0)f^{(k)}(z_0)\bigr)\Bigg].
 \label{unif:eq-darboux-weighted}
\end{align}
\end{theorem}
\begin{proof}
Set $h(t)=f(z_0+t\Delta)$. Since $\phi^{(n+1)}=0$, $n+1$ integrations by
parts give
\[
 \int_0^1\phi h^{(n+1)}
 =\sum_{k=0}^{n}(-1)^{n-k}
 [\phi^{(n-k)}h^{(k)}]_0^1.
\]
Multiply by $(-1)^n$. The $k=0$ term is
$\phi^{(n)}(0)(h(1)-h(0))$, because $\phi^{(n)}$ is constant. Substitute
$h^{(k)}(t)=\Delta^kf^{(k)}(z_0+t\Delta)$ and solve for $f(z)-f(z_0)$.
\end{proof}

Choosing $\phi(t)=(1-t)^n$ recovers Taylor's theorem with integral remainder.
Other choices balance endpoint derivatives and lead to quadrature and
interpolation error formulas.

\section{Algebraic singularity transfer}

\begin{theorem}[Basic transfer formula]
\label{unif:thm-basic-transfer}
If $\rho\ne0$ and
$\alpha\notin\{0,-1,-2,\ldots\}$, then
\begin{align}
 [z^n](1-z/\rho)^{-\alpha}
 &=\rho^{-n}\frac{\Gamma(n+\alpha)}{\Gamma(\alpha)\Gamma(n+1)}
 \notag\\
 &=\frac{\rho^{-n}n^{\alpha-1}}{\Gamma(\alpha)}
 \left(1+\frac{\alpha(\alpha-1)}{2n}+O(n^{-2})\right).
 \label{unif:eq-basic-transfer}
\end{align}
If $F$ is a finite sum of such terms plus a function analytic in $|z|<R$,
$R>|\rho|$, then its coefficients are the corresponding finite sum of transfer
expansions plus $O(R_0^{-n})$ for every $|\rho|<R_0<R$.
\end{theorem}
\begin{proof}
The generalized binomial theorem gives the exact gamma quotient. Stirling's
formula for the quotient $\Gamma(n+\alpha)/\Gamma(n+1)$ gives the asymptotic.
The analytic remainder is bounded by Cauchy's estimate on $|z|=R_0$.
\end{proof}

In particular,
\begin{equation}\label{unif:eq-square-root-transfer}
 [z^n]c(1-z/\rho)^{1/2}
 \sim-\frac{c}{2\sqrt\pi}\rho^{-n}n^{-3/2}.
\end{equation}

\begin{theorem}[Simple critical inverse]
\label{unif:thm-critical-inverse}
If $f$ is analytic near $a$, $b=f(a)$, $f'(a)=0$, and $f''(a)\ne0$, then the
two inverse branches have convergent Puiseux expansions
\begin{equation}\label{unif:eq-critical-inverse-expansion}
 g_{\pm}(z)=a\pm\sqrt{\frac{2(z-b)}{f''(a)}}
 -\frac{f'''(a)}{3f''(a)^2}(z-b)+O((z-b)^{3/2}).
\end{equation}
\end{theorem}
\begin{proof}
Factor
$f(a+u)-b=u^2h(u)$ with $h(0)=f''(a)/2\ne0$. Choose an analytic square root
$q=\sqrt h$ and put $v=uq(u)$. Since $v'(0)\ne0$, $v$ has an analytic inverse
$u=\chi(v)$, and the two branches are
$a+\chi(\pm\sqrt{z-b})$. For the coefficients, set
$f(a+u)-b=Au^2+Bu^3+O(u^4)$ and
$u=c_1s+c_2s^2+O(s^3)$, $s^2=z-b$. Then
$Ac_1^2=1$ and $2Ac_1c_2+Bc_1^3=0$, which gives the displayed terms.
\end{proof}

Together with \eqref{unif:eq-square-root-transfer}, this explains the ubiquitous
$n^{-3/2}$ factor in coefficients of inverse functions at a simple critical
value. The full dented-domain transfer theorem is due to Flajolet and Odlyzko.

\chapter{Pochhammer symbols, polygamma functions, and Barnes' \texorpdfstring{$G$}{G}}

\section{Rising factorials and parameter derivatives}

For the rising Pochhammer symbol,
\begin{equation}\label{unif:eq-pochhammer-gamma}
 (a)_n=\rising an=\frac{\Gamma(a+n)}{\Gamma(a)}.
\end{equation}

\begin{theorem}[Binomial and Vandermonde formulas]
\begin{align}
 (1-z)^{-a}&=\sum_{n\ge0}(a)_n\frac{z^n}{n!},
 \label{unif:eq-pochhammer-binomial}\\
 (a+b)_n&=\sum_{k=0}^{n}\binom nk(a)_k(b)_{n-k}.
 \label{unif:eq-pochhammer-vandermonde}
\end{align}
\end{theorem}
\begin{proof}
The generalized binomial theorem gives the first identity because
$(-1)^n\binom{-a}{n}=(a)_n/n!$. Multiply the series with parameters $a$ and
$b$ and compare coefficients to obtain Vandermonde's convolution.
\end{proof}

Let $\psi=\Gamma'/\Gamma$. Define
\begin{equation}\label{unif:eq-pochhammer-log-derivatives}
 L_j(a,n)=\psi^{(j-1)}(a+n)-\psi^{(j-1)}(a)
 =(-1)^{j-1}(j-1)!\sum_{q=0}^{n-1}(a+q)^{-j}.
\end{equation}

\begin{theorem}[Parameter derivatives through Bell polynomials]
\label{unif:thm-pochhammer-parameter-derivatives}
For $r\ge0$,
\begin{equation}\label{unif:eq-pochhammer-bell-derivative}
 \frac{\partial^r}{\partial a^r}(a)_n
 =(a)_n\BellC r(L_1(a,n),\ldots,L_r(a,n)).
\end{equation}
\end{theorem}
\begin{proof}
Take logarithms of the gamma quotient. Its $j$th derivative is the first
expression in \eqref{unif:eq-pochhammer-log-derivatives}; the gamma recurrence
gives the finite sum. Now apply the exponential special case of Fa\`a di
Bruno to $(a)_n=\exp(\log(a)_n)$.
\end{proof}

\section{Digamma and polygamma series}

\begin{theorem}
Away from the nonpositive integers,
\begin{equation}\label{unif:eq-digamma-series}
 \psi(z)=-\gamma+\sum_{k=0}^{\infty}
 \left(\frac1{k+1}-\frac1{k+z}\right).
\end{equation}
For $m\ge1$ and $\Re z>0$,
\begin{equation}\label{unif:eq-polygamma-series}
 \psi^{(m)}(z)=(-1)^{m+1}m!
 \sum_{k=0}^{\infty}(z+k)^{-m-1},
\end{equation}
and
\begin{equation}\label{unif:eq-digamma-reflection}
 \psi(1-z)-\psi(z)=\pi\cot(\pi z).
\end{equation}
\end{theorem}
\begin{proof}
Logarithmically differentiate Euler's Weierstrass product
\[
 \Gamma(z)^{-1}=ze^{\gamma z}
 \prod_{k\ge1}(1+z/k)e^{-z/k}
\]
to obtain the digamma series. Termwise differentiation is locally uniformly
convergent for $m\ge1$ and yields the polygamma series. Logarithmically
differentiate Euler's reflection formula
$\Gamma(z)\Gamma(1-z)=\pi/\sin(\pi z)$ to obtain reflection.
\end{proof}

Thus
\begin{equation}\label{unif:eq-polygamma-hurwitz}
 \psi^{(m)}(z)=(-1)^{m+1}m!\zeta(m+1,z).
\end{equation}

\section{Barnes' \texorpdfstring{$G$}{G}-function}

Define
\begin{align}
 G(z+1)&=(2\pi)^{z/2}
 \exp\left(-\frac{z+(1+\gamma)z^2}{2}\right)
 \notag\\
 &\quad\times\prod_{n=1}^{\infty}
 \left(1+\frac zn\right)^n
 \exp\left(\frac{z^2}{2n}-z\right).
 \label{unif:eq-barnes-product}
\end{align}

\begin{theorem}[Functional equation]
\label{unif:thm-barnes-functional}
The product converges locally uniformly, defines an entire function, and
satisfies
\begin{equation}\label{unif:eq-barnes-functional}
 G(z+1)=\Gamma(z)G(z),\qquad G(1)=1.
\end{equation}
Hence
\begin{equation}\label{unif:eq-barnes-integers}
 G(n+1)=\prod_{k=1}^{n-1}k!.
\end{equation}
\end{theorem}
\begin{proof}
The logarithm of the $n$th factor is
$n\log(1+z/n)-z+z^2/(2n)=O(n^{-2})$ uniformly on compact sets, proving local
uniform convergence. For the recurrence, take the ratio of the products at
$z+1$ and $z$ and truncate at $N$. The powers telescope:
\[
 \prod_{n=1}^{N}\left(\frac{n+z+1}{n+z}\right)^n
 =\frac{(N+z+1)^N}{\prod_{n=1}^{N}(n+z)}.
\]
After including the exponential prefactors, use
$H_N=\log N+\gamma+1/(2N)+O(N^{-2})$ and Stirling's formula. The numerator is
asymptotic to $N!N^z$, so the ratio tends to
\[
 \lim_{N\to\infty}\frac{N!N^z}{(z+1)\cdots(z+N)}=\Gamma(z+1).
\]
Thus $G(z+2)=\Gamma(z+1)G(z+1)$, equivalent to the stated recurrence.
Iteration gives the integer product.
\end{proof}

\begin{theorem}[Barnes asymptotic expansion]
\label{unif:thm-barnes-asymptotic}
For fixed $p\ge1$ and $z\to+\infty$,
\begin{align}
 \log G(z+1)
 &=\frac{z^2}{2}\log z-\frac{3z^2}{4}
 +\frac z2\log(2\pi)-\frac1{12}\log z+\zeta'(-1)
 \notag\\
 &\quad+\sum_{k=1}^{p-1}
 \frac{\beta_{2k+2}}{4k(k+1)z^{2k}}+O(z^{-2p}).
 \label{unif:eq-barnes-asymptotic}
\end{align}
\end{theorem}
\begin{proof}
For an integer $N$,
\[
 \log G(N+1)=N\log\Gamma(N)-\sum_{j=1}^{N-1}j\log j.
\]
The first term is expanded by \eqref{unif:eq-stirling-expansion}. To expand
the second, apply Euler--Maclaurin to $x^{-s}$, continue the resulting finite
formula to a neighborhood of $s=-1$, and differentiate there. This gives
\begin{align*}
 \sum_{j=1}^{N-1}j\log j
 &=\left(\frac{N^2}{2}-\frac N2+\frac1{12}\right)\log N
 -\frac{N^2}{4}+\frac1{12}-\zeta'(-1)\\
 &\quad-\sum_{k=1}^{p-1}
 \frac{\beta_{2k+2}}{(2k)(2k+1)(2k+2)N^{2k}}
 +O(N^{-2p}).
\end{align*}
Substitution and cancellation yield
\eqref{unif:eq-barnes-asymptotic} for integers. The functional equation and
the uniform Stirling expansion extend it across each bounded fractional
interval, proving the positive-real statement.
\end{proof}

\part{Inverse Examples and Refined Catalan Geometry}

\chapter{Inverse trigonometric functions and Kepler's equation}

\section{Inverse trigonometric coefficient evaluations}

\begin{proposition}
For $|y|<1$,
\begin{align}
 \arcsin y&=\sum_{m\ge0}\frac1{2m+1}\binom{2m}{m}
 \frac{y^{2m+1}}{4^m},
 \label{unif:eq-arcsin-series}\\
 \arctan y&=\sum_{m\ge0}\frac{(-1)^m}{2m+1}y^{2m+1}.
 \label{unif:eq-arctan-series}
\end{align}
Consequently,
\begin{align}
 [x^{2m}]\left(\frac{x}{\sin x}\right)^{2m+1}
 &=4^{-m}\binom{2m}{m},
 \label{unif:eq-sine-power-coefficient}\\
 [x^{2m}]\left(\frac{x}{\tan x}\right)^{2m+1}
 &=(-1)^m.
 \label{unif:eq-tangent-power-coefficient}
\end{align}
\end{proposition}
\begin{proof}
Integrate the binomial series for $(1-y^2)^{-1/2}$ and the geometric series
for $(1+y^2)^{-1}$ to obtain the inverse-function expansions. Lagrange
inversion for $f(x)=\sin x$ and $f(x)=\tan x$ says
\[
 [y^{2m+1}]f^{\langle-1\rangle}(y)
 =\frac1{2m+1}[x^{2m}]\left(\frac{x}{f(x)}\right)^{2m+1}.
\]
Comparison with the two explicit inverse series proves the coefficient
identities.
\end{proof}

\section{Kepler's equation}

For $|e|<1$, Kepler's equation
\begin{equation}\label{unif:eq-kepler}
 E=M+e\sin E
\end{equation}
has a unique real solution because $E\mapsto E-e\sin E$ has derivative
$1-e\cos E>0$.

\begin{theorem}[Lagrange and Fourier--Bessel expansions]
\label{unif:thm-kepler-expansions}
The solution satisfies
\begin{align}
 E&=M+\sum_{n\ge1}\frac{e^n}{n!}
 \frac{\dd^{n-1}}{\dd M^{n-1}}\sin^nM,
 \label{unif:eq-kepler-lagrange}\\
 E&=M+2\sum_{n\ge1}\frac{J_n(ne)}{n}\sin(nM).
 \label{unif:eq-kepler-bessel}
\end{align}
\end{theorem}
\begin{proof}
Set $u=E-M$. Then $u=e\sin(M+u)$, and Lagrange--B\"urmann in the parameter
$e$ gives
\[
 [e^n]u=\frac1n[t^{n-1}]\sin^n(M+t)
 =\frac1{n!}\frac{\dd^{n-1}}{\dd M^{n-1}}\sin^nM.
\]
For the Fourier form, $u(M)$ is odd and $2\pi$-periodic. If
$u=\sum_{n\ge1}b_n\sin(nM)$, integration by parts and $u(\pm\pi)=0$ give
\[
 b_n=\frac1{\pi n}\int_{-\pi}^{\pi}E'(M)\cos(nM)\dd M.
\]
Change variables using $M=E-e\sin E$ and $E'(M)\dd M=\dd E$:
\[
 b_n=\frac1{\pi n}\int_{-\pi}^{\pi}
 \cos(nE-ne\sin E)\dd E=\frac{2J_n(ne)}n,
\]
by the standard integer-order Bessel integral. This proves the second series.
\end{proof}

\chapter{Narayana refinement and associahedral face enumeration}

The earlier Catalan chapter proves the total count of binary trees. The
supplied associahedron and Catalan references also point to two refinements:
Narayana numbers record peaks, and Schr\"oder trees record faces rather than
only vertices.

\section{Narayana numbers}

Let $C(z,u)$ count Dyck paths, with $z$ marking semilength and $u$ marking
peaks. A nonempty path decomposes uniquely as $UADB$, with $A,B$ Dyck paths.
If $A$ is empty, the initial $UD$ is a peak; otherwise no new peak is created.
Thus
\begin{equation}\label{unif:eq-narayana-functional}
 C(z,u)=1+zC(z,u)(C(z,u)+u-1).
\end{equation}

\begin{theorem}[Narayana formula]
\label{unif:thm-narayana}
The number of Dyck paths of semilength $n$ with $k$ peaks is
\begin{equation}\label{unif:eq-narayana}
 N(n,k)=\frac1n\binom nk\binom n{k-1},
 \qquad1\le k\le n.
\end{equation}
Consequently,
\begin{equation}\label{unif:eq-narayana-catalan}
 \sum_{k=1}^{n}N(n,k)=\frac1{n+1}\binom{2n}{n}.
\end{equation}
\end{theorem}
\begin{proof}
Put $Y=C-1$. Equation \eqref{unif:eq-narayana-functional} becomes
$Y=z(1+Y)(u+Y)$. Lagrange inversion gives
\begin{align*}
 [z^nu^k]Y
 &=\frac1n[t^{n-1}u^k](1+t)^n(u+t)^n\\
 &=\frac1n\binom nk\binom n{k-1},
\end{align*}
proving \eqref{unif:eq-narayana}. Summing over $k$ and applying Vandermonde's
identity gives
\[
 \frac1n\sum_k\binom nk\binom n{k-1}
 =\frac1n\binom{2n}{n-1}=\frac1{n+1}\binom{2n}{n}.
\]
\end{proof}

The symmetry $N(n,k)=N(n,n-k+1)$ is immediate from the closed formula. In the
associahedron, these numbers form the $h$-vector; combinatorially, they also
refine binary trees by several equidistributed statistics.

\section{All faces of the associahedron}

Use the normalization of the earlier associahedral inversion theorem: $K_n$
has dimension $n-1$, vertices indexed by full plane binary trees with $n+1$
leaves, and general faces indexed by plane rooted trees with $n+1$ leaves and
internal arity at least $2$.

\begin{theorem}[Face numbers of $K_n$]
\label{unif:thm-associahedron-f-vector}
For $0\le d\le n-1$, the number of $d$-dimensional faces is
\begin{equation}\label{unif:eq-associahedron-f-vector}
 f_d(K_n)=\frac1{n-d}\binom{2n-d}{n-1-d}\binom{n-1}{d}.
\end{equation}
In particular, $f_0(K_n)=C_n$, $f_{n-2}(K_n)=n(n+1)/2-1$, and
$f_{n-1}(K_n)=1$.
\end{theorem}
\begin{proof}
Let $T(z,u)$ count plane rooted trees, with $z$ marking leaves and $u$ marking
internal vertices, every internal vertex having at least two ordered children.
The root is either a leaf or an internal vertex with a sequence of at least two
subtrees, so
\[
 T=z+u\frac{T^2}{1-T}.
\]
Solving for the Lagrange form gives
\[
 T=z\frac{1-T}{1-(1+u)T}.
\]
If a tree has $L$ leaves and $m$ internal vertices, Lagrange inversion yields
\begin{align*}
 [z^Lu^m]T
 &=\frac1L[t^{L-1}u^m]
 \left(\frac{1-t}{1-(1+u)t}\right)^L\\
 &=\frac1L\binom{L+m-1}{m}\binom{L-2}{m-1}.
\end{align*}
A face of $K_n$ has $L=n+1$ leaves. If it has $m$ internal vertices, its
dimension is $d=n-m$. Substitute $L=n+1$ and $m=n-d$ and simplify:
\[
 \frac1{n+1}\binom{2n-d}{n-d}\binom{n-1}{d}
 =\frac1{n-d}\binom{2n-d}{n-1-d}\binom{n-1}{d}.
\]
This is \eqref{unif:eq-associahedron-f-vector}.
\end{proof}

For $K_4$, the formula gives $14$ vertices, $21$ edges, $9$ facets, and one
three-dimensional body, agreeing with the coefficients appearing in the
fourth inverse-series polynomial.

\section{Loday coordinates and a constant-sum proof}

For a full binary tree with $n+1$ leaves, order its $n$ internal vertices from
left to right. At internal vertex $i$, let $\ell_i$ and $r_i$ be the numbers of
leaves in its left and right subtrees. Loday assigns the coordinate
$x_i=\ell_i r_i$.

\begin{proposition}[Loday's constant hyperplane]
\label{unif:prop-loday-hyperplane}
Every such coordinate vector satisfies
\begin{equation}\label{unif:eq-loday-hyperplane}
 \sum_{i=1}^{n}x_i=\binom{n+1}{2}.
\end{equation}
\end{proposition}
\begin{proof}
At vertex $i$, the product $\ell_i r_i$ counts unordered pairs of leaves with
one leaf in each child subtree. Every unordered pair of distinct leaves has a
unique lowest common ancestor, and it is counted precisely at that internal
vertex. Hence the sum counts all pairs of the $n+1$ leaves exactly once.
\end{proof}

Loday proved that the convex hull of these integer vectors is a realization of
the associahedron. The proposition supplies the elementary reason it lies in
an $(n-1)$-dimensional affine hyperplane; the nontrivial part of the realization
theorem is the identification of its face lattice with noncrossing diagonals.


\part{Source Concordance and Verification Record}

\chapter{How the five drafts were unified}
\label{unif:ch-concordance}

The archive contains five substantial but overlapping LaTeX drafts. They share
a common mathematical core, but each has sections, examples, and proof details
that are absent from the others. The present article was produced by a
section-level comparison rather than by concatenation. Repeated definitions,
tables, and derivations were retained once in the strongest available form;
unique material was moved to the chapter where it belongs mathematically.
Notation, boundary conventions, and theorem hypotheses were then normalized
across the whole article.

\begin{longtable}{>{\raggedright\arraybackslash}p{0.28\textwidth}
                  >{\raggedright\arraybackslash}p{0.62\textwidth}}
\toprule
\textbf{Draft} & \textbf{Role in the unified article}\\
\midrule
\endfirsthead
\toprule
\textbf{Draft} & \textbf{Role in the unified article}\\
\midrule
\endhead
\bottomrule
\endlastfoot

\nolinkurl{combinatorial_coefficient_calculus.tex}
& Supplied the broadest dependency-based architecture and the most complete
  core proofs for Stirling, Bell, Eulerian, Fa\`a di Bruno, Lagrange inversion,
  asymptotics, and restricted variants. It became the structural backbone.\\[1ex]

\nolinkurl{Formulae_Monograph.tex} in
\nolinkurl{Formulae_Monograph_Package}
& Contributed independent low-order checks, finite computational certificates,
  additional examples, and several alternative derivations used to audit the
  backbone.\\[1ex]

\nolinkurl{Formulae_Monograph.tex} in
\nolinkurl{Formulae_Monograph_Source_and_PDF}
& Contributed the Kepler-equation example, further tree interpretations,
  source-normalization notes, and additional reversion details.\\[1ex]

\nolinkurl{Combinatorial_Formulae_and_Inversion_Theorems.tex}
& Contributed material on signed Bell evaluations, Poisson shift methods,
  symmetric summations, negative-index extensions, and implementation-oriented
  checks. Overlapping Bell and Eulerian material was deduplicated.\\[1ex]

\nolinkurl{Formulae_Compendium.tex}
& Contributed the multinomial Leibniz rule, coordinate-free Fa\`a di Bruno,
  recursive inverse derivatives, partition-lattice M\"obius inversion, Sheffer
  sequences, exponential Riordan arrays, Abel polynomials, multivariate
  inversion, and the coefficient-algorithm toolkit.\\
\end{longtable}

\section{Use of the supplied reference folders}

The \nolinkurl{Wikipedia} folder was used as a topic checklist and as a source
of complementary interpretations: set partitions, Dyck paths, triangulations,
associahedra, Eulerian descents, and the historical forms of Lagrange reversion.
Statements were not copied mechanically; indexing was reconciled with the
article's conventions and each incorporated formula was proved.

The \nolinkurl{MathWorld} folder supplied independent formula collections and
Mathematica notebooks for Bell and complementary Bell numbers, Bernoulli and
Euler polynomial families, Bernoulli polynomials of the second kind,
Dobi\'nski's formula, Euler--Maclaurin summation, Darboux's weighted formula,
Pochhammer symbols, polygamma functions, Genocchi numbers, Catalan numbers, and
Barnes' $G$-function. These were integrated in the new parts on Bernoulli
families, summation, factorial special functions, and Catalan geometry.

Online primary references were consulted for Good's multivariate inversion,
multivariate Fa\`a di Bruno, singularity analysis, Barnes' $G$-function, the
Riordan group, Bernoulli numbers of the second kind, and Loday's associahedron
realization. Their role is bibliographic and confirmatory; the formulas used in
the article are derived in the text.

\section{Deduplication rules}

\begin{auditbox}
The following editorial rules were applied globally.
\begin{enumerate}
\item A definition appears once, at its first logical use. Later appearances
      are replaced by cross-references or specializations.
\item Equivalent generating functions are retained together only when they
      illuminate different structures, such as ordinary versus exponential
      normalization.
\item Repeated tables are consolidated; a table is not used as a substitute
      for a recurrence or proof.
\item Competing index conventions are translated into one convention before
      formulas are compared. Type-$B$ Eulerian numbers and associahedra required
      particular care.
\item Analytic and formal statements are separated. Formal coefficient
      identities do not inherit unstated convergence claims, while analytic
      versions include their local-invertibility or regularity hypotheses.
\item Historical computations, prime searches, and modular-period reports are
      described as finite computations with certificate procedures rather than
      elevated to symbolic theorems.
\end{enumerate}
\end{auditbox}

\chapter{Proof and consistency audit}

\section{Proof coverage}

Every theorem, lemma, proposition, corollary, and identity introduced in the
unified article has either an adjacent proof or an explicit derivation from a
previously proved result. The only proofless theorem in the structural backbone
was the cube Ehrhart statement; its proof has been inserted by comparing
\[
 \sum_{m\ge0}(m+1)^nt^m
\]
with the Eulerian power-series identity. New reference-driven results---Good's
theorem, Fr\'echet Fa\`a di Bruno, Euler--Maclaurin, Darboux's identity,
Pochhammer parameter differentiation, Barnes' functional equation, Narayana
numbers, and the associahedron face formula---are proved in full in their
respective chapters.

Deep realization or continuation results that are not needed for coefficient
calculations are identified as such. For example, the constant-hyperplane
property of Loday's coordinates is proved here, while the full face-lattice
realization theorem is attributed to Loday rather than silently reproved in an
abbreviated form.

\section{Independent low-order checks}

Several central transformations admit independent certificates.

\paragraph{Composition through fourth order.}
For $F=f\circ g$, the partition formula gives
\begin{align*}
 F'&=f'g',\\
 F''&=f''(g')^2+f'g'',\\
 F'''&=f'''(g')^3+3f''g'g''+f'g''',\\
 F^{(4)}&=f^{(4)}(g')^4+6f'''(g')^2g''
 +3f''(g'')^2+4f''g'g'''+f'g^{(4)}.
\end{align*}
The coefficients $1,6,3,4,1$ are exactly the block-profile counts in
$\Pi_4$, independently verifying the fourth Bell-polynomial row.

\paragraph{Reversion through fourth order.}
For
$A(x)=x+a_1x^2+a_2x^3+a_3x^4+O(x^5)$ and
$B=A^{\langle-1\rangle}$,
\[
 B(x)=x-a_1x^2+(2a_1^2-a_2)x^3
 +(-5a_1^3+5a_1a_2-a_3)x^4+O(x^5).
\]
Direct substitution gives $A(B(x))=x+O(x^5)$; the same coefficients follow
from Lagrange inversion and from the triangular recurrence
\eqref{unif:eq-reversion-triangular}.

\paragraph{Stirling matrix inversion.}
The Riordan product
\[
 (1,e^z-1)(1,\log(1+z))=(1,z)
\]
certifies the inverse relation between the second-kind and signed first-kind
Stirling matrices. This check is independent of the factorial-basis proof.

\paragraph{Bernoulli--Euler--Genocchi check.}
Multiplying
$\GenP n x=n\EulP{n-1}x$ by $t^n/n!$ recovers the Genocchi EGF, while
substitution of \eqref{unif:eq-euler-bernoulli} recovers
$\GenP n0=2(1-2^n)\beta_n$. Thus the two routes through the generating-function
triangle agree coefficientwise.

\section{Boundary conventions}

All triangular arrays are extended by zero outside their natural ranges.
Empty products equal $1$, and the coefficient convention $0^0=1$ is used only
where a generating-function coefficient demands it. Rising and falling
factorials of order zero equal $1$. Binomial coefficients with an out-of-range
lower index vanish. These conventions make recurrences valid at boundaries and
remove the need for repeated exceptional clauses.

\section{Dependency map}

The principal dependency chain is
\[
 \begin{gathered}
 \text{factorial bases}
 \longrightarrow \text{Stirling transforms}
 \longrightarrow \text{partition and Bell calculus},\\
 \text{partition and Bell calculus}
 \longrightarrow \text{composition and inversion}.
 \end{gathered}
\]
Eulerian numbers provide the parallel change of coordinates between powers and
shifted binomial coefficients. Sheffer and Riordan theory packages the same
operations as polynomial and matrix calculus. Bernoulli polynomials then turn
finite differences into summation, Euler--Maclaurin turns summation into
asymptotic analysis, and Lagrange--Good inversion converts implicit equations
into explicit coefficients. The recurring operations are therefore not a
loose catalogue: they are change of basis, labelled decomposition, logarithmic
inversion, and coefficient extraction.


\backmatter

\chapter*{Selected bibliography and source provenance}
\addcontentsline{toc}{chapter}{Selected bibliography and provenance}

\begin{thebibliography}{99}

\bibitem{Bell1934}
E.~T. Bell,
``Exponential polynomials,''
\emph{Annals of Mathematics}, second series, vol.~35 (1934), pp.~258--277.

\bibitem{Broder1984}
A.~Z. Broder,
``The $r$-Stirling numbers,''
\emph{Discrete Mathematics}, vol.~49 (1984), pp.~241--259.

\bibitem{Charalambides2002}
C.~A. Charalambides,
\emph{Enumerative Combinatorics},
Chapman \& Hall/CRC, 2002.

\bibitem{Claesson2001}
A. Claesson,
``Generalized pattern avoidance,''
\emph{European Journal of Combinatorics}, vol.~22 (2001), pp.~961--971.

\bibitem{Comtet1974}
L. Comtet,
\emph{Advanced Combinatorics: The Art of Finite and Infinite Expansions},
translated by J.~W. Nienhuys, D. Reidel, 1974.

\bibitem{FlajoletSedgewick2009}
P. Flajolet and R. Sedgewick,
\emph{Analytic Combinatorics},
Cambridge University Press, 2009.

\bibitem{GKP1994}
R.~L. Graham, D.~E. Knuth, and O. Patashnik,
\emph{Concrete Mathematics}, second edition,
Addison--Wesley, 1994.

\bibitem{Harper1967}
L.~H. Harper,
``Stirling behavior is asymptotically normal,''
\emph{Annals of Mathematical Statistics}, vol.~38 (1967), pp.~410--414.

\bibitem{Henrici1974}
P. Henrici,
\emph{Applied and Computational Complex Analysis}, vol.~1,
Wiley, 1974.

\bibitem{HsuShiue1998}
L.~C. Hsu and P.~J.-S. Shiue,
``A unified approach to generalized Stirling numbers,''
\emph{Advances in Applied Mathematics}, vol.~20 (1998), pp.~366--384.

\bibitem{Johnson2002}
W.~P. Johnson,
``The curious history of Faà di Bruno's formula,''
\emph{American Mathematical Monthly}, vol.~109 (2002), pp.~217--234.

\bibitem{MoserWyman1955}
L. Moser and M. Wyman,
``An asymptotic formula for the Bell numbers,''
\emph{Transactions of the Royal Society of Canada}, section III, vol.~49
(1955), pp.~49--54.

\bibitem{NISTDLMF}
F.~W.~J. Olver, A.~B. Olde Daalhuis, D.~W. Lozier, B.~I. Schneider,
R.~F. Boisvert, C.~W. Clark, B.~R. Miller, B.~V. Saunders,
H.~S. Cohl, and M.~A. McClain, editors,
\emph{NIST Digital Library of Mathematical Functions},
National Institute of Standards and Technology.

\bibitem{Riordan1958}
J. Riordan,
\emph{An Introduction to Combinatorial Analysis},
Wiley, 1958.

\bibitem{Roman1984}
S. Roman,
\emph{The Umbral Calculus},
Academic Press, 1984.

\bibitem{Spivey2008}
M.~Z. Spivey,
``A generalized recurrence for Bell numbers,''
\emph{Journal of Integer Sequences}, vol.~11 (2008), Article 08.2.5.

\bibitem{Stanley2012}
R.~P. Stanley,
\emph{Enumerative Combinatorics}, vol.~1, second edition,
Cambridge University Press, 2012.

\bibitem{Temme1993}
N.~M. Temme,
``Asymptotic estimates of Stirling numbers,''
\emph{Studies in Applied Mathematics}, vol.~89 (1993), pp.~233--243.

\bibitem{Wilf2005}
H.~S. Wilf,
\emph{generatingfunctionology}, third edition,
A K Peters, 2005.

\bibitem{Good1960}
I.~J. Good,
``Generalizations to several variables of Lagrange's expansion, with
applications to stochastic processes,''
\emph{Proceedings of the Cambridge Philosophical Society}, vol.~56 (1960),
pp.~367--380.

\bibitem{ConstantineSavits1996}
G.~M. Constantine and T.~H. Savits,
``A multivariate Fa\`a di Bruno formula with applications,''
\emph{Transactions of the American Mathematical Society}, vol.~348 (1996),
pp.~503--520.

\bibitem{FlajoletOdlyzko1990}
P. Flajolet and A. Odlyzko,
``Singularity analysis of generating functions,''
\emph{SIAM Journal on Discrete Mathematics}, vol.~3 (1990), pp.~216--240.

\bibitem{ShapiroEtAl1991}
L.~W. Shapiro, S. Getu, W.-J. Woan, and L.~C. Woodson,
``The Riordan group,''
\emph{Discrete Applied Mathematics}, vol.~34 (1991), pp.~229--239.

\bibitem{Loday2004}
J.-L. Loday,
``Realization of the Stasheff polytope,''
\emph{Archiv der Mathematik}, vol.~83 (2004), pp.~267--278.

\bibitem{Vardi1988}
I. Vardi,
``Determinants of Laplacians and multiple gamma functions,''
\emph{SIAM Journal on Mathematical Analysis}, vol.~19 (1988), pp.~493--507.

\bibitem{Blagouchine2017}
I.~V. Blagouchine,
``A note on some recent results for the Bernoulli numbers of the second kind,''
\emph{Journal of Integer Sequences}, vol.~20 (2017), Article 17.3.8.

\bibitem{MathWorldDossier}
E.~W. Weisstein, editor,
selected entries from \emph{Wolfram MathWorld}: Barnes $G$-function, Bell and
complementary Bell numbers, Bernoulli and Euler polynomials, Bernoulli
polynomials of the second kind, Catalan and Genocchi numbers, Darboux's formula,
Euler--Maclaurin summation, Pochhammer symbols, and polygamma functions.

\bibitem{WikipediaDossier}
Selected reference snapshots supplied in the \nolinkurl{Wikipedia} folder:
associahedra, Bell and Catalan numbers, Bell polynomials, Eulerian numbers,
Fa\`a di Bruno's formula, Lagrange inversion and reversion, and Stirling
numbers.

\bibitem{SourceArchive}
\emph{Formulae-2.zip}, the five LaTeX drafts and the supplied
\nolinkurl{Wikipedia} and \nolinkurl{MathWorld} reference folders used for this
unified article.

\end{thebibliography}

\printindex
\end{document}
